\documentclass[12pt, a4paper, twoside]{book}
%DIF LATEXDIFF DIFFERENCE FILE
%DIF DEL old-flat.tex   Wed Oct 28 17:21:41 2026
%DIF ADD new-flat.tex   Wed Oct 28 17:22:49 2026

\usepackage[utf8]{inputenc}
\usepackage[T1]{fontenc}
\usepackage[english]{babel}
\usepackage[top=3.0cm, bottom=3.0cm, inner=2.5cm, outer=2.0cm]{geometry}
\usepackage{setspace}
\setstretch{1.4}
\raggedbottom
\usepackage{fancyhdr}
\setlength{\headheight}{1cm}
\addtolength{\topmargin}{-0.5cm}
\pagestyle{fancy}
\fancyhf{}
\fancyhead[RE]{\nouppercase{\leftmark}}
\fancyhead[LO]{\nouppercase{\rightmark}}
\fancyfoot[C]{\thepage}
\renewcommand{\headrulewidth}{0pt}
\usepackage{titlesec}
\titleformat{\chapter}[display]{\normalfont\LARGE\bfseries}{\chaptertitlename\ \thechapter}{16pt}{\LARGE}
\setcounter{secnumdepth}{3}

\usepackage{float}
\usepackage{graphicx}
\usepackage{subcaption}
\usepackage{booktabs}
\usepackage{makecell}
\usepackage{multirow}
\usepackage{threeparttable}
\renewcommand{\arraystretch}{1.1}
\usepackage[linesnumbered,ruled,vlined]{algorithm2e}
\SetCommentSty{textnormal}
\usepackage{enumitem}
\usepackage{pdfpages}

\usepackage{amsmath}
\usepackage{amssymb}
\usepackage{amsthm}
\numberwithin{equation}{chapter}
\DeclareMathOperator*{\argmax}{arg\,max}
\DeclareMathOperator*{\argmin}{arg\,min}
\theoremstyle{plain}
\newtheorem{theorem}{Theorem}[chapter]
\newtheorem{proposition}[theorem]{Proposition}
\newtheorem{lemma}[theorem]{Lemma}
\newtheorem{corollary}[theorem]{Corollary}
\newtheorem{condition}[theorem]{Condition}
\theoremstyle{definition}
\newtheorem{definition}{Definition}[chapter]
\newtheorem{assumption}{Assumption}[chapter]
\theoremstyle{remark}
\newtheorem{remark}{Remark}[chapter]

\usepackage{natbib}
\bibliographystyle{unsrt}
\usepackage{xcolor}
\definecolor{darkblue}{rgb}{0, 0, 0.75}
\definecolor{darkgreen}{rgb}{0, 0.75, 0}
\usepackage[colorlinks=true, linkcolor=black, citecolor=darkgreen, urlcolor=darkblue]{hyperref}

\newenvironment{foldable}{}{}
\newcommand{\frontmattersection}[1]{\newpage\thispagestyle{empty}\noindent{\bfseries \Large #1 \par}\vspace{2em}}
\newcommand*{\ie}{i.e.}
\newcommand*{\eg}{e.g.}
\newcommand*{\etc}{etc.}
\newcommand{\meanstd}[2]{#1$_{\pm{#2}}$}
\newcommand{\meanstdbf}[2]{\textbf{#1}$_{\pm{#2}}$}
\newcommand{\boldtight}[1]{{\fontseries{b}\selectfont #1}}
\newcommand{\whitebox}{\textsuperscript{(W)}}
\newcommand{\blackbox}{\textsuperscript{(B)}}
\newcommand{\dg}{$^\dagger$}
\newcommand*{\bltrg}{$^{\blacktriangle}$}

\usepackage{acro}
\DeclareAcronym{AI}{short = AI, long = artificial intelligence}
\DeclareAcronym{API}{short = API, long = application programming interface}
\DeclareAcronym{BBDFKD}{short = BBDFKD, long = black-box data-free knowledge distillation}
\DeclareAcronym{CNN}{short = CNN, long = convolutional neural network}
\DeclareAcronym{DD}{short = DD, long = dataset distillation}
\DeclareAcronym{DIPKD}{short = DIPKD, long = Diverse Image Priors Knowledge Distillation}
\DeclareAcronym{DivBFKD}{short = DivBFKD, long = Diverse Few-Shot Black-Box Knowledge Distillation}
\DeclareAcronym{FID}{short = FID, long = Fréchet inception distance}
\DeclareAcronym{FLOP}{short = FLOP, long = floating point operations}
\DeclareAcronym{GAN}{short = GAN, long = generative adversarial network}
\DeclareAcronym{GDPR}{short = GDPR, long = General Data Protection Regulation}
\DeclareAcronym{HIPAA}{short = HIPAA, long = Health Insurance Portability and Accountability Act}
\DeclareAcronym{HSB}{short = HSB, long = hard-sample bias}
\DeclareAcronym{IS}{short = IS, long = Inception score}
\DeclareAcronym{KD}{short = KD, long = knowledge distillation}
\DeclareAcronym{KL}{short = KL, long = Kullback-Leibler}
\DeclareAcronym{LLM}{short = LLM, long = large language model}
\DeclareAcronym{ML}{short = ML, long = machine learning}
\DeclareAcronym{NLI}{short = NLI, long = natural language inference}
\DeclareAcronym{NLP}{short = NLP, long = natural language processing}
\DeclareAcronym{OT}{short = OT, long = optimal transport}
\DeclareAcronym{RDM}{short = RDM, long = reweighted distribution matching}
\DeclareAcronym{RNN}{short = RNN, long = recurrent neural network}
\DeclareAcronym{SGD}{short = SGD, long = stochastic gradient descent}
\DeclareAcronym{SPD}{short = SPD, long = symmetric positive definite}
\DeclareAcronym{TAKE}{short = TAKE, long = Trajectory-Aware Knowledge Estimation}
\DeclareAcronym{VAE}{short = VAE, long = variational autoencoder}
\DeclareAcronym{WGAN}{short = WGAN, long = Wasserstein generative adversarial network}
\DeclareAcronym{WGAN-GP}{short = WGAN-GP, long = Wasserstein generative adversarial network with gradient penalty} 
\newcommand{\thesistitle}{Knowledge Distillation under Practical Constraints for Models and Datasets} \newcommand{\authorname}{Tri Nhan Vo}
\newcommand{\authordegrees}{BE(Hons)} \newcommand{\degree}{Doctor of Philosophy}
\newcommand{\university}{Deakin University}
\newcommand{\submissiondate}{June, 2026}

\title{\thesistitle}
\author{\authorname}
\date{\submissiondate}
%DIF PREAMBLE EXTENSION ADDED BY LATEXDIFF
%DIF UNDERLINE PREAMBLE %DIF PREAMBLE
\RequirePackage[normalem]{ulem} %DIF PREAMBLE
\RequirePackage{color}\definecolor{RED}{rgb}{1,0,0}\definecolor{BLUE}{rgb}{0,0,1} %DIF PREAMBLE
\providecommand{\DIFaddtex}[1]{{\protect\color{blue}\uwave{#1}}} %DIF PREAMBLE
\providecommand{\DIFdeltex}[1]{{\protect\color{red}\sout{#1}}} %DIF PREAMBLE
%DIF SAFE PREAMBLE %DIF PREAMBLE
\providecommand{\DIFaddbegin}{} %DIF PREAMBLE
\providecommand{\DIFaddend}{} %DIF PREAMBLE
\providecommand{\DIFdelbegin}{} %DIF PREAMBLE
\providecommand{\DIFdelend}{} %DIF PREAMBLE
\providecommand{\DIFmodbegin}{} %DIF PREAMBLE
\providecommand{\DIFmodend}{} %DIF PREAMBLE
%DIF FLOATSAFE PREAMBLE %DIF PREAMBLE
\providecommand{\DIFaddFL}[1]{\DIFadd{#1}} %DIF PREAMBLE
\providecommand{\DIFdelFL}[1]{\DIFdel{#1}} %DIF PREAMBLE
\providecommand{\DIFaddbeginFL}{} %DIF PREAMBLE
\providecommand{\DIFaddendFL}{} %DIF PREAMBLE
\providecommand{\DIFdelbeginFL}{} %DIF PREAMBLE
\providecommand{\DIFdelendFL}{} %DIF PREAMBLE
%DIF HYPERREF PREAMBLE %DIF PREAMBLE
\providecommand{\DIFadd}[1]{\texorpdfstring{\DIFaddtex{#1}}{#1}} %DIF PREAMBLE
\providecommand{\DIFdel}[1]{\texorpdfstring{\DIFdeltex{#1}}{}} %DIF PREAMBLE
\newcommand{\DIFscaledelfig}{0.5}
%DIF HIGHLIGHTGRAPHICS PREAMBLE %DIF PREAMBLE
\RequirePackage{settobox} %DIF PREAMBLE
\RequirePackage{letltxmacro} %DIF PREAMBLE
\newsavebox{\DIFdelgraphicsbox} %DIF PREAMBLE
\newlength{\DIFdelgraphicswidth} %DIF PREAMBLE
\newlength{\DIFdelgraphicsheight} %DIF PREAMBLE
% store original definition of \includegraphics %DIF PREAMBLE
\LetLtxMacro{\DIFOincludegraphics}{\includegraphics} %DIF PREAMBLE
\newcommand{\DIFaddincludegraphics}[2][]{{\color{blue}\fbox{\DIFOincludegraphics[#1]{#2}}}} %DIF PREAMBLE
\newcommand{\DIFdelincludegraphics}[2][]{% %DIF PREAMBLE
\sbox{\DIFdelgraphicsbox}{\DIFOincludegraphics[#1]{#2}}% %DIF PREAMBLE
\settoboxwidth{\DIFdelgraphicswidth}{\DIFdelgraphicsbox} %DIF PREAMBLE
\settoboxtotalheight{\DIFdelgraphicsheight}{\DIFdelgraphicsbox} %DIF PREAMBLE
\scalebox{\DIFscaledelfig}{% %DIF PREAMBLE
\parbox[b]{\DIFdelgraphicswidth}{\usebox{\DIFdelgraphicsbox}\\[-\baselineskip] \rule{\DIFdelgraphicswidth}{0em}}\llap{\resizebox{\DIFdelgraphicswidth}{\DIFdelgraphicsheight}{% %DIF PREAMBLE
\setlength{\unitlength}{\DIFdelgraphicswidth}% %DIF PREAMBLE
\begin{picture}(1,1)% %DIF PREAMBLE
\thicklines\linethickness{2pt} %DIF PREAMBLE
{\color[rgb]{1,0,0}\put(0,0){\framebox(1,1){}}}% %DIF PREAMBLE
{\color[rgb]{1,0,0}\put(0,0){\line( 1,1){1}}}% %DIF PREAMBLE
{\color[rgb]{1,0,0}\put(0,1){\line(1,-1){1}}}% %DIF PREAMBLE
\end{picture}% %DIF PREAMBLE
}\hspace*{3pt}}} %DIF PREAMBLE
} %DIF PREAMBLE
\LetLtxMacro{\DIFOaddbegin}{\DIFaddbegin} %DIF PREAMBLE
\LetLtxMacro{\DIFOaddend}{\DIFaddend} %DIF PREAMBLE
\LetLtxMacro{\DIFOdelbegin}{\DIFdelbegin} %DIF PREAMBLE
\LetLtxMacro{\DIFOdelend}{\DIFdelend} %DIF PREAMBLE
\DeclareRobustCommand{\DIFaddbegin}{\DIFOaddbegin \let\includegraphics\DIFaddincludegraphics} %DIF PREAMBLE
\DeclareRobustCommand{\DIFaddend}{\DIFOaddend \let\includegraphics\DIFOincludegraphics} %DIF PREAMBLE
\DeclareRobustCommand{\DIFdelbegin}{\DIFOdelbegin \let\includegraphics\DIFdelincludegraphics} %DIF PREAMBLE
\DeclareRobustCommand{\DIFdelend}{\DIFOaddend \let\includegraphics\DIFOincludegraphics} %DIF PREAMBLE
\LetLtxMacro{\DIFOaddbeginFL}{\DIFaddbeginFL} %DIF PREAMBLE
\LetLtxMacro{\DIFOaddendFL}{\DIFaddendFL} %DIF PREAMBLE
\LetLtxMacro{\DIFOdelbeginFL}{\DIFdelbeginFL} %DIF PREAMBLE
\LetLtxMacro{\DIFOdelendFL}{\DIFdelendFL} %DIF PREAMBLE
\DeclareRobustCommand{\DIFaddbeginFL}{\DIFOaddbeginFL \let\includegraphics\DIFaddincludegraphics} %DIF PREAMBLE
\DeclareRobustCommand{\DIFaddendFL}{\DIFOaddendFL \let\includegraphics\DIFOincludegraphics} %DIF PREAMBLE
\DeclareRobustCommand{\DIFdelbeginFL}{\DIFOdelbeginFL \let\includegraphics\DIFdelincludegraphics} %DIF PREAMBLE
\DeclareRobustCommand{\DIFdelendFL}{\DIFOaddendFL \let\includegraphics\DIFOincludegraphics} %DIF PREAMBLE
%DIF AMSMATHULEM PREAMBLE %DIF PREAMBLE
\makeatletter %DIF PREAMBLE
\let\sout@orig\sout %DIF PREAMBLE
\renewcommand{\sout}[1]{\ifmmode\text{\sout@orig{\ensuremath{#1}}}\else\sout@orig{#1}\fi} %DIF PREAMBLE
\makeatother %DIF PREAMBLE
%DIF COLORLISTINGS PREAMBLE %DIF PREAMBLE
\RequirePackage{listings} %DIF PREAMBLE
\RequirePackage{color} %DIF PREAMBLE
\lstdefinelanguage{DIFcode}{ %DIF PREAMBLE
%DIF DIFCODE_UNDERLINE %DIF PREAMBLE
  moredelim=[il][\color{red}\sout]{\%DIF\ <\ }, %DIF PREAMBLE
  moredelim=[il][\color{blue}\uwave]{\%DIF\ >\ } %DIF PREAMBLE
} %DIF PREAMBLE
\lstdefinestyle{DIFverbatimstyle}{ %DIF PREAMBLE
	language=DIFcode, %DIF PREAMBLE
	basicstyle=\ttfamily, %DIF PREAMBLE
	columns=fullflexible, %DIF PREAMBLE
	keepspaces=true %DIF PREAMBLE
} %DIF PREAMBLE
\lstnewenvironment{DIFverbatim}{\lstset{style=DIFverbatimstyle}}{} %DIF PREAMBLE
\lstnewenvironment{DIFverbatim*}{\lstset{style=DIFverbatimstyle,showspaces=true}}{} %DIF PREAMBLE
\lstset{extendedchars=\true,inputencoding=utf8}

%DIF END PREAMBLE EXTENSION ADDED BY LATEXDIFF

\begin{document}

\frontmatter
\clearpage{}\begin{titlepage}
  \centering
  \vspace*{\fill}
  {\LARGE \fontfamily{lmr}\bfseries Knowledge Distillation under Practical Constraints \\ for Models and Datasets \par}
  \vspace{2cm}
  {\large by \par}
  \vspace{1cm}
  {\Large \authorname \\ \authordegrees \par}
  \vspace{3cm}
  {\large Submitted in fulfilment of the requirements for the degree of \\ \degree \par}
  \vspace{2cm}
  {\Large \university \par}
  \vspace{0.5cm}
  {\large \submissiondate \par}
  \vspace{2.5cm}
  \vfill
\end{titlepage}

\includepdf[pages=-]{./extra_pages/02-candidature-signed.pdf}

\newpage
\thispagestyle{empty}
\vspace*{\fill}
{
  \centering
  \large
  This research was supported by \\ a Deakin University Postgraduate Research (DUPR) Scholarship.
  \par
}
\vspace*{\fill}


\frontmattersection{ }
{
  \centering
  \textit{
    To my father, my mother, my fiancée-to-be, \\
    You are the lighthouse, that guides this lost me. \\
    From the other hemisphere your light shines through, \\
    No storm I cannot weather, because of you.
  }
  \par
}

\frontmattersection{Acknowledgements}

\begin{foldable}
  A PhD is rarely a solo endeavour, and this one is certainly no exception.
  This work exists because of the generosity, patience, and kindness of more people than I can properly do justice to here, but I will try.
\end{foldable}

\begin{foldable}
  My deepest thanks go to my primary supervisor, Dr.\ Dang~Nguyen, whose guidance, patience, and quiet confidence in me shaped both this work and my growth as a researcher.
  His steady encouragement carried me through all the highs and lows.
  My co-supervisor, Prof.\ Sunil~Gupta, brought sharp expertise and thoughtful questions that helped steer my research in the right direction and consistently pushed the work to be better.
  Dr.\ Kien~Do and Dr.\ Trung~Le were also helpful mentors and generous sounding boards throughout, always willing to engage and offering more than was asked.
\end{foldable}

\begin{foldable}
  The reading groups, hosted by Prof.\ Truyen~Tran and Dr.\ Julian~Berk, kept me connected to the broader field when it was easy to lose perspective inside a narrow problem.
  I would like to give an extra nod to Prof.\ Tran's on his enthusiasm and vision on the landscape of AI, which at times reminds me how my research portraits in the grand picture, and why it matters.
\end{foldable}

\begin{foldable}
  This work was carried out within the Applied AI Initiative at Deakin University, and I am grateful to the people who keep things running behind the scenes.
  Shweta Gupta, Ha Bui, Kim Ngan Pham, and Davina Montgomery provided administrative support with great care and patience.
  I also owe genuine thanks to Josh Cole and Aiden Molluso, who kept the HPC cluster in working order for all the compute resources I needed.
\end{foldable}

\begin{foldable}
  To my post-docs, fellow researchers and friends along the way; some now alumni, some former housemates, some both.
  Thank you for the company, the conversations, and the camaraderie that made this period of life memorable.
  You have helped made not just research but also life in general more enjoyable, and I am grateful for that.
  Specifically, I am grateful to Kien Do, Thao Le, Tri Nguyen, Hung Tran, Kha Pham, Long Dang, Giang Ngo, Dat Phan, Linh Le, Duy Tran, Duc Kieu, Hiep Nguyen, Dai Do, Khoa Le, Hung Le, Toan Nguyen, Thong Bach, Thang Nguyen, Kishan Nagiredla, Maxence Hussonnois, Ajsal Palattuparambil, Azhar Mohammed, Ragja Palakkadavath, and many more.
\end{foldable}

\begin{foldable}
  I also want to mention my two friends from high school whom I never expected to reunite on a different continent, Hy Nguyen (also a fellow researcher) and Khoi Vo.
  The occasional visits and their company reminded me that life extends well beyond just research.
\end{foldable}

\begin{foldable}
  To my pa and ma.
  Thank you for your unwavering belief and support across the years and the distance.
  Your sacrifices extend beyond what words can express.
  You are home.
\end{foldable}

\begin{foldable}
  And to Charisse, my fiancée-to-be.
  Thank you for your love, your faith, and your patience through the long stretches of this journey.
  You have been my anchor throughout.
\end{foldable}

\frontmattersection{List of Publications}

The following peer-reviewed publications form the core research of this thesis and are integrated throughout the chapters.

\begin{enumerate}
  \item {
      Tri-Nhan Vo, Dang Nguyen, Kien Do, and Sunil Gupta.
      \textit{Improving Diversity in Black-box Few-shot Knowledge Distillation.}
      Accepted at \textit{Proceedings of the European Conference on Machine Learning and Principles and Practice of Knowledge Discovery in Databases} (ECML-PKDD), 2024.
      Forms Chapter~\ref{ch:30-research01}.
    }
  \item {
      Tri-Nhan Vo, Dang Nguyen, Trung Le, Kien Do, and Sunil Gupta.
      \textit{Diverse Image Priors for Black-box Data-free Knowledge Distillation.}
      Accepted at \textit{International Conference on Computational Collective Intelligence} (ICCCI), 2026.
      Forms Chapter~\ref{ch:40-research02}.
    }
  \item {
      Tri-Nhan Vo, Dang Nguyen, and Sunil Gupta.
      \textit{TAKE: Trajectory-Aware Knowledge Estimation for Text Dataset Distillation.}
      Accepted at \textit{Proceedings of the European Conference on Machine Learning and Principles and Practice of Knowledge Discovery in Databases} (ECML-PKDD), 2026.
      Forms Chapter~\ref{ch:50-research03}.
    }
\end{enumerate}

\newpage

\includepdf[pages=-]{./extra_pages/06-authorship-statement-research01-signed.pdf}
\includepdf[pages=-]{./extra_pages/06-authorship-statement-research02-signed.pdf}
\includepdf[pages=-]{./extra_pages/06-authorship-statement-research03-signed.pdf}

\chapter*{Abstract} \label{sec:03-abstract}
\addcontentsline{toc}{chapter}{Abstract}

\begin{foldable}
  Recent advances in \acl{AI} have been driven by increasingly large and complex models, trained on increasingly large datasets.
  Yet, model and dataset size are practical bottlenecks across all stages of a model's life cycle, from data preparation and training through to deployment and storage.
  Among compression techniques for both models and datasets, \ac{KD} stands out because it treats compression as a knowledge transfer problem rather than a reduction problem.
  Specifically, model distillation trains a small student model to mimic the outputs of a larger teacher model.
  Dataset distillation, in turn, synthesizes a small, informative dataset that \DIFdelbegin \DIFdel{yields }\DIFdelend \DIFaddbegin \DIFadd{aims to yield }\DIFaddend downstream performance comparable to its larger counterpart.
\end{foldable}

\begin{foldable}
  However, several practical constraints limit the applicability of \ac{KD} in real-world settings.
  Model distillation methods typically assume full access to the original training dataset and to the internal parameters of the teacher network.
  This assumption routinely fails in practice, where data is limited or unavailable because of proprietary or privacy restrictions, and where teacher models are exposed only as black-box interfaces.
  Dataset distillation, on the other hand, has matured for image domains but remains underexplored for text, where the discrete and compositional nature of language poses distinct challenges.
  The distilled text dataset must also withstand extreme compression ratios while remaining human-readable for auditability and interpretability.
\end{foldable}

\begin{foldable}
  This thesis addresses \ac{KD} under three restrictive settings: (i) few-shot access to real data with a black-box teacher, (ii) complete absence of real data with a black-box teacher, and (iii) text dataset distillation, where the goal is to synthesize a tiny surrogate dataset that preserves model training utility.
  We motivate our approaches through the lens of \textit{distributional fidelity}: whether the compact \DIFdelbegin \DIFdel{artefact}\DIFdelend \DIFaddbegin \DIFadd{artifact}\DIFaddend , the student model or the surrogate corpus, faithfully represents the original.
  For the first setting, we propose \textit{\acl{DivBFKD}} (\ac{DivBFKD}), a framework that improves the diversity of the limited training data by generating high-quality synthetic samples, supporting black-box \ac{KD} from only a few samples per class.
  For the second setting, we propose \textit{\acl{DIPKD}} (\ac{DIPKD}), a data-free black-box \ac{KD} method that synthesizes training images by incorporating structured image priors, then optimizes them for diversity via a contrastive objective.
  For the third setting, we propose \textit{\acl{TAKE}} (\ac{TAKE}), which distills a large corpus into a compact dataset of high-influence, human-readable synthetic texts.
  Across all three settings, our methods reach state-of-the-art performance on \ac{KD} benchmarks, showing that effective distillation is achievable under severe practical constraints.
\end{foldable}

\cleardoublepage
\tableofcontents
\listoffigures
\listoftables
\printacronyms
\clearpage{}

\mainmatter
\clearpage{}\acresetall

\chapter{Introduction} \label{ch:10-introduction}

\section{Large-Scale Models and Knowledge Distillation} \label{sec:11-introduction-largescale}
\begin{foldable}
  The history of modern \ac{AI} illustrates the bitter lesson~\cite{sutton2019bitter}: general, compute-driven approaches eventually outperform hand-crafted solutions.
  When AlexNet~\cite{krizhevsky2012imagenet} appeared in 2012, its roughly 60\,M parameters and strong performance on ImageNet signalled that raw model capacity, combined with enough data and compute, could surpass decades of hand-crafted feature engineering.
  An arms race followed that has not slowed since.
  Vision models scaled rapidly: parameter counts of VGGNet~\cite{simonyan2014very} and ResNet~\cite{he2016deep} reached hundreds of millions, while vision transformers soon crossed the 1\,B scale~\cite{dosovitskiy2020image,zhai2022scaling}.
  Language models followed an even steeper trajectory after the Transformer~\cite{vaswani2017attention} architecture arrived, with BERT~\cite{devlin2019bert} featuring 340\,M parameters, followed by GPT-3~\cite{brown2020language} reaching 175\,B parameters the year after.
  As of writing, flagship \ac{AI} models~\cite{anthropic2026assessing,openai2025introducing} are all speculated to have reached trillions of parameters, especially when open-source counterparts too fall in this range~\cite{meta2025llama4,yang2025qwen3}, backed by compute infrastructure demands~\cite{elmeleegy2024demystifying}.
  Their training data have expanded in tandem to tens of trillions of tokens, estimated to have exhausted 80\% of the high-quality English-language surface web~\cite{villalobos2024position}.
  The empirical result is clear: larger models, trained on larger datasets, consistently achieve state-of-the-art performance across vision, language, and multimodal tasks.
\end{foldable}

\begin{foldable}
  Yet scaling success has created a deployment paradox: the models with the strongest benchmark performance are exactly the ones real-world systems struggle most to accommodate.
  Practical deployment is governed not by leaderboard rankings but by hard physical, economic, and legal constraints.
  These constraints operate on two distinct fronts.
\end{foldable}

\begin{foldable}
  The first front concerns the model at inference time.
  Edge devices (smartphones, microcontrollers, or embedded sensors) are subject to fixed limits on silicon area, working memory, power draw, and thermal dissipation.
  A model requiring tens of gigabytes of weights and billions of floating-point operations per forward pass simply cannot execute within these envelopes.
  Safety-critical and real-time applications further compound this with latency requirements.
  For example, an autonomous vehicle processing sensor data cannot tolerate a round-trip to a remote server, a clinical decision-support system must respond within the timescales of patient care, and a fraud detection system must act before a transaction clears.
  Even in cloud settings, where hardware constraints are nominally relaxed, the economics of serving a multi-billion-parameter model at consumer scale are severe.
  Offline deployments such as disaster response, rural healthcare, and maritime or aerospace operations face the same problem in its starkest form: no server is reachable, so the model must run locally or not at all.
  The research community has responded with a family of model compression techniques: pruning, quantization, low-rank adaptation, and \ac{KD}, all aiming to produce smaller, faster models with minimal accuracy loss.
  \ac{KD} in particular transfers knowledge from a teacher model to a student model, requiring no weight access and no structural modification.
\end{foldable}

\begin{foldable}
  The second front concerns the data volume required to build and adapt these models.
  Pre-training on billions of images and documents demands massive infrastructure (storage, pipelines, distributed clusters) available only to a handful of \DIFdelbegin \DIFdel{organisations }\DIFdelend \DIFaddbegin \DIFadd{organizations }\DIFaddend worldwide.
  This incurs expense on training and fine-tuning, often requiring hundreds of hours and terabytes of task-specific data.
  For \DIFdelbegin \DIFdel{organisations }\DIFdelend \DIFaddbegin \DIFadd{organizations }\DIFaddend adapting models frequently, cumulative costs rival the original pre-training expense.
  Energy expenditure compounds this further: a single training run emits carbon dioxide comparable to several passenger vehicles~\cite{strubell2019energy}.
  Beyond these economic pressures, data access is independently constrained by legal and institutional barriers.
  Privacy legislation such as the \ac{GDPR}~\cite{eu2016gdpr} in Europe and the \ac{HIPAA}~\cite{us1996hipaa} in the United States restrict data centralization.
  Datasets are also proprietary, with \DIFdelbegin \DIFdel{organisations }\DIFdelend \DIFaddbegin \DIFadd{organizations }\DIFaddend unwilling to share data that underpins their competitive advantage.
  This creates friction for collaboration and research, since researchers must either forgo real-world data or work through complex legal and ethical requirements to access it.
  Dataset compression techniques, coreset selection, prototype-based methods, and \ac{DD}, address these constraints.
  \ac{DD} in particular optimizes a synthetic set so that models trained on it \DIFdelbegin \DIFdel{match }\DIFdelend \DIFaddbegin \DIFadd{approach }\DIFaddend full-dataset performance, capturing essential knowledge within a minimal data budget.
\end{foldable}

\begin{foldable}
  Both compression dimensions are independently valuable, but neither is sufficient on its own: a practitioner who compresses only the model still faces the full cost of training data at adaptation time, while one who compresses only the data still deploys a model too large for the target hardware.
  Together, model and dataset distillation\footnote{This thesis uses \textit{knowledge distillation} as an umbrella term for both \textit{model distillation} and \textit{dataset distillation}; the two terms are used explicitly when the contexts are ambiguous.} address the full spectrum of these constraints, in a way that sets them apart from other compression families.
  Both treat compression as a knowledge transfer problem through different mechanisms: model distillation via teacher output, dataset distillation via synthetic data.
  Both are therefore concerned with the fidelity of what is transferred, not just the efficiency of the transfer.
  This concern is what makes the practical gaps non-trivial: when data access is restricted, when the teacher is a black box, or when the target modality is discrete, the central question becomes whether the distilled signal remains faithful to the original.
  The following sections examine where these gaps arise and how the contributions of this thesis address them.
\end{foldable}
 \section{Model Distillation and Practical Gaps} \label{sec:12-introduction-model}
\begin{foldable}
  The research community has responded with model compression techniques that recover deployable models without giving up too much of the performance earned through scale.
  \textit{Pruning} removes redundant weights or entire sub-networks, reducing parameter counts and, with appropriate hardware support, inference time~\cite{han2015learning,molchanov2016pruning}.
  \textit{Quantization} reduces the numerical precision of weights and activations from standard 32-bit floating point to lower-bit representations, cutting memory footprint and accelerating arithmetic on compatible hardware~\cite{hubara2018quantized,jacob2018quantization}.
  \textit{Low-rank adaptation} decomposes weight matrices into products of lower-rank factors, reducing the number of trainable parameters while preserving model expressivity~\cite{hu2022lora}.
  Each of these approaches has yielded meaningful compression ratios, and in practice they are often combined.
  Among them, \textit{\acl{KD}}~\cite{hinton2015distilling} occupies a distinct position: it frames compression as a learning problem in which a compact \textit{student} network is trained to replicate the \DIFdelbegin \DIFdel{behaviour }\DIFdelend \DIFaddbegin \DIFadd{behavior }\DIFaddend of a large, pre-trained \textit{teacher} network.
  \ac{KD} is architecture-agnostic: the student need not share the teacher's topology or depth.
  It never touches the teacher's internal parameters, giving the freedom to design students suited to their deployment targets.
  What makes the approach work is what the teacher conveys beyond its final decision.
\end{foldable}

\begin{foldable}
  Transferring knowledge as hard labels are sparse: all probability mass concentrates on a single correct class, discarding any signal about the relative plausibility of the rest.
  The teacher's output distribution, by contrast, spreads probability mass across many classes in a way that reflects its learned understanding of the input: for a given image of a cat, assigning most mass to \texttt{cat}, a smaller amount to \texttt{dog}, and a tiny but non-zero amount to \texttt{leopard} or \texttt{tiger}, while assigning essentially nothing to \texttt{car} or \texttt{aeroplane}.
  These non-zero probabilities encode inter-class similarity, also termed \textit{dark knowledge}: it is present in every teacher query but discarded when supervision is collapsed to a hard label~\cite{hinton2015distilling}.
  Training the student to match these \textit{soft targets} transfers this relational structure directly.
  This is why the \ac{KD} objective combines a \ac{KL} divergence on temperature-scaled teacher and student distributions with a cross-entropy term on hard labels.
  Yet this effectiveness rests on two critical assumptions: (1)~the full training dataset is available for querying the teacher, and (2)~the teacher exposes its complete output distribution rather than a truncated prediction.
  In many real deployments, both assumptions fail.
\end{foldable}

\begin{foldable}
  Consider first the assumption on having data access.
  Training data for high-performing models is often the most commercially and legally sensitive asset held by an \DIFdelbegin \DIFdel{organisation}\DIFdelend \DIFaddbegin \DIFadd{organization}\DIFaddend .
  Privacy legislation such as \ac{GDPR}\cite{eu2016gdpr} and \ac{HIPAA}\cite{us1996hipaa} in the United States strictly limits the collection, retention, and redistribution of personal or medical records, making it legally impossible to share the data used to train a clinical imaging model with a third-party deployer.
  Even where regulation is not a binding constraint, intellectual property law and commercial confidentiality create strong incentives to withhold training corpora: a company that has spent years curating a proprietary dataset will not release it merely to facilitate compression of the model it trained.
  Beyond legal and strategic barriers, raw cost is a practical obstacle: large-scale datasets in computer vision and natural language processing can occupy terabytes of storage and require expensive curation pipelines, making them infeasible to distribute or even store locally at the deployment site.
  The net effect is that knowledge transfer is gated by access to the original training data.
  This constraint exists on a spectrum, from \textit{few-shot} settings in which only a handful of samples are available, to fully \textit{data-free} settings in which no original data exists at all, with the latter being the strictly harder regime.
\end{foldable}

\begin{foldable}
  The second assumption, white-box access to the teacher, is equally fragile in practice.
  The dominant commercial model for deploying large neural networks has shifted decisively toward black-box inference services.
  GPT-style language models, vision classification services, and multimodal foundation models are made available as hosted endpoints: a client submits an input and receives a prediction, but the model weights, architecture, and internal activations never leave the server.
  Even the output is often severely truncated; many production services return only the top-predicted class label or a short ranked list, rather than the full probability distribution over thousands of categories.
  This hard-label regime is a deliberate design choice, since returning full softmax distributions would expose the model to higher risks of extraction attacks and reverse-engineering.
  The consequence for distillation is severe: without access to the teacher's soft probability vector, the student loses the inter-class similarity information that makes \ac{KD} more effective than training on one-hot labels alone.
\end{foldable}

\begin{foldable}
  Taken together, these two gaps force compounding compromises: absent or incomplete training data degrades the coverage of the distillation signal, while black-box access degrades its quality.
  The effect is more than a practical inconvenience.
  Each constraint reshapes the statistical character of whatever signal remains: the student is left with a signal that is at once narrower in scope and poorer in relational content than the original \ac{KD} formulation assumes.
  This raises a basic question: how much of the teacher's knowledge can be faithfully recovered under such restrictions?
  \S\ref{sec:14-introduction-fidelity} frames this question through the lens of distributional fidelity and examines what it demands of any principled solution.
\end{foldable}
 \section{Dataset Distillation and Practical Gaps} \label{sec:13-introduction-dataset}
\begin{foldable}
  Alongside model compression, a parallel line of research targets the dataset side of the problem.
  Several families of techniques have emerged for reducing the volume of training data while preserving its utility.
  \textit{Prototype} and \textit{clustering}-based methods represent the full dataset by a small set of centroids, discarding individual samples once the cluster structure is captured~\cite{macqueen1967some,dempster1977maximum}.
  \textit{Coreset selection} methods score and rank individual samples by their estimated importance to the learned model, then retain only the top-ranked subset~\cite{sener2018active}.
  \textit{Dataset distillation} (\ac{DD}) takes a fundamentally different approach: rather than selecting from existing samples, it synthesizes an entirely new, compact dataset whose training signal matches that of the full corpus~\cite{wang2018dataset}.
  This end-to-end optimization of the distilled set, rather than the selection or \DIFdelbegin \DIFdel{summarisation }\DIFdelend \DIFaddbegin \DIFadd{summarization }\DIFaddend of existing ones, is what distinguishes \ac{DD} from all prior dataset compression methods and constitutes the central subject of this section.
\end{foldable}

\begin{foldable}
  The intuition behind \ac{DD} mirrors that of \ac{KD}: a large dataset encodes knowledge not merely in its individual labels, but in the statistical relationships among its samples.
  Those include the frequencies, co-occurrence patterns, decision boundaries, and distributional structure that a model \DIFdelbegin \DIFdel{internalises }\DIFdelend \DIFaddbegin \DIFadd{internalizes }\DIFaddend over the course of training.
  The central question \ac{DD} poses is whether a much smaller synthetic set can preserve this distributed knowledge faithfully enough that a model trained on it behaves comparably to one trained on the full dataset.
  This is a stronger requirement than coreset selection, which asks only that the retained samples be individually representative.
  Specifically, \ac{DD} asks that the distilled set collectively reproduce the learning dynamics of the original dataset.
\end{foldable}

\begin{foldable}
  The objective is intuitive: create a small synthetic dataset such that a model trained on it learns just as well as one trained on the full corpus.
  The seminal approach, \textit{gradient matching}~\cite{wang2018dataset}, formalizes this by optimizing the distilled set so that the gradients it produces during training, computed from the training loss at each step, closely match those produced by the full dataset.
  Since gradients drive parameter updates, matching them ensures the model follows nearly identical optimization steps regardless of which dataset it trains on.
  In other words, the distilled set steers the learning process along nearly the same path as the full data would, even though it contains far fewer examples.
  The idea is simple but strong: if the learning trajectory stays the same, the final model should be nearly identical.
  Subsequent refinements have explored matching full training trajectories~\cite{cazenavette2022dataset} or embedding distributions~\cite{zhao2023dataset} rather than individual gradients, but all share this core principle.
\end{foldable}

\begin{foldable}
  Despite this progress, existing \ac{DD} methods have been developed almost exclusively for image data, and their core mechanisms do not transfer to the text modality.
  The mismatch is not superficial: it reflects two structural incompatibilities.
  The first is that text operates over \textit{discrete tokens} that resist interpolation.
  Gradient matching, trajectory matching, and distribution matching all rely on computing gradients through the distilled samples with respect to their continuous pixel values; a small perturbation to a pixel produces a smooth, differentiable change in the loss.
  Text tokens admit no such perturbation: changing a token is an all-or-nothing substitution, and omitting even a single rare but semantically critical token constitutes a complete loss of the knowledge it encodes.
  The second incompatibility is that text exhibits \textit{non-local dependencies}: the meaning of a sentence can hinge on a word several positions away, and the coherence of a paragraph depends on relationships that span its full length.
  Image \ac{DD} methods originally treat distilled image samples as continuous pixel values, which enables gradient-based optimization directly.
  For texts, however, tokens are discrete and non-differentiable, and workarounds that optimize in continuous embedding space produce architecture-locked vectors that cannot be audited or reused, sacrificing the human-readability that makes a distilled corpus practical.
\end{foldable}

\begin{foldable}
  The second gap is more fundamental and intensifies at extreme compression ratios: the smaller the distilled set, the higher the knowledge content placed on each retained sample.
  In the text modality, this burden is compounded by discrete sparsity.
  A rare token that appears only a handful of times in the full corpus may be absent entirely from a small distilled set, yet the knowledge it encodes may be irreplaceable.
  Standard importance-scoring approaches, which rank samples independently by their estimated influence on the learned model, cannot correct for this.
  Independent scoring captures the importance of each sample in isolation but ignores dataset-level structure, and will systematically omit rare but critical linguistic patterns in \DIFdelbegin \DIFdel{favour }\DIFdelend \DIFaddbegin \DIFadd{favor }\DIFaddend of high-frequency, easy-to-score examples.
  At the most aggressive compression ratios, this failure mode becomes categorical.
  Each distilled sample must carry more knowledge than any single sample in the original dataset, yet the scoring mechanism provides no way to enforce that the selected set collectively spans the full distribution.
\end{foldable}

\begin{foldable}
  Together, these gaps define the requirements for a principled text \ac{DD} method: it must handle discrete tokens without continuous interpolation, preserve the distributional structure of the original corpus collectively rather than sample-by-sample, and produce human-readable output that can be audited and reused.
  \S\ref{sec:14-introduction-fidelity} examines the distributional consequences of failing to meet these requirements and identifies the unifying principle, \textit{distributional fidelity}, that connects the gaps in model distillation and dataset distillation alike.
\end{foldable}
 \section{Distributional Fidelity of the Training Signal} \label{sec:14-introduction-fidelity}
\begin{foldable}
  The restrictions described in the previous subsections do not merely reduce the quantity of available supervision, but instead reshape its statistical character.
  A student trained under a limited data constraint, or a model fine-tuned on a compressed corpus can only learn from whatever training signal is actually present.
  If that signal fails to faithfully represent the distribution the teacher was trained on, then the student inherits a systematically distorted picture of the teacher's knowledge.
  The term \textit{distributional fidelity} names this property: the degree to which the training signal available to the student covers the target distribution in a representative and unbiased way.
  \DIFaddbegin \DIFadd{In this thesis it takes two related forms: for Chapters~\ref{ch:30-research01} and~\ref{ch:40-research02}, coverage of the useful parts of the teacher's input or representation space; for Chapter~\ref{ch:50-research03}, matching an importance-weighted source distribution in an embedding space.
  }\DIFaddend When distributional fidelity is low, distillation fails because the information it needs is not present in the signal it receives.
\end{foldable}

\begin{foldable}
  In model distillation, distributional fidelity shows up as a coverage problem that worsens as data access diminishes.
  Gaps in coverage correspond directly to gaps in the student's knowledge domain: regions of the teacher's input distribution that the student never trains on remain unknown to it, regardless of model capacity.
  In the few-shot setting, a handful of real samples cannot span the full diversity of the teacher's input distribution, and synthetic data generated from them inherits the same distributional limits.
  Generative models can extend synthesis beyond these limits, but without external guidance they remain tethered to the modes the few-shot set already represents.
  Diversity of the generated set, specifically its ability to reach regions the few-shot examples do not represent, is therefore the binding challenge.
  In the data-free setting, the problem is more severe: without any real-data anchor, synthetic images serve as a substitute, seeding synthesis with structured diversity before any teacher query.
  Yet this does not guarantee coverage: the central challenge is ensuring that synthetic samples collectively capture \textit{diverse patterns\DIFdelbegin \DIFdel{and semantics}\DIFdelend }\DIFaddbegin \DIFadd{, and ideally semantics, }\DIFaddend of the teacher's knowledge domain, rather than merely satisfying its hard labels.
\end{foldable}

\begin{foldable}
  The black-box constraint compounds the limited data setting by further degrading the signal quality.
  In practice, black-box access exists on a spectrum: some services return a truncated probability distribution over the top-$k$ classes, while others return only a single integer identifying the predicted class.
  Even the softer end of this spectrum degrades distributional fidelity: a truncated distribution discards the long tail of inter-class relationships that the teacher's full softmax encodes.
  At the hard-label extreme, the distilled prediction distribution collapses to a one-hot spike, carrying only the identity of the predicted class and none of the relational structure that distinguishes \ac{KD} from ordinary supervised training.
  The central challenge is therefore to recover knowledge in the form of soft labels, the inter-class relational structure that grounds the original \ac{KD} formulation~\cite{hinton2015distilling}, from a teacher interface that deliberately withholds it.
\end{foldable}

\begin{foldable}
  Dataset distillation tackles distributional fidelity from a different angle: preserving knowledge for downstream utility.
  Compared with prototype and coreset methods, dataset distillation approximates the original data with synthetic samples, free of the constraint of the existing data support.
  In principle, this can raise the fidelity and downstream quality of the distilled data.
  At extreme compression ratios, however, preserving fidelity is still hard: each distilled sample must carry a disproportionately higher informational load.
  A principled method must therefore perform \textit{knowledge estimation} on each sample and ensure that the distilled set collectively covers the full distribution, rather than over-representing some regions and under-representing or omitting others.
  The optimization objective must therefore be \textit{global}: selecting samples independently by any per-sample score cannot guarantee coverage, since it ignores relationships between chosen samples and leaves redundant or unrepresented regions uncorrected.
  The \textit{discrete} nature of text also prevents the original gradient-based optimization from transferring directly.
  This constraint demands a discrete optimization method that handles text tokens without sacrificing human-readability.
\end{foldable}

\begin{foldable}
  All three settings exhibit the same underlying problem: failure modes that undermine distributional fidelity.
  In model distillation, coverage deficit arises through generative collapse in \textit{black-box few-shot} or lack of knowledge anchor \textit{black-box data-free}.
  In dataset distillation, it arises through compression and interpretability requirements of \textit{text datasets}.
  Across all cases, the training signal fails to faithfully represent the target distribution: it over-represents some regions and under-represents or entirely omits others.
  The diagnostic question is identical in every setting: does the data the student actually trains on give sufficient, representative coverage of the knowledge the teacher possesses?
  \DIFdelbegin \DIFdel{Fidelity deficit is }\DIFdelend \DIFaddbegin \DIFadd{We take fidelity deficit to be }\DIFaddend the root cause of these failures, not a symptom.
  \DIFdelbegin \DIFdel{No algorithmic refinement can }\DIFdelend \DIFaddbegin \DIFadd{Algorithmic refinement alone cannot }\DIFaddend recover knowledge that is absent from the training signal.
  \DIFdelbegin \DIFdel{Recognising }\DIFdelend \DIFaddbegin \DIFadd{Recognizing }\DIFaddend this unifies what would otherwise appear as three separate technical problems.
\end{foldable}

\begin{foldable}
  Treating distributional fidelity as a first-class design objective, rather than a collateral outcome, is the commitment that motivates all three contributions of this thesis.
  Each contribution targets a distinct failure mode, unified by one principle: distillation quality is bounded by signal representativeness, not algorithmic sophistication.
  \S\ref{sec:15-introduction-contributions} describes how each contribution realizes this principle in concrete, experimentally validated form.
\end{foldable}
 \section{Thesis Contributions} \label{sec:15-introduction-contributions}
\begin{foldable}
  The three restricted settings identified in \S\ref{sec:12-introduction-model} and \S\ref{sec:13-introduction-dataset} are few-shot and data-free black-box model distillation, and text dataset distillation.
  Each imposes a distinct form of information scarcity on the distillation process.
  Structurally, these settings map to two complementary axes of \ac{KD}: the first operates on the model axis (Chapters~\ref{ch:30-research01}--\ref{ch:40-research02}), compressing the learned function under restricted teacher access; the second operates on the data axis (Chapter~\ref{ch:50-research03}), compressing the training signal under discrete-token and coverage constraints.
  Across both axes, \DIFdelbegin \DIFdel{distributional fidelity is }\DIFdelend \DIFaddbegin \DIFadd{this thesis treats distributional fidelity as }\DIFaddend the binding constraint.
  The thesis addresses all three settings under a \DIFdelbegin \DIFdel{single design philosophy}\DIFdelend \DIFaddbegin \DIFadd{shared design principle}\DIFaddend : treat distributional fidelity as a first-class objective, derive practical constraints from the deployment setting, and enforce representational coverage of the distilled signal explicitly.
  The three contributions are as follows.
\end{foldable}

\begin{foldable}
  \boldtight{Diversity-Aware Black-Box Few-Shot Knowledge Distillation} (\textbf{\acs{DivBFKD}}).
  In Chapter~\ref{ch:30-research01}, \ac{DivBFKD} addresses the setting in which only a small number of \DIFaddbegin \DIFadd{labeled }\DIFaddend real training images are available and the teacher is accessible only as a black box, returning \DIFdelbegin \DIFdel{final predictions }\DIFdelend \DIFaddbegin \DIFadd{class probabilities }\DIFaddend rather than internal representations.
  The method trains a \ac{WGAN} to synthesize additional training images, augmenting this process with a high-confidence image scheme that identifies and retains synthetic samples whose teacher-assigned confidence exceeds class-specific adaptive thresholds computed via the $q$-quantile of the score distribution.
  This threshold acts as a quality gate, admitting only \DIFdelbegin \DIFdel{sufficiently informative and class-coherent synthetic images into the distillation pipeline, and directly counters the semantic narrowness }\DIFdelend \DIFaddbegin \DIFadd{synthetic images that the teacher classifies confidently, and counters the narrow coverage }\DIFaddend that arises when few real examples are available.
  We evaluate across seven image classification benchmarks and reach state-of-the-art performance \DIFaddbegin \DIFadd{among the compared few-shot methods}\DIFaddend , with a gain of \DIFdelbegin \DIFdel{over }\DIFdelend 17\% \DIFdelbegin \DIFdel{relative }\DIFdelend \DIFaddbegin \DIFadd{compared }\DIFaddend to standard \ac{KD} on CIFAR-10 and closing to within \DIFaddbegin \DIFadd{about }\DIFaddend 2\% of a fully supervised student upper bound on \DIFdelbegin \DIFdel{the simpler datasets}\DIFdelend \DIFaddbegin \DIFadd{MNIST, FMNIST, and SVHN}\DIFaddend .
  This work was published at ECML-PKDD 2024.
\end{foldable}

\begin{foldable}
  \boldtight{Diverse Image Priors for Black-Box Data-Free Knowledge Distillation} (\textbf{\acs{DIPKD}}).
  Later, in Chapter~\ref{ch:40-research02}, \ac{DIPKD} addresses the strictest image-domain setting: the student has access to neither real training data nor teacher internals, receiving only top-1 hard labels in response to queries.
  The method is \DIFdelbegin \DIFdel{organised }\DIFdelend \DIFaddbegin \DIFadd{organized }\DIFaddend as a three-phase pipeline.
  In the first phase, a diverse pool of synthetic images is built using hierarchical noise injection, nonlinear \DIFdelbegin \DIFdel{photometric transforms}\DIFdelend \DIFaddbegin \DIFadd{filters}\DIFaddend , and semantic cutmixing \DIFaddbegin \DIFadd{whose masks are irregular regions rather than object segments}\DIFaddend , producing candidate images that cover a broad region of the input space.
  In the second phase, a \DIFdelbegin \DIFdel{lightweight }\DIFdelend primer student is trained on these candidates under \DIFdelbegin \DIFdel{a }\DIFdelend \DIFaddbegin \DIFadd{an instance-level }\DIFaddend contrastive objective that \DIFdelbegin \DIFdel{maximises inter-class separation}\DIFdelend \DIFaddbegin \DIFadd{pushes each synthetic image away from the others}\DIFaddend , sharpening the diversity of the pool before full distillation begins.
  In the third phase, the trained primer student supplies surrogate \DIFdelbegin \DIFdel{logits as soft distillation labels }\DIFdelend \DIFaddbegin \DIFadd{soft supervision }\DIFaddend to supplement the teacher's hard labels, \DIFdelbegin \DIFdel{recovering much of }\DIFdelend \DIFaddbegin \DIFadd{partly compensating for }\DIFaddend the information lost to black-box access.
  \ac{DIPKD} reaches state-of-the-art results on 11 of 12 benchmarks, with a gain of over 21\% on Imagenette relative to the best prior baseline, and shows strong performance on medical imaging datasets from MedMNIST.
  This work has been submitted to ICCCI 2026.
\end{foldable}

\begin{foldable}
  \boldtight{Trajectory-Aware Knowledge Estimation for Text Dataset Distillation} (\textbf{\acs{TAKE}}).
  Finally, in Chapter~\ref{ch:50-research03}, \ac{TAKE} addresses dataset distillation in the natural language domain, where the goal is to distill a \DIFdelbegin \DIFdel{compact subset of a }\DIFdelend large text corpus \DIFdelbegin \DIFdel{that can surrogates }\DIFdelend \DIFaddbegin \DIFadd{into a compact synthetic surrogate corpus, drawn from an LLM-generated pool, that approximates }\DIFaddend the performance of training on the full dataset.
  Existing text dataset distillation methods \DIFaddbegin \DIFadd{each }\DIFaddend fall short on \DIFaddbegin \DIFadd{at least one of }\DIFaddend three fronts: their weighting schemes are uniform and uninformative about task-aligned importance, their objectives optimize batch by batch \DIFaddbegin \DIFadd{or per sample }\DIFaddend rather than enforcing global distributional coverage, \DIFdelbegin \DIFdel{and }\DIFdelend \DIFaddbegin \DIFadd{or }\DIFaddend their outputs are embeddings rather than human-readable text.
  \ac{TAKE} addresses all three through a multi-stage methodology.
  First, it quantifies each sample's downstream contribution via influence functions and integrates these estimates across the full training trajectory.
  Second, it generates a pool of human-readable synthetic candidates through language model fine-tuning, then distills a compact corpus via discrete optimal transport that enforces global distributional coverage rather than relying on per-sample independence.
  Theoretically, \ac{TAKE} formalizes the hard-sample bias in text distillation and derives formal bounds on the downstream loss gap.
  Empirically, \ac{TAKE} matches or exceeds baseline methods across six text datasets at extreme compression ratios, producing human-readable corpora that generalize across diverse model architectures.
  This work has been submitted to ECML-PKDD 2026.
\end{foldable}

\begin{foldable}
  Operationally, the design philosophy becomes concrete mechanisms tailored to the information scarcity each setting imposes.
  \ac{DivBFKD} (image, black-box few-shot) filters synthetic candidates through generation and adaptive confidence thresholds, admitting only \DIFdelbegin \DIFdel{semantically coherent, informative examples }\DIFdelend \DIFaddbegin \DIFadd{examples the teacher classifies confidently }\DIFaddend into distillation.
  \ac{DIPKD} (image, black-box data-free) synthesizes diverse image priors, applies contrastive learning to sharpen distinction between samples, then distills via a primer student that supplies soft signals, addressing the diversity deficit that arises without real data.
  \ac{TAKE} (text dataset) concentrates budget through knowledge-based influence weighting and discrete optimal transport, identifying high-value samples (not uniform, uninformative ones), while trajectory integration corrects hard-sample bias and enforces global coverage.
  Across all three, the unifying insight is the same: representational coverage and quality are not implicit properties of the training signal, but design objectives that require explicit measurement and enforcement.
  Together, these contributions show that \ac{KD} under restricted access is feasible and can match or exceed \DIFaddbegin \DIFadd{several compared }\DIFaddend methods that operate under \DIFdelbegin \DIFdel{far }\DIFdelend less constrained conditions.
\end{foldable}
 \section{Thesis Structure} \label{sec:16-introduction-structure}
\begin{foldable}
  This thesis is \DIFdelbegin \DIFdel{organised }\DIFdelend \DIFaddbegin \DIFadd{organized }\DIFaddend around two complementary tracks unified by \textit{distributional fidelity}.
  \textbf{Model Track} addresses model distillation under restricted teacher access, progressing from few-shot to data-free settings.
  \textbf{Dataset Track} addresses dataset distillation for text, producing compact, human-readable corpora under extreme compression.
\end{foldable}

\begin{foldable}
  \textbf{Chapter~\ref{ch:20-background}} establishes shared foundations, then divides into two parallel strands for model distillations and dataset distillation, accordingly.
  Both tracks converge on \textit{distributional fidelity}: whether the compressed artifact, student model or surrogate corpus, faithfully represents the original training distribution, enabling effective distillation under restricted access.
  \textbf{Model Track (Chapters~\ref{ch:30-research01}--\ref{ch:40-research02})} develops model distillation with a black-box teacher, with an emphasis on image classification.
  Chapter~\ref{ch:30-research01} addresses few-shot constraints, where a small number of real training images are available.
  Chapter~\ref{ch:40-research02} removes access to real training data, forcing distillation from black-box signals alone.
  \textbf{Dataset Track (Chapter~\ref{ch:50-research03})} develops dataset distillation, specifically distilling a large training corpus into a compact set of human-readable examples.
  Discrete tokens and extreme compression ratios introduce distinct constraints, but the unifying principle is the same: distributional fidelity under information scarcity.
  \textbf{Chapter~\ref{ch:60-conclusion}} draws the contributions across both tracks together, surfacing common themes and placing them within the trajectory of \ac{KD} research.
  The chapter discusses limitations and identifies directions for future work.
  Readers may follow the thesis sequentially or by track (Model Track in Chapters~\ref{ch:30-research01}--\ref{ch:40-research02}, Dataset Track in Chapter~\ref{ch:50-research03}).
\end{foldable}
 \clearpage{}
\clearpage{}\acresetall

\chapter{Background} \label{ch:20-background}

\begin{foldable}
  This chapter sets out a two-track overview in which \textit{distributional fidelity} is the unifying design criterion.
  \S\ref{sec:21-background-deeplearning}--\S\ref{sec:22-background-compression} establish the shared foundation: deep learning architectures across vision and language, and the compression landscape.
  From there the chapter divides into two parallel tracks.
\end{foldable}

\begin{foldable}
  \textbf{Model Track}:
  \S\ref{sec:23-background-modeldistill}--\S\ref{sec:24-background-kdrestricted} develops the model distillation pipeline from standard \ac{KD} (soft-target transfer, dark knowledge, temperature scaling) through the taxonomy of methods (response/feature/relation-based) to the restricted-access setting that arises in practice.
  Two restrictions define this landscape: teacher access and data availability.
  As restrictions tighten, the distillation signal degrades: few-shot or data-free combined with black-box access creates a double bind where both coverage and signal quality are simultaneously constrained.
\end{foldable}

\begin{foldable}
  \textbf{Dataset Track}:
  \S\ref{sec:26-background-datasetdistill}--\S\ref{sec:28-background-ot} develops the parallel problem on the data axis: compressing a training corpus into a small surrogate that induces similar model behavior.
  Text dataset distillation under extreme compression requires knowledge-informed weighting and distributional coverage for utility, and must remain human-readable for interpretability.
\end{foldable}

\begin{foldable}
  \textbf{Unifying Principle}:
  Both tracks are connected by \S\ref{sec:25-background-distribution}, which establishes \textit{distributional fidelity} as the binding constraint: whether the compressed \DIFdelbegin \DIFdel{artefact }\DIFdelend \DIFaddbegin \DIFadd{artifact }\DIFaddend is a student model (Model Track) or a surrogate corpus (Dataset Track), it must \DIFdelbegin \DIFdel{faithfully recover }\DIFdelend \DIFaddbegin \DIFadd{adequately represent }\DIFaddend the original distribution.
\end{foldable}

\begin{foldable}
  \S\ref{sec:29-background-summary} draws both tracks together and states three open problems that Chapters~\ref{ch:30-research01}--\ref{ch:50-research03} address: few-shot black-box \ac{KD}, data-free black-box \ac{KD}, and text dataset distillation.
\end{foldable}

\section{Deep Learning Foundations} \label{sec:21-background-deeplearning}
\subsection{Neural Networks and Supervised Learning} \label{sec:21-background-deeplearning-nn}
\begin{foldable}
  A feedforward neural network maps an input vector $x \in \mathbb{R}^{d}$ to an output $\hat{\mathbf{y}}$ through a sequence of \DIFdelbegin \DIFdel{parameterised }\DIFdelend \DIFaddbegin \DIFadd{parameterized }\DIFaddend transformations.
  Each layer applies an affine map followed by a pointwise non-linearity: $h = \sigma(Wh_{\text{prev}} + b)$, where $W$ and $b$ are learnable weight matrices and bias vectors, and $\sigma$ is the activation function~\cite{rumelhart1986learning}.
  Stacking many such layers lets the network compose simple transformations into more abstract representations, a property that underlies the empirical success of deep architectures across vision, language, and beyond.
  Early networks used sigmoid or hyperbolic tangent activations, both of which saturate for large input magnitudes and produce near-zero gradients that impede learning in deep stacks~\cite{bengio1994learning}.
  The rectified linear unit $\text{ReLU}(z) = \max(0, z)$ largely resolved this issue: it is non-saturating for positive inputs, cheap to compute, and has become the default activation in modern networks~\cite{glorot2011deep}.
  Variants such as Leaky ReLU~\cite{maas2013rectifier} and GELU~\cite{hendrycks2017gaussian} have since been introduced to handle the dying-neuron problem and improve performance in specific settings, but the core design principle remains unchanged.
\end{foldable}

\begin{foldable}
  In supervised learning, a labeled dataset $\mathcal{D} = \{(x_i, y_i)\}_{i=1}^{N}$ is used to \DIFdelbegin \DIFdel{minimise }\DIFdelend \DIFaddbegin \DIFadd{minimize }\DIFaddend an empirical loss with respect to the network parameters $\theta$:
  \begin{equation}
    \label{eq:21-background-deeplearning-loss-supervised}
    \mathcal{L}(\theta) = \frac{1}{N}\sum_{i=1}^{N} \ell(f_{\theta}(x_i), y_i).
  \end{equation}
  For classification tasks, $\ell$ takes the form of the cross-entropy loss between the network's predicted class probabilities $\hat{\mathbf{p}} = \text{softmax}(f_\theta(x))$ and the ground-truth one-hot label $\mathbf{y}$: $\ell = -\sum_c y_c \log \hat{p}_c$.
  \DIFdelbegin \DIFdel{Minimising }\DIFdelend \DIFaddbegin \DIFadd{Minimizing }\DIFaddend this objective encourages the network to assign high probability to the correct class and low probability to all others.
  Parameters are updated via backpropagation, which applies the chain rule to compute the gradient of $\mathcal{L}$ with respect to every parameter in the network.
  Stochastic gradient descent~\cite{robbins1951stochastic} and its adaptive variants such as Adam~\cite{kingma2014adam} perform these updates on randomly sampled mini-batches rather than the full dataset, trading gradient accuracy for computational efficiency and often improving generalization through the implicit noise introduced by mini-batch sampling.
\end{foldable}

\begin{foldable}
  Generalization does not follow automatically from fitting the training data well~\cite{zhang2017understanding}.
  A model with insufficient capacity underfits: it cannot represent the true input-output mapping even on training data.
  A model with excess capacity relative to the training set overfits: it memorizes training examples without learning transferable structure, and generalizes poorly.
  This bias-variance trade-off means that capacity and data must scale together; increasing one without the other yields diminishing returns.
  Regularization techniques such as weight decay, dropout, and data augmentation partially mitigate overfitting by constraining the hypothesis space or artificially inflating the effective training set, but they cannot substitute for real data diversity.
  The dominant finding of the deep learning era is that scaling both model capacity and dataset size simultaneously yields the most reliable accuracy gains.
  The empirical record is unambiguous: from LeNet~\cite{lecun1998gradient} on handwritten digits to AlexNet~\cite{krizhevsky2012imagenet} on large-scale images, each step change in performance has been accompanied by a corresponding increase in both model size and training data.
  The practical consequence, relevant throughout this thesis, is that high-accuracy models are invariably large and data-hungry, and both the model and the data required to produce them are costly to acquire, store, and deploy.
\end{foldable}

\subsection{Convolutional Neural Networks} \label{sec:21-background-deeplearning-cnn}
\begin{foldable}
  \Acp{CNN} are the standard architecture for image classification~\cite{lecun1998gradient}.
  Rather than applying fully connected transformations to the entire input, a convolutional layer applies a bank of small learnable filters that share weights across spatial positions.
  This inductive bias, translation equivariance and local connectivity, substantially reduces the parameter count relative to a fully connected layer of equivalent capacity, while also encoding the spatial structure of natural images: nearby pixels are more informative about each other than distant ones, and the same edge detector is useful everywhere in the image, not only at one location.
  Pooling layers subsample the resulting feature maps, progressively reducing spatial resolution while expanding the region of the input each subsequent filter can see.
  The result is a hierarchy: early layers detect low-level features such as edges and textures; later layers compose these into semantic representations such as object parts and categories.
  How far up the input this hierarchy reaches is governed by the receptive field.
\end{foldable}

\begin{foldable}
  The receptive field of a neuron is the region of the input that can influence its activation.
  In a network with only $3\times3$ convolutions and no pooling, the receptive field grows by two pixels per layer, one on each side, so reaching a large input region requires many layers.
  Pooling and strided convolutions accelerate this growth geometrically: each pooling step doubles the spatial extent that subsequent filters cover, so a modest number of layers suffices to integrate information over the full image.
  As illustrated in Figure~\ref{fig:21-background-deeplearning-cnn}, a single activation in a deep layer traces back to a progressively larger patch of the input, allowing the network to capture local texture at early layers and global spatial layout at late layers within a single end-to-end trained model.
  Building this hierarchy therefore requires depth, and depth directly drives parameter and \ac{FLOP} counts: more layers mean more filters, more weights, and more multiply-accumulate operations per forward pass.
\end{foldable}

\begin{figure}[ht]
  \centering
  \includegraphics[width=0.70\linewidth]{./figures/20-background/21-background-deeplearning-cnn-receptive.pdf}
  \caption[Receptive field growth across convolutional layers.]{
    Receptive field growth across convolutional layers.
    A single cell in Layer~3 (red) corresponds to a $3{\times}3$ region in Layer~2 (orange), which in turn maps back to a larger $5{\times}5$ region in Layer~1 (orange).
    Each additional layer expands the input region that influences any given activation, allowing deeper layers to integrate information over progressively larger spatial extents.
  }
  \label{fig:21-background-deeplearning-cnn}
\end{figure}

\begin{foldable}
  The first practical \DIFdelbegin \DIFdel{realisation }\DIFdelend \DIFaddbegin \DIFadd{realization }\DIFaddend of this design was LeNet5~\cite{lecun1998gradient}, originally trained to classify handwritten digits.
  Figure~\ref{fig:21-background-deeplearning-lenet5} shows its architecture: seven trainable layers alternating between convolutions and subsampling, beginning with 32$\times$32 input images and ending with 10 output units (one per digit class \texttt{0}--\texttt{9}).
  The network hierarchy is evident: early layers capture low-level features (edges, corners), middle layers compose these into stroke fragments, and later layers \DIFdelbegin \DIFdel{recognise }\DIFdelend \DIFaddbegin \DIFadd{recognize }\DIFaddend digit patterns.
  LeNet5 was both a proof of concept and a practical success, demonstrating that the convolutional inductive bias is a fundamental choice to vision tasks over fully connected networks.
\end{foldable}

\begin{figure}[ht]
  \centering
  \includegraphics[width=0.95\columnwidth]{./figures/20-background/21-background-deeplearning-lenet5.pdf}
  \caption[LeNet5 architecture.]{
    LeNet5 architecture for handwritten digit classification.
    Input: 32$\times$32 image.
    Alternating convolutions (C) and subsampling/pooling layers (S) progressively reduce spatial resolution while increasing feature map depth.
    The final three layers (C5, F6, F) are fully connected and produce the 10-way classification output.
    This architecture established the template for modern convolutional networks: local feature extraction followed by spatial hierarchy.
  }
  \label{fig:21-background-deeplearning-lenet5}
\end{figure}

\begin{foldable}
  LeNet-5's success motivated researchers to push beyond its modest depth.
  As larger datasets and more computational resources became available, and it was hypothesized that deeper networks could capture more abstract patterns and achieve higher accuracy on large-scale problems.
  AlexNet~\cite{krizhevsky2012imagenet} and VGG~\cite{simonyan2014very} validated this hypothesis by pushing to much deeper architectures: AlexNet with 8 trainable layers and VGG with up to 19 layers.
  Yet the deeper networks remained fundamentally harder to train: optimization slowed, convergence became fragile, and the network's ability to learn useful features in early layers degraded.
  The root cause was the \textit{vanishing gradient} problem: the gradient of the loss with respect to early-layer parameters shrinks exponentially in depth for typical weight \DIFdelbegin \DIFdel{initialisations}\DIFdelend \DIFaddbegin \DIFadd{initializations}\DIFaddend , leaving early layers with near-zero update signals and preventing them from learning.
  Careful choice of \DIFdelbegin \DIFdel{initialisation }\DIFdelend \DIFaddbegin \DIFadd{initialization }\DIFaddend and activation functions can reduce this effect, but do not eliminate it, and depth remained fundamentally limited.
\end{foldable}

\begin{foldable}
  Skip connections, introduced in ResNets~\cite{he2016deep}, resolved this by creating additive shortcut paths that allow gradients to flow directly from the loss to early layers, bypassing the multiplicative chain through intermediate layers.
  Each residual block learns only the residual $\mathcal{F}(\mathbf{x}) = \mathcal{H}(\mathbf{x}) - \mathbf{x}$ relative to its input, with the identity shortcut $\mathbf{x}$ added back at the output.
  The effective gradient magnitude at early layers then depends on the shortest path to the loss, not on total network depth.
  This structural change removed the depth bottleneck entirely: ResNet-152 (152 layers) trains successfully, and deeper variants continue to improve accuracy.
  The training failure at depth is therefore not a capacity problem but a gradient propagation problem, and skip connections are its structural solution.
\begin{foldable}

\end{foldable}
  Skip connections~\cite{he2016deep}, illustrated in Figure~\ref{fig:21-background-deeplearning-skipconnection}, resolve this structurally by creating additive shortcut paths through which gradients flow directly from the loss to early layers, bypassing the multiplicative chain.
  Each residual block learns only the residual $\mathcal{F}(\mathbf{x}) = \mathcal{H}(\mathbf{x}) - \mathbf{x}$ relative to its input, with the identity shortcut $\mathbf{x}$ added back at the output; the effective gradient magnitude at early layers then depends on the length of the shortest path to the loss, not on total network depth.
  The training failure at depth is therefore not a capacity problem but a gradient propagation problem, and skip connections are its structural solution.
\end{foldable}

\begin{figure}[t]
  \centering
  \includegraphics[width=0.75\linewidth]{./figures/20-background/21-background-deeplearning-cnn-skipconnection.pdf}
  \caption[Residual block with identity skip connection.]{
    Residual block with identity skip connection~\cite{he2016deep}.
    The input $x$ bypasses two weight layers via an identity shortcut and is added to the residual output $\mathcal{F}(x)$ before the final activation.
    This additive path allows gradients to flow directly to earlier layers, enabling stable training of very deep networks.
  }
  \label{fig:21-background-deeplearning-skipconnection}
\end{figure}

\begin{foldable}
  With this foundation in place, the architectures used in this thesis can be read as a fifteen-year progression in which each generation pushes depth further and pays a higher parameter cost to do so.
  LeNet5~\cite{lecun1998gradient} (two convolutional layers, 60\,K parameters) established the core template, convolutional feature extraction followed by fully connected classification, and remains a standard benchmark on simple datasets.
  AlexNet~\cite{krizhevsky2012imagenet} scaled this template to five convolutional layers and 60\,M parameters, working on slightly harder datasets.
  VGGNet~\cite{simonyan2014very} and ResNet~\cite{he2016deep} are deeper and larger, suited to more complex datasets.
  These architectures remain the most widely used teacher architectures in the \ac{KD} literature and appear throughout Chapters~\ref{ch:30-research01} and~\ref{ch:40-research02}.
\end{foldable}

\begin{foldable}
  Accuracy scales smoothly with depth, width, and input resolution across the \ac{CNN} family, but parameter and compute costs scale faster than accuracy gains.
  Doubling depth increases parameter count roughly linearly, yet the representational gains and the compute required to \DIFdelbegin \DIFdel{realise }\DIFdelend \DIFaddbegin \DIFadd{realize }\DIFaddend them grow more than proportionally for many architectures, so the gap between what accuracy demands and what deployment allows widens as task difficulty increases.
  The need for more efficient deployable models holds regardless of where a model sits on this spectrum.
  This is the fundamental motivation for model compression, taken up in \S\ref{sec:22-background-compression}.
\end{foldable}

\subsection{Foundational Natural Language Processing} \label{sec:21-background-deeplearning-foundationalnlp}
\begin{foldable}
  Just as vision progressed from hand-crafted filters to learned convolutional hierarchies, language processing followed a parallel arc, from hand-crafted rules through statistical learning to distributed neural representations.
  These are driven by the same underlying force: scale and data consistently outperform manual design.
\end{foldable}

\begin{foldable}
  Early \ac{NLP} was built on hand-crafted rules~\cite{chomsky1956three}.
  Symbolic systems encoded linguistic knowledge as explicit grammars, lexicons, and inference rules, producing interpretable pipelines capable of handling carefully scoped tasks such as syntactic parsing or limited-domain question answering~\cite{winograd1971procedures}.
  The core limitation was scalability: every new domain or linguistic phenomenon required manual engineering, and the resulting systems were brittle outside the narrow conditions they were designed for~\cite{manning1999foundations}.
  The coverage that real language understanding demands cannot be manually encoded.
\end{foldable}

\begin{foldable}
  The shift to statistical methods replaced hand-crafted rules with data-driven models.
  Language models based on $n$-gram statistics~\cite{shannon1948mathematical}, sequence labelers based on hidden Markov models~\cite{rabiner2002tutorial}, and structured predictors based on conditional random fields~\cite{lafferty2001conditional} each brought a different inductive bias.
  They shared a common workflow: hand-engineered features were extracted from text and fed to a probabilistic model trained to \DIFdelbegin \DIFdel{maximise }\DIFdelend \DIFaddbegin \DIFadd{maximize }\DIFaddend likelihood on labeled corpora.
  Performance improved substantially over symbolic baselines, and the data-driven approach proved far more portable across domains.
  The persistent limitation, however, was representation: features were sparse, high-dimensional, and task-specific, so each new task required a fresh feature-engineering effort and no knowledge transferred between tasks~\cite{bengio2003neural}.
\end{foldable}

\begin{foldable}
  Neural methods dissolved this limitation by learning representations directly from raw text.
  Word2Vec~\cite{mikolov2013distributed} and GloVe~\cite{pennington2014glove} showed that dense, low-dimensional word embeddings trained on large unlabeled corpora encode meaningful semantic and syntactic relationships, transferring usefully to downstream tasks without task-specific feature design.
  \Acp{RNN} and later long-short term memory networks~\cite{hochreiter1997long}, extended this to variable-length sequences, enabling models to read a sentence left to right and accumulate a hidden state.
  Seq2Seq architectures~\cite{sutskever2014sequence} added a decoder, producing the encoder-decoder template used for machine translation and \DIFdelbegin \DIFdel{summarisation}\DIFdelend \DIFaddbegin \DIFadd{summarization}\DIFaddend .
  Two limitations remained: vanishing gradients over long sequences caused the hidden state to lose early context, and the fixed-length encoder vector was forced to compress an entire input sentence into a single representation regardless of length, a design that degraded on long inputs.
\end{foldable}

\begin{foldable}
  All of these models operate on discrete token sequences, making the choice of tokenization a foundational design decision.
  Early systems tokenized on whitespace or characters; modern systems use subword methods such as byte-pair encoding~\cite{sennrich2016neural} and WordPiece~\cite{wu2016google}, which segment words into frequent subword units and thereby handle rare and out-of-vocabulary words gracefully.
  The choice of tokenizer is not neutral: different segmentations partition the vocabulary differently, and any method that operates on token-level statistics, including the importance-scoring procedure in Chapter~\ref{ch:50-research03}, inherits the systematic biases of the tokenizer it is applied to.
  The trajectory of this era traces a clear arc: from manually encoded rules, through hand-crafted statistical features, to fully learned distributed representations. Yet every model in this lineage remained task-specific, and no unified architecture capable of general-purpose language understanding had appeared.
  That unification came with the Transformer.
\end{foldable}

\subsection{The Transformer Era} \label{sec:21-background-deeplearning-transformer}
\begin{foldable}
  The Transformer~\cite{vaswani2017attention} resolved the sequential limitation of recurrent models by replacing recurrence entirely with self-attention.
  Rather than reading a sequence token by token, self-attention computes pairwise interactions between all positions simultaneously: each token attends to every other token, weighted by a learned compatibility score, and aggregates their representations accordingly.
  Multi-head attention runs several such attention functions in parallel, allowing the model to jointly attend to information from different representation subspaces.
  Positional encodings are added to the input embeddings to supply the order information that recurrence formerly provided implicitly.
  Each Transformer block pairs multi-head attention with a feed-forward network (two fully-connected layers with a nonlinearity between them), and residual connections around both sub-layers \DIFdelbegin \DIFdel{stabilise }\DIFdelend \DIFaddbegin \DIFadd{stabilize }\DIFaddend training and preserve gradient flow.
  The result is a model that processes sequences in parallel, eliminates the vanishing-gradient-over-distance problem, and scales far more efficiently on modern hardware than any recurrent alternative.
  Figure~\ref{fig:21-background-deeplearning-attention} illustrates the full encoder-decoder architecture with these components stacked in sequence.
\end{foldable}

\begin{figure}[ht]
  \centering
  \includegraphics[width=0.5\columnwidth]{./figures/20-background/21-background-deeplearning-attention.png}
  \caption[The Transformer model architecture]{
    Transformer encoder-decoder architecture showing multi-head attention blocks, feed-forward networks, positional encodings, and residual connections (Add \& Norm) stacked in sequence.
  }
  \label{fig:21-background-deeplearning-attention}
\end{figure}

\begin{foldable}
  The Transformer's parallelism made large-scale pre-training practical, and pre-training transformed the field.
  Rather than training a model from scratch for each task, the paradigm trains a single large model on a massive unlabeled corpus and then fine-tunes it on a small task-specific dataset.
  Three architecture families followed, distinguished by the attention mask used during pre-training.
  Encoder-only models such as BERT~\cite{devlin2019bert} apply bidirectional attention and are pre-trained via masked language \DIFdelbegin \DIFdel{modelling}\DIFdelend \DIFaddbegin \DIFadd{modeling}\DIFaddend , making them well-suited to understanding tasks such as classification and named entity recognition.
  Decoder-only models such as GPT~\cite{radford2018improving} apply causal (left-to-right) attention and are pre-trained via next-token prediction, making them natural generators.
  Encoder-decoder models such as T5~\cite{raffel2020exploring} combine both, handling sequence-to-sequence tasks such as translation and \DIFdelbegin \DIFdel{summarisation}\DIFdelend \DIFaddbegin \DIFadd{summarization}\DIFaddend .
  This distinction is not merely taxonomic: Chapter~\ref{ch:50-research03} operates on encoder-only models fine-tuned for classification, and the dataset distillation method it proposes is designed specifically around token-level statistics produced by this architecture class.
\end{foldable}

\begin{foldable}
  Fine-tuned pre-trained models are evaluated across a range of downstream tasks that probe different aspects of language understanding.
  Text classification assigns a label to a document or sentence; question answering requires locating or generating an answer given a passage.
  Whereas, \ac{NLI} asks whether a hypothesis is entailed by, contradicts, or is neutral with respect to a premise.
  \ac{NLI} is particularly useful as an evaluation task because it requires the model to represent semantic relationships between sentence pairs rather than individual sentences, and its three-way label space provides a discriminative test of fine-grained understanding.
  These tasks appear in Chapter~\ref{ch:50-research03} as the benchmarks on which dataset distillation for text is evaluated.
\end{foldable}

\begin{foldable}
  Neural scaling laws showed that Transformer loss decreases smoothly and predictably as model size, dataset size, and compute increase together~\cite{kaplan2020scaling}.
  GPT-2~\cite{radford2019language} demonstrated coherent long-form generation at 1.5\,B parameters, and GPT-3~\cite{brown2020language} solved many tasks from in-context examples alone at 175\,B.
  At sufficient scale, models began exhibiting qualitatively new \DIFdelbegin \DIFdel{behaviours }\DIFdelend \DIFaddbegin \DIFadd{behaviors }\DIFaddend (in-context learning, multi-step reasoning, and few-shot task adaptation) that were absent in smaller counterparts and not predictable by extrapolation, often called emergent abilities~\cite{wei2022emergent}.
  Models soon crossed the 1\,T mark with the help of mixture-of-experts architectures~\cite{fedus2022switch}, but each step up this ladder multiplied training and serving costs, making frontier models accessible only to a small number of \DIFdelbegin \DIFdel{organisations}\DIFdelend \DIFaddbegin \DIFadd{organizations}\DIFaddend .
  The practical response has been to train a large model and distil it into a compact, deployable variant.
  This workflow places distillation, not scale itself, at the \DIFdelbegin \DIFdel{centre }\DIFdelend \DIFaddbegin \DIFadd{center }\DIFaddend of the deployment problem this thesis addresses.
\end{foldable}
 \section{From Compression to Distillation} \label{sec:22-background-compression}
\subsection{Model Compression} \label{sec:22-background-compression-model}
\begin{foldable}
  The gap between the accuracy achievable by large neural networks and the computational budget available at deployment has motivated a substantial body of work on model compression.
  The field broadly divides into four families: pruning, quantization, low-rank factorization, and \acl{KD}.
  Each makes a different trade-off between compression ratio, hardware portability, and the degree of access required to the teacher model, and this last dimension turns out to be the most consequential in practice.
  Figure~\ref{fig:22-background-compression-modelcompress} illustrates these four families and their relationship to the original teacher.
\end{foldable}

\begin{figure}[ht]
  \centering
  \includegraphics[width=0.85\columnwidth]{./figures/20-background/22-background-compression-modelcompress.pdf}
  \caption[Four model compression families.]{
    Four model compression families.
    (Left) Teacher network: the original large model.
    (Top middle) Pruning: removes nodes and connections, indicated by dashed circles.
    (Top right) Quantization: reduces numerical precision, represented by square nodes.
    (Bottom middle) Low-rank adaptation: factors weight matrices into lower-rank components, shown as trapezoidal modules.
    (Bottom right) Knowledge distillation: trains a separate compact student network (blue) to match teacher outputs.
  }
  \label{fig:22-background-compression-modelcompress}
\end{figure}

\begin{foldable}
  \textbf{Pruning} removes parameters or entire structural units from a trained network~\cite{han2015learning}.
  Unstructured pruning zeros individual weights below a magnitude threshold, yielding sparse weight tensors.
  Over 90\% of parameters in AlexNet and VGG can be pruned with negligible accuracy loss when combined with retraining.
  However, irregular sparsity does not translate directly to wall-clock speedups without \DIFdelbegin \DIFdel{specialised }\DIFdelend \DIFaddbegin \DIFadd{specialized }\DIFaddend sparse hardware or libraries.
  Structured pruning removes whole filters, channels, or layers, producing dense sub-networks that accelerate inference on commodity hardware~\cite{molchanov2016pruning}.
  The coarser granularity typically requires a larger pruning ratio to achieve the same compression as unstructured pruning.
  Pruning requires direct access to the network's weight matrices and structure.
\end{foldable}

\begin{foldable}
  \textbf{Quantization} reduces the numerical precision of weights and activations~\cite{jacob2018quantization}.
  Fixed-point (\texttt{INT8}) quantization is now standard in production inference pipelines and is directly supported by most modern hardware accelerators.
  It typically recovers near-full-precision accuracy with a four-fold memory reduction.
  Extreme approaches such as binary neural networks represent each weight as a single bit~\cite{hubara2018quantized}, achieving large memory savings at the cost of accuracy on complex tasks.
  Quantization demands access to internal weight representations and activation statistics.
\end{foldable}

\begin{foldable}
  \textbf{Low-rank factorization} decomposes weight matrices into products of smaller matrices~\cite{hu2022lora}.
  For example, $W \approx UV^{\top}$ where $\text{rank}(U) \ll \min(m,n)$.
  This reduces parameter count and arithmetic cost at the price of an approximation error.
  The technique is most effective for large fully connected and embedding layers, where a low-rank approximation captures most of the variance.
  It is less suitable for the many small convolutional filters in modern CNNs, where the weight tensors are not well approximated by low-rank structures and the approximation error accumulates across layers.
  Low-rank factorization, like the two preceding methods, requires full access to the model's weight matrices.
\end{foldable}

\begin{foldable}
  \textbf{\Acl{KD}} takes quite a different approach~\cite{hinton2015distilling}.
  Rather than modifying the teacher's parameters, it trains a separate, compact student network to reproduce the teacher's output predictions.
  The student is \DIFdelbegin \DIFdel{initialised }\DIFdelend \DIFaddbegin \DIFadd{initialized }\DIFaddend independently and never reads the teacher's weights, gradients, or internal activations.
  The only interface it uses is the teacher's output, which is a vector of class probabilities for each input query.
  This output-only interface is what makes \ac{KD} compatible with settings where the teacher is accessible only as a black-box \ac{API}.
\end{foldable}

\begin{foldable}
  The access requirements of the four families draw a sharp boundary.
  Pruning, quantization, and low-rank factorization all operate on the model's internal weight representation: they read weight matrices, modify them, and require at least one subsequent pass of calibration or retraining.
  This is not an incidental implementation detail but a fundamental prerequisite: without direct access to parameters and activations, none of these transformations can be applied.
  Any deployment setting in which the model is accessible only as an opaque input-output function rules out all three simultaneously.
  Yet \ac{KD} does not cross this boundary: it queries the teacher, observes its outputs, and trains the student entirely from those observations.
  When the teacher is a black-box, \ac{KD} is not a preference; it is the only viable path.
\end{foldable}

\subsection{Dataset Compression} \label{sec:22-background-compression-dataset}
\begin{foldable}
  Meanwhile, dataset compression takes the complementary route, reducing the volume of training data required to reach a given model quality while leaving the model architecture unchanged.
  Three families of dataset compression have been proposed, distinguished by whether they retain real samples and how they measure representativeness.
  Figure~\ref{fig:22-background-compression-datasetcompress} illustrates these three families.
\end{foldable}

\begin{figure}[ht]
  \centering
  \includegraphics[width=0.99\columnwidth]{./figures/20-background/22-background-compression-datasetcompress.pdf}
  \caption[Three dataset compression families.]{
    Three dataset compression families.
    (Left) Prototype centroids (orange) summarize each cluster from original dataset (blue).
    (Middle) Coreset selects only a subset of representative real samples.
    (Right) Dataset distillation optimizes synthetic samples to match the original distribution.
  }
  \label{fig:22-background-compression-datasetcompress}
\end{figure}

\begin{foldable}
  \textbf{Prototype and clustering methods} summarize the dataset by a set of representative centroids.
  Algorithms such as $k$-means~\cite{macqueen1967some} and Gaussian mixture models~\cite{dempster1977maximum} partition the training data and represent each partition by a single centroid.
  Training proceeds on these centroids rather than the original data, yielding a compact surrogate that captures the global cluster structure.
  The approach is computationally cheap and scales to large datasets, but it treats each sample as interchangeable within its cluster and therefore ignores which samples are most informative for the learning objective.
\end{foldable}

\begin{foldable}
  \textbf{Coreset selection} retains a subset of real training samples chosen so that a model trained on the subset closely approximates one trained on the full dataset.
  Selection criteria are grounded in the learning objective: geometry-based methods select samples that span the input space~\cite{sener2018active}, while gradient-based methods select samples whose gradients collectively best approximate the full-data gradient~\cite{mirzasoleiman2020coresets}.
  Because coresets consist of original samples, they preserve the distributional properties of the data and require no generative process.
  The principal limitation is that the selected subset is constrained to what already exists in the dataset; no sample can be more informative than the most informative real example.
\end{foldable}

\begin{foldable}
  \textbf{Dataset distillation}~\cite{wang2018dataset} takes a different approach.
  Rather than selecting or \DIFdelbegin \DIFdel{summarising }\DIFdelend \DIFaddbegin \DIFadd{summarizing }\DIFaddend existing samples, it optimizes a small set of synthetic samples end-to-end so that models trained on the distilled set match the performance of models trained on the full dataset.
  The synthetic samples are not drawn from any real distribution; they are free parameters jointly optimized with a surrogate training objective, typically by matching model parameters or gradients~\cite{zhao2021dataset} between the distilled set and the full dataset.
  Because the distilled set is not constrained to the support of the original data, it can be far more compact than any coreset.
\end{foldable}

\begin{foldable}
  The three families differ not only in mechanism but in what they preserve.
  Prototype methods preserve cluster geometry; coreset methods preserve the loss or gradient surface of the full dataset; dataset distillation preserves model \DIFdelbegin \DIFdel{behaviour}\DIFdelend \DIFaddbegin \DIFadd{behavior}\DIFaddend , specifically end-to-end performance on the target task.
  All three act on the training signal rather than the learned parameters, distinguishing them collectively from model compression.
  Prototype and coreset methods are bounded by the real samples that exist: no selected or summarized sample can carry more knowledge than the most informative real example.
  When the corpus is insufficient or biased, selection alone cannot recover coverage; dataset distillation, unconstrained to the original data support, is the only family that can.
\end{foldable}

\subsection{Why Distillation?} \label{sec:22-background-compression-whydistill}
\begin{foldable}
  This thesis adopts a distillation approach to both compression axes.
  Throughout, the term \textbf{\acl{KD}} serves as an umbrella encompassing both model distillation (compressing a learned model) and dataset distillation (compressing a training corpus): both solve compression by transferring knowledge from a richer source to a more constrained target, one on the model axis and one on the data axis.
  The explicit distinction between model and dataset distillation will be invoked where the two paradigms diverge mechanically or in scope; where context renders it unambiguous, the umbrella term is used.
\end{foldable}

\begin{foldable}
  Within model compression, model distillation is uniquely \textit{architecture-agnostic}: the student is designed independently and can differ from the teacher in architecture, depth, or model family, whereas pruning and quantization operate on the teacher's existing weight structure and produce compressed variants of it.
  Beyond this, model distillation \textit{does not require internal access to the source}: the teacher is a fixed oracle queried through its outputs alone, and its weights, activations, and gradients are never read.
  Pruning, quantization, and low-rank factorization each depend on exactly this access (reading weight matrices, computing activation statistics, calibrating layer-wise ranges), so any setting where the teacher is exposed only as an \ac{API} eliminates all three simultaneously.
  Most importantly, model distillation \textit{preserves learned knowledge} rather than surface proxies: the teacher's output probability vector encodes not only the predicted class but the relational structure among all classes, the so-called dark knowledge~\cite{hinton2015distilling} absent from a hard label.
  A one-hot label records which class won; a soft output records by how much and how every other class ranked, a signal that hard-label training permanently discards at no savings in annotation cost.
\end{foldable}

\begin{foldable}
  On the data axis, the alternatives fail for analogous reasons.
  Prototype methods summarize the dataset by geometric centroids in input space: a centroid may be Euclidean-representative while bearing no relationship to the decision boundary, the analogue of a hard label that records position without recording importance.
  Coreset selection grounds selection in the learning objective and is strictly more informative, but remains bounded by the existing data: no selected sample can carry more knowledge than the most informative real example, and selection can only redistribute importance across the existing support, not extend it.
  Dataset distillation addresses both shortcomings: synthetic parameters are \textit{not bounded by the original data}, extending beyond the existing support.
  \DIFdelbegin \DIFdel{Optimising }\DIFdelend \DIFaddbegin \DIFadd{Optimizing }\DIFaddend against the training objective \textit{preserves learned knowledge} rather than surface proxies.
  Both failures thus have a common cause: compression that ignores the learning objective cannot faithfully represent the original distribution.
\end{foldable}

\begin{foldable}
  Both verdicts compound into a single principle.
  Model and dataset distillation survive their respective access restrictions (black-box \acp{API} on the model axis, insufficient or biased corpora on the data axis) because they treat compression as a \textbf{knowledge transfer} problem rather than a parameter reduction problem.
  The non-distillation alternatives fail for the same structural reason: hard labels discard inter-class relational structure; geometric summaries discard task-relevant importance; coreset selection is bounded by the existing distribution.
  The shared failure mode, and the shared design objective, is \textbf{distributional fidelity}: whether the compact \DIFdelbegin \DIFdel{artefact}\DIFdelend \DIFaddbegin \DIFadd{artifact}\DIFaddend , a student model or surrogate corpus, faithfully represents the original distribution under practical constraints.
  The practical constraints this thesis addresses, limited real data, black-box teacher access, and discrete token modality, each degrade distributional fidelity in a specific way, and recovering it under each is the unifying objective of Chapters~\ref{ch:30-research01} through~\ref{ch:50-research03}.
\end{foldable}
 \section{Model Distillation} \label{sec:23-background-modeldistill}
\subsection{The Original Formulation} \label{sec:23-background-modeldistill-hinton}
\begin{foldable}
  The idea of training a smaller student network to mimic a larger teacher originated with model compression~\cite{bucila2006model}.
  This was formalized as \textit{knowledge distillation} (\ac{KD})~\cite{hinton2015distilling}, proposing that students learn not from hard labels but from soft probabilistic targets that encode the teacher's learned structure.
  Model distillation trains a compact \textit{student} network $S$ to replicate the \DIFdelbegin \DIFdel{behaviour }\DIFdelend \DIFaddbegin \DIFadd{behavior }\DIFaddend of a large, pre-trained \textit{teacher} network $T$, illustrated in Figure~\ref{fig:23-background-modeldistill-overview}.
  The goal is to transfer the teacher's learned \textit{decision-boundary structure} and relational knowledge, the structured understanding of the task and its geometry acquired over the full training budget, rather than requiring pointwise input-output fidelity on a fixed set of training examples.
  This distinguishes distillation from model approximation, which targets pointwise structural matching, and from label smoothing, which perturbs labels uniformly without reference to the teacher's learned geometry.
\end{foldable}

\begin{figure}[ht]
  \centering
  \includegraphics[width=0.80\linewidth]{./figures/20-background/23-background-modeldistill-overview.pdf}
  \caption[The original knowledge distillation formulation.]{
    The original knowledge distillation formulation: a student network is trained on the same input data as the teacher, learning to match the teacher's soft output distribution in addition to the hard labels alone.
  }
  \label{fig:23-background-modeldistill-overview}
\end{figure}

\begin{foldable}
  The obstacle is that standard supervised training discards the very information distillation needs.
  A one-hot label records which class won the annotation decision and assigns zero probability to every other class.
  A large network trained on such labels does not merely learn to classify correctly; it implicitly acquires a structured posterior over the class space, shaped by the statistical regularities of the full training distribution.
  When such a network evaluates an image of a cat, it assigns high probability to \texttt{cat}, a smaller amount to \texttt{dog}, and a trace to \texttt{leopard}, while assigning essentially nothing to \texttt{car}, encoding which classes share visual structure and which do not.
  Similarly, a digit \texttt{8} receives non-trivial mass on \texttt{3} and near-zero on \texttt{1}, reflecting the perceptual similarity the training data induced.
  This inter-class relational structure, termed \textit{dark knowledge}~\cite{hinton2015distilling}, is present in every forward pass of the teacher yet discarded by hard-label supervision.
\end{foldable}

\begin{foldable}
  \Acl{KD} recovers dark knowledge by training the student to \DIFdelbegin \DIFdel{minimise }\DIFdelend \DIFaddbegin \DIFadd{minimize }\DIFaddend the \ac{KL} divergence between the teacher's output distribution $q$ and the student's $p$:
  \begin{equation}
    \mathcal{L}_{\text{KD}} = \mathcal{L}_\text{KL}(q \| p).
    \label{eq:kd-loss-core}
  \end{equation}
  In practice, a highly confident teacher concentrates nearly all mass on the correct class, making dark knowledge numerically invisible even though it exists.
  \textit{Temperature scaling} addresses this: a parameter $\tau \geq 1$ applied to the pre-softmax logits softens the distribution and amplifies inter-class structure:
  \begin{equation}
    q_k^{(\tau)} = \frac{\exp(z_k^{T} / \tau)}{\sum_{j=1}^{K} \exp(z_j^{T} / \tau)}.
    \label{eq:softmax-temperature}
  \end{equation}
  A hard-label cross-entropy term is added to keep the student anchored to ground truth, yielding the full objective~\cite{hinton2015distilling}:
  \begin{equation}
    \mathcal{L}_{\text{KD}}
    = \alpha \tau^{2} \mathcal{L}_\text{KL}(q^{(\tau)} \| p^{(\tau)})
    + (1 - \alpha) \mathcal{L}_{\text{CE}}(p^{(1)}, y),
    \label{eq:kd-loss}
  \end{equation}
  where $\alpha \in [0,1]$ balances the two terms and $\tau^{2}$ restores gradient scale after softening.
\end{foldable}

\begin{foldable}
  The teacher's output is not merely a prediction: it is a compressed summary of everything the teacher learned from the full training distribution, encoded in a probability vector.
  Training the student against this posterior transfers the teacher's generalization directly: its responses on hard examples, its calibrated uncertainty, and its implicit encoding of class relationships.
  Empirically, a student of identical architecture to the teacher consistently outperforms the same model trained on hard labels alone~\cite{hinton2015distilling}, confirming that the gain comes from signal quality rather than capacity.
  Subsequent work has shown that the teacher's intermediate representations and the relational structure among samples carry additional signal, motivating a broader taxonomy of distillation methods.
\end{foldable}

\subsection{Taxonomy of Model Distillation Methods} \label{sec:23-background-modeldistill-taxonomy}
\begin{foldable}
  The original formulation transfers knowledge through the teacher's output distribution, but the field has since expanded to exploit richer sources of teacher signal.
  The design space of \ac{KD} decomposes along two orthogonal dimensions: (1) the \textit{source of knowledge} transferred (response-based, feature-based, or relation-based methods), and (2) the \textit{training scheme} that structures the teacher-student relationship (offline, online, self, or co-distillation).
  We first inventory the knowledge sources by examining what each family transfers and what teacher access is required, then discuss training schemes.
  Together they \DIFdelbegin \DIFdel{characterise }\DIFdelend \DIFaddbegin \DIFadd{characterize }\DIFaddend the full design space.
\end{foldable}

\begin{figure}[ht]
  \centering
  \includegraphics[width=0.70\linewidth]{./figures/20-background/23-background-modeldistill-responsekd.pdf}
  \caption[Response-based knowledge distillation.]{
    Response-based distillation matches the teacher's output distribution, requiring only output-level access.
  }
  \label{fig:23-background-modeldistill-responsekd}
\end{figure}

\begin{foldable}
  \textbf{Response-based distillation}~\cite{hinton2015distilling} trains the student to match the teacher's output class distribution:
  \begin{equation}
    \mathcal{L}_{\text{KD}}^{\text{response}} = \alpha \tau^2 \mathcal{L}_{\text{KL}}(q^{(\tau)} \| p^{(\tau)}) + (1-\alpha) \mathcal{L}_{\text{CE}}(p, y),
    \label{eq:kd-response}
  \end{equation}
  as formalized in Equation~\ref{eq:kd-loss}.
  Variants include regression on the raw un-softened logit vector and born-again networks, in which a same-architecture student is sequentially re-distilled from the previous generation, yielding ensembles that outperform the original teacher~\cite{ba2014deep}.
  Response-based methods require only the teacher's output probability vector and are therefore compatible with any setting in which that vector is accessible, including black-box \acp{API}.
\end{foldable}

\begin{figure}[ht]
  \centering
  \includegraphics[width=0.65\linewidth]{./figures/20-background/23-background-modeldistill-featurekd.pdf}
  \caption[Feature-based knowledge distillation.]{
    Feature-based distillation matches intermediate layer representations via a learned projector, requiring internal teacher access.
  }
  \label{fig:23-background-modeldistill-featurekd}
\end{figure}

\begin{foldable}
  \textbf{Feature-based distillation} transfers knowledge from the teacher's intermediate representations by matching aligned feature maps at selected layers:
  \begin{equation}
    \mathcal{L}_{\text{KD}}^{\text{feature}} = \left\| W f^T - f^S \right\|^2,
    \label{eq:kd-feature}
  \end{equation}
  where $f^T$ and $f^S$ are teacher and student feature maps at a chosen intermediate layer, and $W$ is a learned projection matrix.
  This loss can be summed over multiple selected layer pairs.
  FitNets~\cite{romero2015fitnets} introduced matching intermediate hidden layers via a thin regressor, enabling students deeper and thinner than the teacher.
  Attention Transfer~\cite{zagoruyko2016paying} distils spatial attention maps, capturing \textit{where} the teacher focuses rather than \textit{what} it predicts.
  PKT~\cite{passalis2018learning} matches kernel density estimates of teacher and student feature distributions, encoding distributional geometry rather than individual activation values.
  Feature-based methods consistently improve on response-based baselines, but their incompatibility with opaque \acp{API} is fundamental.
  Without access to the teacher's activation dimensionality and representation geometry, the projector bridging teacher and student spaces is entirely ill-posed: there is no anchor in the teacher's representation space to which the student can be steered.
\end{foldable}

\begin{figure}[ht]
  \centering
  \includegraphics[width=0.85\linewidth]{./figures/20-background/23-background-modeldistill-relationkd.pdf}
  \caption[Relation-based knowledge distillation.]{
    Relation-based distillation preserves pairwise distance and angle relationships between samples in the teacher's embedding space.
  }
  \label{fig:23-background-modeldistill-relationkd}
\end{figure}

\begin{foldable}
  \textbf{Relation-based distillation} transfers higher-order structure: pairwise or group-wise relationships among samples in the teacher's embedding space:
  \begin{equation}
    \mathcal{L}_{\text{KD}}^{\text{rel}} = \sum_{i,j} \left[ (d_T(f_i^T, f_j^T) - d_S(f_i^S, f_j^S))^2 + (\Delta\theta_T - \Delta\theta_S)^2 \right],
    \label{eq:kd-relation}
  \end{equation}
  where $f_i^T$ and $f_i^S$ are feature vectors for samples $i$ and $j$ extracted from an intermediate layer, $d$ measures pairwise Euclidean distance, and $\Delta\theta$ is the angle between feature vectors, with subscripts $T$ and $S$ denoting teacher and student respectively.
  RKD~\cite{park2019relational} penalizes discrepancies in distance and angle relationships between sample pairs, arguing that relational information is more transferable across architectures than pointwise feature values.
  Contrastive distillation methods extend this by using self-supervised objectives to preserve the teacher's instance discrimination structure.
  Relation-based methods therefore share the access problem of feature-based methods, as pairwise statistics require intermediate-level teacher outputs.
  Beyond access, a student trained solely on hard labels develops an embedding geometry incommensurable with the teacher's, making the relational target unreachable.
\end{foldable}

\begin{foldable}
  The second axis \DIFdelbegin \DIFdel{organises }\DIFdelend \DIFaddbegin \DIFadd{organizes }\DIFaddend methods by the \textit{training scheme}: the relationship between teacher and student during training, illustrated in Figure~\ref{fig:23-background-modeldistill-schemes}.
  \textit{Offline distillation} uses a fixed, pre-trained teacher queried through its outputs; the student is trained entirely from the teacher's frozen predictions.
  \textit{Online distillation} trains teacher and student simultaneously, allowing the teacher to evolve during the student's training.
  \textit{Self-distillation} requires no separate teacher: the student distils knowledge from its own earlier or shallower representations.
  \textit{Co-distillation} extends this to an ensemble of students that mutually supervise one another.
  All work in this thesis operates in the offline regime.
\end{foldable}

\begin{figure}[ht]
  \centering
  \includegraphics[width=0.99\linewidth]{./figures/20-background/23-background-modeldistill-schemes.pdf}
  \caption[Knowledge distillation training schemes.]{
    Knowledge distillation training schemes: offline (frozen pre-trained teacher), online (jointly trained), self-distillation, and co-distillation.
  }
  \label{fig:23-background-modeldistill-schemes}
\end{figure}

\begin{foldable}
  Across all three source-of-knowledge families, two conditions are implicit: the full training dataset $\mathcal{D}$ is available to supply inputs to the teacher, and the teacher exposes at least its output probability vector.
  When either condition fails, the picture changes dramatically.
  Response-based methods' black-box compatibility becomes their sole advantage when these assumptions break: feature-based and relation-based methods become inaccessible, while response-based distillation survives, albeit degraded, under hard-label-only settings.
  The next subsection formalizes these two assumptions and \DIFdelbegin \DIFdel{characterises }\DIFdelend \DIFaddbegin \DIFadd{characterizes }\DIFaddend what happens when they fail.
\end{foldable}

\subsection{Standard Assumptions and Their Limitations} \label{sec:23-background-modeldistill-assumptions}
\begin{foldable}
  \textbf{Assumption 1: the full training dataset is available.}
  In many realistic settings, the original training data is inaccessible: it may contain private records (medical imaging, financial transactions), it may be proprietary to the model provider, or it may have been discarded after training to reduce storage costs.
  Data regulations such as \ac{GDPR}~\cite{eu2016gdpr} or \ac{HIPAA}~\cite{us1996hipaa} further restrict the retention and transfer of training sets across organizational boundaries.
  When $\mathcal{D}$ is unavailable, the student must be trained on a surrogate: a small held-out subset, synthetic data from a generative model, or public proxy data from a related domain.
  This constraint exists on a spectrum, from \textit{few-shot} settings in which only a handful of samples are available to fully \textit{data-free} settings in which no original data exists; the data-free end is strictly harder, as no real distribution signal remains.
\end{foldable}

\begin{foldable}
  \textbf{Assumption 2: the teacher is accessible at the logit level or deeper.}
  All three \ac{KD} families require at minimum the teacher's output probability vector over all $K$ classes; feature-based and relation-based methods additionally require internal activations.
  Modern production models are routinely deployed behind \acp{API} that return only a top-1 predicted label, or with associated confidence scores, to protect proprietary model internals and limit adversarial probing.
  In this setting, the access restrictions eliminate two of the three families outright: feature-based and relation-based methods are impossible without internal activations.
  Response-based methods remain the sole viable family, though they are severely degraded by the loss of soft probability distributions: a single hard label carries $\log_2 K$ bits per query, whereas the teacher's full soft target encodes the inter-class similarity structure it learned.
\end{foldable}

\begin{foldable}
  Together, these two failures \DIFdelbegin \DIFdel{characterise }\DIFdelend \DIFaddbegin \DIFadd{characterize }\DIFaddend the landscape of practical \ac{KD}.
  Absent original data, the distillation signal loses coverage of the true input distribution; absent soft targets, it loses the relational structure that distinguishes \ac{KD} from ordinary supervised training.
  \S\ref{sec:24-background-kdrestricted} surveys these two restriction axes, frames distributional fidelity as the binding constraint, and identifies the structural gaps that motivate Chapters~\ref{ch:30-research01}--\ref{ch:50-research03}.
\end{foldable}
 \section{Model Distillation Under Restricted Access} \label{sec:24-background-kdrestricted}

Two independent dimensions of restriction define the practical setting of model distillation: teacher access (white-box $\to$ black-box) and data availability (full $\to$ few-shot $\to$ data-free).
This section tightens both dimensions in turn, showing how each constraint plays out, and narrows to the two most challenging and practically motivated settings: black-box few-shot \ac{KD} (Chapter~\ref{ch:30-research01}) and black-box data-free \ac{KD} (Chapter~\ref{ch:40-research02}).
The \DIFdelbegin \DIFdel{organisation }\DIFdelend \DIFaddbegin \DIFadd{organization }\DIFaddend is pedagogical: we first survey white-box distillation across all data regimes, setting a baseline for how data constraints behave when teacher internals remain accessible; we then remove white-box access entirely and survey the black-box regime across data-availability scales.

\subsection{White-Box Model Distillation} \label{sec:24-background-kdrestricted-whitebox}
\begin{foldable}
  White-box teacher access is a natural extension of standard \ac{KD}~\cite{hinton2015distilling}: the student has full access to teacher internals (weights, gradients, activations, probabilities), enabling feature-based and relation-based methods in addition to response-based distillation.
  That situation gives the baseline case: white-box access with a full, diverse training dataset.
  This subsection examines what happens as data availability is reduced step by step, from few-shot to proxy/in-the-wild datasets to fully data-free, while white-box access remains intact.
  The progression shows how each data constraint plays out, before removing white-box access entirely and surveying the qualitatively different problems that arise in the black-box regime.
\end{foldable}

\begin{foldable}
  In the few-shot white-box setting, the student is trained on a small labeled subset, and the teacher returns its full probability vector for any queried image.
  Standard augmentation techniques (rotation, flipping, cropping, \etc) extend the training set within the convex hull of the few-shot examples, encouraging smoother decision boundaries~\cite{shorten2019survey}.
  Since the teacher remains white-box, the student receives rich per-sample soft supervision that partially compensates for the limited data coverage.
  The dual channel, informative soft targets paired with modest data coverage, yields a trade-off: the distillation signal quality is high, but its scope is inherently narrow.
\end{foldable}

\begin{foldable}
  When the original training set is unavailable, a related dataset from a different domain can substitute~\cite{chen2021learning, tang2023distribution}.
  White-box access preserves the teacher's soft targets, the probability distributions over all classes that encode the inter-class structure the teacher learned.
  However, proxy datasets introduce a domain mismatch bias: a model trained on ImageNet transfers poorly to medical imaging, and even within-domain shifts (\eg, synthetic images as a proxy for real-world deployment) distort the student's learned decision surface~\cite{torralba2011unbiased}.
  The student recovers a model that approximates the teacher's posterior on the proxy distribution, not the original.
  This bias is unavoidable when the true data is inaccessible; white-box access cannot overcome a domain mismatch that is epistemic, not technical.
\end{foldable}

\begin{foldable}
  Data-free white-box distillation eliminates the requirement for real training data entirely, replacing it with synthetic image generation.
  The core strategy is \textit{model inversion}: training a generator to synthesize images that the teacher \DIFdelbegin \DIFdel{recognises }\DIFdelend \DIFaddbegin \DIFadd{recognizes }\DIFaddend as informative, using the teacher's internal structure as a guide.
  DAFL~\cite{chen2019data} \DIFdelbegin \DIFdel{maximises }\DIFdelend \DIFaddbegin \DIFadd{maximizes }\DIFaddend activation strength and enforces diversity by penalizing the student's output entropy.
  DFAD~\cite{fang2019data} introduces an adversarial distillation mechanism where a generator produces hard samples to enlarge the discrepancy between student and teacher, while both act as discriminators to reduce it.
  DeepInversion~\cite{yin2020dreaming} inverts the teacher by optimizing inputs to \DIFdelbegin \DIFdel{maximise }\DIFdelend \DIFaddbegin \DIFadd{maximize }\DIFaddend class-specific activation while matching the teacher's batch normalization statistics, anchoring synthesis to the teacher's training distribution.
  CMI~\cite{fang2021contrastive} applies contrastive learning to model inversion, distinguishing class-specific features and improving the diversity of synthesized images.
  NAYER~\cite{tran2024nayer} generates synthetic data by injecting noise into real examples, treating data generation as a layer in the student network for data-free distillation.
  GAN-based approaches further extend model inversion by training generative models to produce synthetic data.
  DeGAN~\cite{addepalli2020degan} uses a discriminator to improve image quality and diversity in data-free \ac{KD}.
  Momentum Adversarial Distillation~\cite{do2022momentum} \DIFdelbegin \DIFdel{stabilises }\DIFdelend \DIFaddbegin \DIFadd{stabilizes }\DIFaddend synthetic data generation through momentum updates and adversarial training, handling large distribution shifts in the data-free setting.
  All of these methods share one prerequisite: they require gradient flow through the teacher or read access to internal statistics (BatchNorm means and variances).
  Without white-box access, none of these techniques apply.
\end{foldable}

\begin{foldable}
  Removing white-box access is the dividing line between two qualitatively different problem settings.
  When teacher internals are inaccessible, the student loses every technical handle on the teacher's training distribution: no gradients to back-propagate synthesis refinement, no internal statistics to anchor generation, no layer activations to invert.
  The teacher becomes a sealed function that takes an image and returns a label, and the generator must reconstruct distributional coverage of the teacher's input space from the single integer the teacher outputs per query.
  This is not a harder version of model inversion; it is a different problem, requiring a different strategy.
  The design space collapses to response-based distillation with \DIFdelbegin \DIFdel{hard-label queriesonly}\DIFdelend \DIFaddbegin \DIFadd{output-only queries}\DIFaddend , and the information bottleneck becomes the binding constraint on what can be recovered.
\end{foldable}

\subsection{Black-Box Model Distillation: Few-Shot} \label{sec:24-background-kdrestricted-bbfewshot}
\begin{foldable}
  The loss of white-box access eliminates every technical handle on the teacher's distribution.
  The student's only recourse is to treat the teacher as a query interface over whatever training data is available, receiving only \DIFdelbegin \DIFdel{hard labels (single predicted class per image) rather than soft targets (probability distributions over all classes)}\DIFdelend \DIFaddbegin \DIFadd{its output class probabilities rather than logits or internal features}\DIFaddend .
  This setting is manageable when data is abundant: the student receives diverse coverage of the teacher's decision space, and the signal quality deficit is the only remaining constraint.
  But as data availability shrinks to few-shot scales, the problem changes fundamentally.
\end{foldable}

\begin{foldable}
  \textbf{Few-shot setting:} the student has access to only a small number of real training samples per class.
  Combined with black-box access\DIFdelbegin \DIFdel{that provides only hard labels}\DIFdelend \DIFaddbegin \DIFadd{, which exposes only the teacher's output probabilities and no logits or internal features}\DIFaddend , this creates a double bind: signal quality and data coverage are both constrained, and neither can compensate for the other.
  The student receives \DIFdelbegin \DIFdel{hard labels only (low }\DIFdelend \DIFaddbegin \DIFadd{only output probabilities (limited }\DIFaddend signal quality) on a narrow set of examples (low coverage), with no obvious mechanism to escape this bind.
  The few-shot setting is therefore qualitatively harder than black-box with full data access, and the strategy must change.
\end{foldable}

\begin{foldable}
  The degradation mechanism is straightforward.
  Standard \ac{KD} relies on querying the teacher on a large, diverse set of training examples to collect soft targets that span the input distribution.
  When $\hat{\mathcal{D}}$ is small, the set of queries is narrow: it covers only a few modes of the teacher's training distribution and misses the inter-class boundary regions where soft targets are most informative.
  The student effectively receives a compressed and biased view of the teacher's knowledge, rich for the modes present in $\hat{\mathcal{D}}$ and absent everywhere else.
  This coverage deficit leads to overfitting on the few-shot examples and poor generalization on the full test distribution.
\end{foldable}

\begin{foldable}
  \textbf{Standard augmentation} attempts to alleviate this by increasing the number of training examples through synthetic transformations.
  MixUp~\cite{zhang2018mixup} interpolates between pairs of training images and their labels, creating convex combinations that encourage linear \DIFdelbegin \DIFdel{behaviour }\DIFdelend \DIFaddbegin \DIFadd{behavior }\DIFaddend between training points.
  CutMix~\cite{yun2019cutmix} mixes regions of different images, producing examples that combine local patches from multiple classes.
  These methods increase the effective sample count and encourage smoother decision boundaries, but they operate within the convex hull of the few-shot training set: the synthesized images are still anchored to the observed examples and do not explore unseen input regions.
  Semantic diversity, the coverage of the teacher's full input distribution, is therefore not improved by standard augmentation.
\end{foldable}

\begin{foldable}
  \textbf{Generative augmentation} uses generative models to synthesize new images that extend beyond the few-shot support.
  Prior methods such as FSKD~\cite{kimura2018few} and WaGe~\cite{kong2020learning} required white-box access---FSKD uses trainable synthetic images as learnable parameters optimized via adversarial signals from the teacher's gradients, while WaGe explores the data space via Wasserstein loss on teacher logits---making them inapplicable in the black-box setting.
  BBKD~\cite{wang2020neural} constructs synthetic images using MixUp and trains a conditional GAN on $\hat{\mathcal{D}}$ with active learning to select the most informative examples for student training.
  FS-BBT~\cite{nguyen2022black} trains a conditional variational autoencoder (CVAE) as a secondary synthetic image source, conditioning both the GAN and CVAE generators on teacher feedback to refine images toward high-confidence regions.
  These methods produce visually plausible images that expand the training set beyond $\hat{\mathcal{D}}$.
  However, GAN generators trained on few-shot data inherit the distributional limitations of their training set: they can only model the modes present in their training data and largely ignore the long tail of the class distribution.
  The synthetic images therefore cluster around the same modes as $\hat{\mathcal{D}}$, inheriting rather than correcting the diversity deficit.
\end{foldable}

\begin{foldable}
  Synthetic data augmentation represents real progress, but it does not dissolve the core tension.
  Generative augmentation improves over MixUp and CutMix by escaping the strict convex hull of $\hat{\mathcal{D}}$, yet the generator's learned distribution remains tied to the same few-shot support through which all its training signal has passed.
  The improvement is one of degree, not kind: the synthesized images are more varied, but their diversity is still bounded by the diversity that $\hat{\mathcal{D}}$ could communicate to the generator during training.
  In distributional terms, the few-shot generator is caught in a loop: it can only synthesize from what it has seen, and what it has seen is already insufficient to cover the teacher's full input distribution.
  Standard and generative augmentation both operate inside the information horizon defined by the few-shot set; neither provides a way to step outside it.
  Addressing the coverage deficit therefore requires a signal that is not derived from $\hat{\mathcal{D}}$ at all, one that operates outside the few-shot support and can anchor synthesis to regions of the teacher's distribution that the few-shot examples do not represent.
\end{foldable}

\begin{foldable}
  \textbf{Open problem (Ch.~\ref{ch:30-research01}):}
  In black-box few-shot \ac{KD}, how can GAN-based synthesis be made to produce data that is \textit{both} high-fidelity \textit{and} semantically diverse, without access to teacher internals, when the generator itself is trained on a few-shot dataset that already lacks diversity?
\end{foldable}

\subsection{Black-Box Model Distillation: Data-Free} \label{sec:24-background-kdrestricted-bbdatafree}
\begin{foldable}
  In the data-free regime, no real training data is available, and the teacher returns only a top-1 hard label $T(\mathbf{x}) \in \{1, \ldots, K\}$ for each queried image.
  This matches the deployment of large \ac{AI} systems (commercial vision \acp{API}, cloud-hosted language models) in which the provider exposes only a label interface to protect model internals and limit adversarial probing.
\end{foldable}

\begin{foldable}
  ZSDB3~\cite{wang2021zero} assumes class robustness based on distance to the decision boundary and explores the teacher's decision boundary by generating query images near the hyperplanes that separate class regions, then using the resulting hard labels to train the student.
  The method initializes with random pixel-wise uniform noise and uses gradient estimation to refine it.
  Its limitation is scalability: decision boundary exploration in high-dimensional image spaces becomes increasingly sparse for complex datasets with many fine-grained classes, such that obtaining diverse image samples scales exponentially worse with high-resolution images or numerous classes, limiting success to small-scale, simple datasets.
\end{foldable}

\begin{foldable}
  IDEAL~\cite{zhang2022ideal} takes a different approach: a generator network is trained using the \textit{student's} gradients as a proxy feedback signal.
  Since the teacher is a black box, IDEAL substitutes the student's current predictions for the teacher's: the generator is rewarded for producing images on which the student and teacher disagree, using in-class similarity and class-balance objectives.
  This enables end-to-end training without white-box teacher access, but it introduces a dependency on the student's current state: the generator is pulled toward the student's current failure modes rather than the teacher's distribution.
  Relying on the signals of an immature student is not reliable enough to train the generator effectively.
  As the student improves, the generator feedback signal shifts, causing a form of sequential mode collapse: the generated images at any given training stage are diverse only with respect to the current student checkpoint, resulting in possible mode collapse and insufficient diversity.
\end{foldable}

\begin{foldable}
  DFHL-RS~\cite{yuan2024data} introduces a random search variant of the generator update, adding stochastic exploration to partially counteract the mode collapse observed in IDEAL.
  The method improves coverage on moderate-scale benchmarks but retains the hard-label bottleneck: since the teacher returns only top-1 labels, there is no soft distillation signal, and the student is trained on the least informative possible supervision.
\end{foldable}

\begin{foldable}
  Two structural deficits are shared by all existing \ac{BBDFKD} methods.
  First, the synthetic data lacks \textit{domain coverage}: without real images to anchor the generator, the synthesized distribution is not constrained to cover the teacher's training distribution.
  The generator may produce images that are plausible under some learned prior but do not reflect the diversity of the real-world inputs on which the teacher's knowledge is defined.
  Second, the hard-label interface is an \textit{information bottleneck}: a single integer per query provides $\log_2 K$ bits of supervision, compared to the $K$-dimensional probability vector that contains the inter-class relational structure that makes \ac{KD} effective.
\end{foldable}

\begin{foldable}
  To make the domain-coverage deficit concrete, consider a teacher trained on ImageNet, which comprises 1{,}000 classes and approximately 1.2 million training images spanning a highly multimodal natural-image manifold.
  A \ac{BBDFKD} generator with no access to real images must implicitly reconstruct this manifold from hard-label feedback alone, in effect inverting a 1{,}000-way classifier using only its top-1 decisions.
  The difficulty of this inversion is severe: natural images occupy a tiny, low-dimensional manifold within pixel space, and the fraction of pixel-space samples that fall anywhere near that manifold is vanishingly small.
  A generator \DIFdelbegin \DIFdel{initialised }\DIFdelend \DIFaddbegin \DIFadd{initialized }\DIFaddend from a generic prior and updated solely by top-1 label agreement will gravitate toward the modes that its prior already \DIFdelbegin \DIFdel{favours}\DIFdelend \DIFaddbegin \DIFadd{favors}\DIFaddend , because those are the regions where label feedback is most consistent.
  The long tail of each class, atypical viewpoints, unusual lighting, and boundary-region examples that lie close to adjacent classes, receives little or no gradient signal and remains uncovered.
  The practical consequence is a biased distillation corpus: the student is trained overwhelmingly on visually prototypical exemplars and is never exposed to the discriminative boundary structure on which the teacher's generalization depends.
  Distributional fidelity between the synthetic corpus and the teacher's training distribution is therefore violated at the level of class coverage, not just at the level of individual image quality.
\end{foldable}

\begin{foldable}
  The information-bottleneck deficit is equally concrete.
  A soft target vector encodes the full inter-class similarity structure: the probability mass assigned to visually related classes carries relational information that shapes how the student \DIFdelbegin \DIFdel{organises }\DIFdelend \DIFaddbegin \DIFadd{organizes }\DIFaddend its own representation space.
  A hard top-1 label discards all of this, recording only which class won while assigning zero to every other.
  The difference is qualitative, not merely quantitative: two images from the same coarse class but distinct fine-grained subclasses, a golden retriever and a German shepherd both labeled \texttt{dog}, receive identical supervision under hard-label access, while a soft target would assign them different relative probabilities over related breeds and thereby expose the within-class substructure.
  A student trained exclusively on hard labels cannot recover this within-class geometry; its representation must treat the two images as equivalent inputs carrying the same learning signal.
  This is the specific failure mode that motivates the \DIFaddbegin \DIFadd{surrogate }\DIFaddend soft-label \DIFdelbegin \DIFdel{recovery }\DIFdelend component developed in Chapter~\ref{ch:40-research02}: \DIFdelbegin \DIFdel{restoring }\DIFdelend \DIFaddbegin \DIFadd{supplying }\DIFaddend even an approximation of the soft probability vector partially repairs the distributional fidelity of the supervision signal, independent of whether the synthetic images themselves fully cover the real data manifold.
\end{foldable}

\begin{foldable}
  \textbf{Open problem (Ch.~\ref{ch:40-research02}):}
  In black-box data-free \ac{KD}, how can a student simultaneously (i) generate synthetic images that cover the teacher's full input distribution, and (ii) \DIFdelbegin \DIFdel{recover a }\DIFdelend \DIFaddbegin \DIFadd{obtain a surrogate }\DIFaddend soft distillation signal \DIFdelbegin \DIFdel{approximating }\DIFdelend \DIFaddbegin \DIFadd{in place of }\DIFaddend the teacher's class probabilities, both without any real training data and with only hard-label access to the teacher?
\end{foldable}
 \section{Distributional Fidelity and Synthetic Data} \label{sec:25-background-distribution}
\begin{foldable}
  All three contributions operate under restricted access: few-shot data (Chapter~\ref{ch:30-research01}), no data (Chapter~\ref{ch:40-research02}), or an oversized corpus (Chapter~\ref{ch:50-research03}).
  In each case the available training signal is insufficient to directly support high-fidelity distillation.
  The unifying thread is \textit{distributional fidelity}: whether the compact training signal, synthetic images or a selected text surrogate, adequately covers the original distribution.
  This section develops that thread as the strategic frame for Chapters~\ref{ch:30-research01}--\ref{ch:50-research03}, moving from the threat (\S\ref{sec:25-background-distribution-deficit}) to the toolkit (\S\ref{sec:25-background-distribution-mechanisms}) to diagnosis (\S\ref{sec:25-background-distribution-diagnosis}) to remedy (\S\ref{sec:25-background-distribution-remedy}).
\end{foldable}

\subsection{The Fidelity Deficit} \label{sec:25-background-distribution-deficit}
\begin{foldable}
  Restricted access degrades the training signal in a specific way: it reduces distributional coverage, not just signal quantity.
  A smaller or noisier signal misses portions of the original distribution; no downstream algorithm can recover knowledge absent from the training signal.
  This is not just a data scarcity problem; it is a structural problem in what information remains available after restriction.
  A few-shot dataset may contain high-quality examples, but if they cluster in a single mode of the teacher's distribution, no amount of augmentation or model capacity can teach the student about the other modes.
  Similarly, a data-free generator may produce thousands of plausible images, but if they concentrate on prototypical exemplars, they omit the boundary-region examples where the student most needs guidance.
  This mismatch between signal quantity and distributional coverage is the defining constraint of restricted-access distillation.
\end{foldable}

\begin{foldable}
  \textbf{Coverage as a component of distributional fidelity:}
  In model distillation (Chapters~\ref{ch:30-research01}--\ref{ch:40-research02}), the student's accuracy on held-out data is bounded by the fraction of the teacher's distribution covered by the synthetic training set.
  Even if the student has infinite capacity and each synthetic example carries high-quality supervision, it cannot learn structure it has never seen.
  If the teacher's decision boundary has structure in a region $\mathcal{R}$ that the synthetic set omits, the student will generalize randomly there.
  In dataset distillation (Chapter~\ref{ch:50-research03}), the surrogate's utility is bounded by how well it represents the full corpus distribution in the learning objective's space.
  A surrogate composed entirely of easy-to-classify examples may cover the input space densely yet fail to represent the corpus distribution over the learning trajectories that matter\DIFdelbegin \DIFdel{: and }\DIFdelend \DIFaddbegin \DIFadd{, which requires both }\DIFaddend representative examples that capture the distributional core \DIFdelbegin \DIFdel{, }\DIFdelend \DIFaddbegin \DIFadd{and }\DIFaddend boundary-adjacent examples that refine decisions.
  Coverage is therefore a necessary condition for distributional fidelity, not a separate objective\DIFdelbegin \DIFdel{; and for }\DIFdelend \DIFaddbegin \DIFadd{. For }\DIFaddend dataset distillation (Chapter~\ref{ch:50-research03}), fidelity to the full corpus distribution is the objective itself.
\end{foldable}

\begin{foldable}
  \textbf{Three instantiations of the same deficit:}
  Few-shot access leaves entire regions of the teacher's distribution unsampled, creating a convex-hull barrier: the student cannot generalize beyond the modes present in the few-shot set, no matter how well it fits them.
  Data-free access forces synthesis without any real-data anchor, causing the generator to gravitate toward the modes its implicit prior \DIFdelbegin \DIFdel{favours}\DIFdelend \DIFaddbegin \DIFadd{favors}\DIFaddend , abandoning regions where label feedback is sparse or contradictory.
  An oversized corpus, filtered naively by importance scores (gradient norms, loss values), concentrates on outliers and mislabeled examples at convergence, missing the representative core because noisy examples retain large gradients while clean examples' gradients vanish.
  All three are coverage failures under different constraints: first, insufficient data; second, no ground-truth reference; third, misaligned importance metrics.
\end{foldable}

\begin{foldable}
  This shared structure motivates a unified treatment: synthetic data (or a selected surrogate) must be evaluated and enforced on global fidelity, not just on fidelity of individual examples.
  The diagnostic distinction between Density (quality of individual examples) and Coverage (breadth of distributional support) must be maintained throughout design, optimization, and evaluation.
\end{foldable}

\subsection{Generation Mechanisms} \label{sec:25-background-distribution-mechanisms}
\begin{foldable}
  Synthetic data is generated differently depending on modality and access constraints, but both modalities share a common representational ground.
  Understanding the generation mechanism matters because the affordances and failure modes of each mechanism directly constrain which coverage-enforcement strategies are viable.
\end{foldable}

\begin{foldable}
  \textbf{For images (Chapters~\ref{ch:30-research01}--\ref{ch:40-research02}):}
  \Acp{GAN}~\cite{goodfellow2014generative} frame synthesis as a minimax game between generator $G: \mathcal{Z} \to \mathcal{X}$ and discriminator $D: \mathcal{X} \to [0,1]$.
  The generator samples from a latent distribution and produces candidate images; the discriminator learns to distinguish real from generated.
  Vanilla GANs suffer training instability when real and generated distributions are disjoint; the Jensen-Shannon divergence provides poor gradient signals and the generator collapses to a few high-density modes.
  This mode collapse is particularly severe in few-shot and data-free settings: when the generator is trained on few examples or receives only hard-label feedback, it has no signal to explore regions beyond the modes it has seen or can easily guess.
  \Ac{WGAN-GP}~\cite{gulrajani2017improved} addresses this by replacing Jensen-Shannon divergence with Wasserstein-1 distance, which provides more stable gradients even when distributions are disjoint.
  \Ac{WGAN-GP} \DIFdelbegin \DIFdel{stabilises }\DIFdelend \DIFaddbegin \DIFadd{stabilizes }\DIFaddend training and improves mode coverage by rewarding the discriminator for producing valid gradients everywhere in the input space, not just where the two distributions overlap.
  \Ac{WGAN-GP} is the generative backbone for \DIFdelbegin \DIFdel{Chapters~\ref{ch:30-research01}and~\ref{ch:40-research02}}\DIFdelend \DIFaddbegin \DIFadd{Chapter~\ref{ch:30-research01}}\DIFaddend , chosen for its stability and natural alignment with distributional fidelity objectives.
\end{foldable}

\begin{foldable}
  \textbf{For text (Chapter~\ref{ch:50-research03}):}
  \Acp{LLM} generate text autoregressively, sampling tokens sequentially from learned conditional distributions $p(t_i | t_1, \ldots, t_{i-1})$.
  Unlike image generation via continuous pixel optimization, text generation is inherently discrete: there is no gradient with respect to token indices, and straight-through estimators introduce severe approximation errors when applied to long sequences.
  Fine-tuning an \ac{LLM} on the task corpus and sampling from it produces fluent, human-readable synthetic examples, sidestepping the gradient-flow problem that makes pixel-space optimization inapplicable.
  The autoregressive sampling process ensures that generated text inherits the distributional properties of the training corpus: if the corpus is biased toward certain sentence structures or vocabulary, the \ac{LLM}-generated samples will reflect that bias.
  Fine-tuning on the task corpus steers generation toward task-relevant vocabulary and sentence structure, improving distributional alignment over generic \ac{LLM} sampling.
  The remaining challenge is distributional fidelity: even \DIFdelbegin \DIFdel{an }\DIFdelend \DIFaddbegin \DIFadd{a }\DIFaddend well-fine-tuned \ac{LLM} may oversample frequent patterns and neglect rare but informative examples, producing a fluent yet \DIFdelbegin \DIFdel{distributional-unfaithful }\DIFdelend \DIFaddbegin \DIFadd{distributionally unfaithful }\DIFaddend surrogate.
\end{foldable}

\begin{foldable}
  \textbf{Embedding space as common ground:}
  In both modalities, generated data is projected into a modality-specific embedding space via a pre-trained encoder, \eg, ResNet for images, BERT for text.
  This continuous representation decouples evaluation from the raw format of the generated data and provides a common geometric space in which distributional structure, mode collapse, and coverage gaps become measurable.
  The embedding space therefore serves as the shared interface between generation (this subsection), diagnosis (\S\ref{sec:25-background-distribution-diagnosis}), and active enforcement (\S\ref{sec:25-background-distribution-remedy}).
\end{foldable}

\subsection{Diagnosing Coverage Failures} \label{sec:25-background-distribution-diagnosis}
\begin{foldable}
  Measuring coverage requires separating \textbf{\DIFaddbegin \DIFadd{sample }\DIFaddend fidelity} (quality of individual examples) from \textbf{diversity} (breadth of distributional support); standard metrics conflate the two.
  \Ac{IS}~\cite{salimans2016improved} and \ac{FID}~\cite{heusel2017gans} are the dominant metrics for evaluating generative models.
  \ac{IS} measures fidelity via class-conditional sharpness: samples should have high predicted class confidence (sharp posterior) and collectively span many classes.
  \ac{FID} measures distributional distance by computing the 2-Wasserstein distance between real and generated embeddings, balancing fidelity (proximity of generated samples) and diversity (spread across feature space).
  Both metrics operate in embedding space and aggregate fidelity and diversity into a single number.
  This aggregation is the core problem: a generator can achieve low \ac{FID} by concentrating on high-density modes (high fidelity, narrow diversity) while leaving low-density regions empty.
  For example, in ImageNet, some classes are visually prototypical (\texttt{chihuahua} or \texttt{golden retriever}) while others have high variance (\texttt{dog}, which includes hundreds of distinct breeds and poses).
  A generator achieving low \ac{FID} might produce crisp, classifiable images of prototypical dogs while missing the long tail of atypical poses, unusual lighting, and boundary-region examples that appear in the real dataset.
  This coverage failure is invisible in the aggregate \ac{FID} score.
\end{foldable}

\begin{foldable}
  \textit{Density} and \textit{Coverage}~\cite{naeem2020reliable} decompose these explicitly.
  Density measures \DIFaddbegin \DIFadd{sample }\DIFaddend fidelity: the average distance from each generated sample to its nearest real neighbor in embedding space (lower is better; generated samples are close to the real distribution).
  Coverage measures recall: the fraction of the real distribution that has at least one generated sample nearby (higher is better; the real distribution is well-represented).
  A dataset can have high Density but low Coverage (a few crisp, real-looking samples missing most of the distribution) or high Coverage but low Density (samples spread across the distribution but perceptually distorted).
  Coverage is the primary diagnostic across all three settings.
  Qualitative diagnosis via t-SNE~\cite{van2008visualizing} or UMAP~\cite{mcinnes2018umap} in embedding space reveals whether synthetic samples cluster narrowly (high density, low coverage: mode collapse) or span the full distributional support (high coverage: good diversity).
  This structural information, which modes are covered, which are missed, where clusters occur, is essential for debugging and is invisible in aggregate scores.
\end{foldable}

\begin{foldable}
  Low \ac{FID} does not imply high Coverage.
  Conversely, high Coverage does not guarantee high Density: a broad sample set may under-represent the high-density modes of the real distribution.
  Coverage must be treated as an independent objective, not a consequence of good generation.
  This dissociation points to a core requirement: evaluation alone cannot prevent coverage failures.
  Measuring coverage after synthesis is too late; the damage is structural.
  Diversity must be actively enforced during training, generation, and selection, upstream of the final distillation, before the distribution of the resulting set is ever measured.
  A generator trained without explicit coverage incentives will naturally concentrate on high-reward modes (high confidence from the discriminator, or high agreement with the teacher), abandoning regions where feedback is ambiguous or absent.
\end{foldable}

\subsection{Active Fidelity Recovery} \label{sec:25-background-distribution-remedy}
\begin{foldable}
  Each instantiation of the fidelity deficit (\S\ref{sec:25-background-distribution-deficit}) requires an active mechanism matched to its constraints; diagnosis alone cannot correct a coverage failure.
  The challenge is that each setting imposes different structural constraints: few-shot has a data anchor but limited diversity; data-free has no anchor; text has discrete tokens and importance-weighted selection objectives.
  Three concrete mechanisms span the design space, each tailored to its setting's affordances and constraints.
\end{foldable}

\begin{foldable}
  \textbf{Few-shot model distillation (Chapter~\ref{ch:30-research01}).}
  The generator is trained on few-shot data and inherits the distribution of that data.
  Mode collapse occurs not because the generator is inherently limited, but because the gradient signal available during training covers only the few-shot modes.
  The generator cannot be improved by training longer: it learns the few-shot distribution perfectly and has nothing left to learn.
  Correcting this requires external signal from outside the few-shot set.
  Adaptive confidence filtering provides this signal: after generating candidates, the method filters via teacher querying, retaining only synthetic samples on which the teacher has high confidence (above a per-class adaptive threshold).
  High teacher confidence indicates that the sample is in a region the teacher \DIFdelbegin \DIFdel{knows well; rejection sampling on low-confidence samples pushes the generator toward unexplored regions where the }\DIFdelend \DIFaddbegin \DIFadd{recognizes well.
  Admitting such samples as additional real images gives the discriminator signal from outside the few-shot set, with confidence serving as a practical proxy for the }\DIFaddend teacher's \DIFdelbegin \DIFdel{decision boundary is uncertain}\DIFdelend \DIFaddbegin \DIFadd{input distribution}\DIFaddend .
  Per-class thresholds prevent imbalance: a class with few prototype modes might naturally have fewer high-confidence synthetics, skewing the final distribution.
  This pushes coverage beyond the convex hull of the few-shot set by using the teacher as a guide to which regions are under-explored.
\end{foldable}

\begin{foldable}
  \textbf{Data-free model distillation (Chapter~\ref{ch:40-research02}).}
  Without any real-data anchor, the generator gravitates toward the modes its implicit prior (\eg, uniform noise, natural image statistics) \DIFdelbegin \DIFdel{favours}\DIFdelend \DIFaddbegin \DIFadd{favors}\DIFaddend .
  These priors concentrate mass on simple, prototypical structures because they generalize well and are easy to learn from minimal training signal.
  Multi-scale noise synthesis addresses this by injecting structured noise at multiple frequency scales before each teacher query.
  Low-frequency noise introduces coarse-level diversity (different global scene structures); high-frequency noise adds fine-level variation (texture, details).
  This prevents the generator from collapsing to a single reconstruction.
  Complementing this, a contrastive objective operates on the generated candidates in embedding space, pushing them apart before any teacher interaction.
  The contrastive loss encourages samples to cluster in different regions of the embedding space, spreading the synthetic set to cover more modes.
  Only after this structural diversity is enforced are samples queried to the teacher for label feedback.
  This two-stage approach (diversity-first, then label-conditioning) corrects the failure mode where hard-label feedback alone drives convergence to high-confidence modes.
\end{foldable}

\begin{foldable}
  \textbf{Text dataset distillation (Chapter~\ref{ch:50-research03}).}
  The problem is two-fold: single-checkpoint gradient-norm importance scores are unreliable at convergence, and independent ranking-based selection cannot enforce distributional objectives.
  Trajectory-aware scoring integrates importance signals across multiple training checkpoints rather than a single convergence point.
  Early checkpoints produce large gradients for all examples; late checkpoints produce large gradients mainly for hard/noisy examples.
  Averaging across the trajectory dilutes the late-stage noise emphasis, producing calibrated scores that reflect each example's true contribution to learning.
  However, corrected scores alone are insufficient: naive ranking selects the top-$K$ examples by score, ignoring relationships and redundancy.
  Discrete optimal transport reformulates selection as a distributional problem: given importance weights on the full corpus, select a $K$-sized subset that \DIFdelbegin \DIFdel{minimises }\DIFdelend \DIFaddbegin \DIFadd{minimizes }\DIFaddend the Wasserstein distance between the weighted source distribution and the selected subset.
  This allows the solver to trade off individual importance scores for joint distributional coverage, deselecting redundant high-importance examples if that makes room for under-represented but collectively important regions.
\end{foldable}

\begin{foldable}
  The pattern is the same across all three: name the fidelity deficit $\to$ understand its structural cause $\to$ apply an upstream mechanism to correct it $\to$ verify with Coverage metrics.
  Distributional fidelity is an explicit design objective enforced during training and selection, not a property hoped for after the fact.
  Each mechanism operates before the final distillation loss is computed, ensuring that the training signal available to the student covers the target distribution.
\end{foldable}

\begin{foldable}
  The same fidelity concern applies when the training data itself is compressed rather than the model: a selected text surrogate must cover the original corpus distribution just as a synthetic image set must cover the teacher's distribution.
  Dataset Track (\S\ref{sec:26-background-datasetdistill}--\S\ref{sec:28-background-ot}) develops this parallel, introducing the dataset distillation landscape and the importance-scoring and optimal-transport tools that \DIFdelbegin \DIFdel{operationalise }\DIFdelend \DIFaddbegin \DIFadd{operationalize }\DIFaddend distributional fidelity in the discrete text domain.
  The mechanisms differ, selection vs.\ generation and discrete vs.\ continuous, but the principle is identical: distributional fidelity must be actively enforced, not assumed.
\end{foldable}
 \section{Dataset Distillation} \label{sec:26-background-datasetdistill}
\subsection{Problem Definition} \label{sec:26-background-datasetdistill-definition}
\begin{foldable}
  Dataset distillation~\cite{wang2018dataset} (also called dataset condensation) seeks a small surrogate dataset $\tilde{\mathcal{D}}$ with $|\tilde{\mathcal{D}}| = K \ll N = |\mathcal{D}|$ such that a model trained on $\tilde{\mathcal{D}}$ performs approximately as well as one trained on the full dataset $\mathcal{D}$:
  \begin{equation}
    \tilde{\mathcal{D}}^{*} = \arg\min_{|\tilde{\mathcal{D}}|=K}
    \mathbb{E}_{\theta \sim \mathcal{A}(\tilde{\mathcal{D}})}\!\left[\mathcal{L}(\theta; \mathcal{D})\right],
    \label{eq:dd-objective}
  \end{equation}
  where $\mathcal{A}(\tilde{\mathcal{D}})$ denotes the learning algorithm applied to $\tilde{\mathcal{D}}$.
  This objective is bold: the student model (trained on $\tilde{\mathcal{D}}$) is evaluated on the full held-out test set, not on the original dataset.
  A model that memorizes the surrogate is worthless; the surrogate must capture the generalization properties of the original dataset.
  The compression ratio $N/K$ is typically large, ranging from $10\times$ to $100\times$ in practice, and the surrogate may consist of either selected real examples (coreset selection) or synthesized examples that do not appear in $\mathcal{D}$.
  Figure~\ref{fig:26-background-datasetdistill-overview} illustrates the core objective: a model trained on the small distilled dataset should achieve similar performance to a model trained on the full dataset.
\end{foldable}

\begin{figure}[ht]
  \centering
  \includegraphics[width=0.75\columnwidth]{./figures/20-background/26-background-datasetdistill-datasetdistill.pdf}
  \caption[Dataset distillation.]{
    Dataset distillation objective.
    The original dataset $\mathcal{D}$ (left) is compressed into a small distilled dataset $\tilde{\mathcal{D}}$ (middle, green).
    A model trained on $\tilde{\mathcal{D}}$ (upper path) should achieve similar performance to a model trained on the full dataset $\mathcal{D}$ (lower path), measured by the distillation loss $\mathcal{L}_{\text{distill}}$ (right).
    The key challenge is preserving the generalization properties of the original dataset in a far smaller surrogate.
  }
  \label{fig:26-background-datasetdistill-overview}
\end{figure}

\begin{foldable}
  The objective~\eqref{eq:dd-objective} is intractable when the dataset is large: computing the full training trajectory for every candidate surrogate is prohibitively expensive.
  A tractable relaxation replaces loss-trajectory matching with distribution matching:
  \begin{equation}
    \tilde{\mathcal{D}}^{*} = \arg\min_{|\tilde{\mathcal{D}}|=K} d(P_{\tilde{\mathcal{D}}}, P_{\mathcal{D}}),
    \label{eq:dm-abstract}
  \end{equation}
  where $d$ measures distributional divergence in a task-relevant feature space (\eg, gradient space, embedding space, or optimal transport cost).
  Different choices of $d$ and feature space can lead to qualitatively different approaches to tackle the \ac{DD} problem.
  The assumption underlying the relaxation is that distributional proximity in the chosen feature space implies similar learning outcomes---an assumption that holds well empirically for gradient-matching but requires verification for each new distance metric.
\end{foldable}

\begin{foldable}
  Dataset distillation and model distillation are complementary but distinct in scope.
  Model distillation compresses a \textit{model}: given a trained teacher, it produces a smaller student with similar predictive \DIFdelbegin \DIFdel{behaviour}\DIFdelend \DIFaddbegin \DIFadd{behavior}\DIFaddend .
  Dataset distillation compresses a \textit{dataset}: given a full training corpus, it produces a smaller surrogate that induces similar model \DIFdelbegin \DIFdel{behaviour }\DIFdelend \DIFaddbegin \DIFadd{behavior }\DIFaddend when used for training.
  Model distillation evaluates success by student accuracy on held-out test data (how well the student predicts).
  Dataset distillation evaluates success by training a model on the surrogate and measuring its accuracy on held-out test data (how well a model trained on the surrogate generalizes).
  Both problems share the distributional fidelity requirement established in \S\ref{sec:25-background-distribution}, that the student model and the surrogate corpus must represent their source distributions faithfully, but they differ in their input (a trained model vs.\ a raw dataset) and output (a model vs.\ a corpus).
\end{foldable}

\begin{foldable}
  A complete dataset distillation method for text should satisfy three properties.
  First, it should be \textit{knowledge-informed}: the selection or synthesis of $\tilde{\mathcal{D}}$ should reflect each example's contribution to model learning, not just its frequency or surface features.
  Second, it should ensure \textit{distributional coverage}: the surrogate should represent the full range of semantic and task-relevant variation in $\mathcal{D}$, not just its high-frequency modes.
  Third, for text specifically, it should produce \textit{human-readable output}: unlike image distillation where synthetic pixels need not be interpretable, text examples must be fluent and coherent to be usable for fine-tuning and auditable for privacy inspection.
\end{foldable}

\subsection{Image Dataset Distillation} \label{sec:26-background-datasetdistill-imagedd}
\begin{foldable}
  Image dataset distillation methods synthesize a small set of images by optimizing them to preserve some property of training on the full dataset.
  Three optimization targets have been explored.
  \textit{Gradient matching}~\cite{zhao2021dataset} (DC) synthesizes $\tilde{\mathcal{D}}$ such that the model gradients computed on $\tilde{\mathcal{D}}$ match those computed on $\mathcal{D}$ at each training step:
  \begin{equation}
    \min_{\tilde{\mathcal{D}}} \sum_{t} \left\| \nabla_\theta \mathcal{L}(\theta_t;
    \tilde{\mathcal{D}}) - \nabla_\theta \mathcal{L}(\theta_t; \mathcal{D}) \right\|^2.
    \label{eq:dc-objective}
  \end{equation}
  Because the gradient matching loss is differentiable with respect to the pixel values of the synthesized images, the images can be optimized directly by backpropagation.
  The resulting synthetic images often appear unnatural, encoding gradient information rather than visual realism, but they reliably reproduce the training dynamics of the full dataset on simple benchmarks such as MNIST and CIFAR-10.
\end{foldable}

\begin{foldable}
  \textit{Trajectory matching}~\cite{cazenavette2022dataset,kim2022dataset} extends gradient matching from single-step gradients to entire parameter trajectories.
  Given expert parameters $\theta^*_t$ obtained by training on $\mathcal{D}$, the synthetic set is optimized so that $N$ steps of training on $\tilde{\mathcal{D}}$ starting from $\theta^*_t$ reach a point close to the expert checkpoint $M$ steps further along the trajectory:
  \begin{equation}
    \min_{\tilde{\mathcal{D}}} \sum_{t} \frac{\left\| \hat\theta_{t+N}(\tilde{\mathcal{D}}) - \theta^*_{t+M} \right\|^2}{\left\| \theta^*_t - \theta^*_{t+M} \right\|^2},
    \label{eq:mtt-objective}
  \end{equation}
  where $\hat\theta_{t+N}(\tilde{\mathcal{D}})$ is the parameter reached after $N$ gradient steps on $\tilde{\mathcal{D}}$ from $\theta^*_t$, and the denominator normalizes by the expert displacement so the loss is dimensionless across stages of training.
  Matching the entire trajectory rather than single-step gradients improves performance on larger and more complex datasets.
\end{foldable}

\begin{foldable}
  \textit{Distribution matching}~\cite{zhao2023dataset,wang2022cafe} takes a softer approach: instead of matching gradients or trajectories, it aligns the feature-space distributions of real and synthetic data.
  It instantiates~\eqref{eq:dm-abstract} with $d$ = squared L2 of mean embeddings under random feature maps $\psi_\vartheta$:
  \begin{equation}
    \min_{\tilde{\mathcal{D}}} \mathbb{E}_\vartheta \left\| \frac{1}{|\mathcal{D}|} \sum_{x \in \mathcal{D}} \psi_\vartheta(x) - \frac{1}{|\tilde{\mathcal{D}}|} \sum_{\tilde{x} \in \tilde{\mathcal{D}}} \psi_\vartheta(\tilde{x}) \right\|^2.
    \label{eq:dm-objective}
  \end{equation}
  Distribution matching is more computationally scalable than gradient or trajectory matching because it does not require unrolling the training computation graph.
\end{foldable}

\begin{foldable}
  All three methods share a basic incompatibility with the text domain.
  They synthesize examples by gradient-based optimization in the continuous pixel space of images.
  Text tokens are discrete: there is no well-defined gradient of a loss with respect to an integer token index.
  Straight-through estimators and continuous relaxations partially address this, but they introduce approximation errors that compound across the typically long token sequences of text examples, and the resulting `soft' token sequences are not human-readable.
\end{foldable}

\subsection{Text Dataset Distillation} \label{sec:26-background-datasetdistill-textdd}
\begin{foldable}
  The discrete and sequential nature of text requires a different approach to dataset distillation.
  Text operates over discrete tokens (integers indexing a vocabulary) with no natural interpolation between adjacent tokens, and semantic relationships span non-local dependencies that cannot be captured by per-token gradient information.
  These constraints invalidate the gradient matching, trajectory matching, and distribution matching approaches developed for continuous image data.
  Specifically: gradient-matching on pixels is differentiable ($\frac{\partial}{\partial x_{ij}}$ is well-defined for pixel values); gradient-matching on token indices is not ($\frac{\partial}{\partial t_i}$ where $t_i$ is a discrete token is undefined without approximation).
  Straight-through estimators and continuous relaxations partially address this by replacing discrete tokens with soft token distributions, but they introduce approximation errors that compound across token sequences (typically hundreds of tokens) and the resulting soft sequences are not human-readable or usable for downstream applications.
  Text imposes a constraint absent from image \ac{DD}: the synthetic outputs have to be \textit{semantically coherent}, not just gradient-aligned, so an objective that matches gradients while producing token-soup is not a valid solution.
  Prior methods for text dataset distillation fall into distinct families, each satisfying some but not all of the three properties outlined above.
\end{foldable}

\begin{foldable}
  \textbf{Soft-embedding methods} apply continuous optimization in the embedding space, instead of the discrete token space.
  SLDD~\cite{sucholutsky2021soft} optimizes a set of soft-label, continuous embedding vectors that, when used as training inputs to a language model, preserve model performance; it partially recovers interpretability by mapping each distilled embedding to its nearest bag-of-words in vocabulary space.
  DDTC~\cite{li2021data} similarly learns embedding sequences optimized by gradient matching against a differentiable fine-tuning objective.
  DDAL~\cite{maekawa2023dataset} takes a different angle within the same continuous-optimization regime: instead of learning embedding sequences directly, it optimizes the attention labels of a BERT model to match class-wise dataset statistics.
  By operating in embedding space, they sidestep the discrete-token gradient problem entirely.
  These methods achieve good compression ratios on simple classification tasks and satisfy property (3) in a minimal sense: the output is interpretable as importance scores on real examples.
  However, they satisfy neither property (2) nor the stronger interpretation of property (3).
  The optimized embeddings are \textit{architecture-locked}: they are tailored to a specific model's embedding layer and non-transferable to other models, making them brittle for deployment.
  None produce natural language text---the output is a set of continuous vectors in embedding space that are uninterpretable as sequences, cannot be shared or audited for privacy, and cannot be used outside the model training pipeline.
  For practical applications where auditing (\eg, privacy inspection, fairness analysis) or reuse (\eg, fine-tuning on different models) is required, this is a severe limitation.
\end{foldable}

\begin{foldable}
  \textbf{Clustering-based methods} select real examples by clustering the full dataset in embedding space and taking cluster representatives.
  DaLLME~\cite{tao2024textual} embeds all training examples using a language model and applies $k$-means clustering, selecting the example closest to each centroid.
  This approach is fast (linear in $N$), architecture-agnostic (operates on pre-computed embeddings), and produces human-readable output (real examples rather than soft embeddings), strongly satisfying property (3).
  However, $k$-means assigns uniform weight to each cluster---it selects exactly one representative per cluster, regardless of that cluster's contribution to model learning.
  A cluster containing critical, hard-to-learn boundary examples receives the same selection budget as a cluster of redundant, easy examples.
  This violates property (1), knowledge-informed weighting.
  Moreover, the selection criterion (Euclidean proximity to centroid) is geometric, not learning-objective-aware.
  The selection does not formally \DIFdelbegin \DIFdel{minimise }\DIFdelend \DIFaddbegin \DIFadd{minimize }\DIFaddend any measure of distributional divergence from the weighted source, violating property (2)---the method has no principled dataset-level objective, only independent per-cluster decisions.
  Clustering is a form of geometric summarization, effective for input-space coverage but not for learning-objective coverage.
\end{foldable}

\begin{foldable}
  \textbf{\ac{LLM}-based generation.}
  An alternative to selecting from $\mathcal{D}$ is to generate a synthetic candidate pool by fine-tuning a language model on $\mathcal{D}$ and sampling from it.
  DiLM~\cite{maekawa2024dilm} is an instance of this family: it generates readable candidate texts from a fine-tuned \ac{LLM} and then selects among them by gradient matching against the training corpus, followed by $k$\DIFdelbegin \DIFdel{-centre }\DIFdelend \DIFaddbegin \DIFadd{-center }\DIFaddend selection.
  This approach is attractive for several reasons.
  First, it avoids the gradient-flow problem entirely: the generation process is autoregressive and produces fluent, human-readable text automatically, satisfying property (3) by design.
  Second, generation provides privacy: rather than extracting examples from the original dataset, the surrogate consists of new, synthetic examples, minimizing the risk of direct data leakage.
  Third, the generated pool can be large (thousands of samples), providing a rich candidate set before selection.
  However, generation alone does not solve the selection problem; it shifts it.
  A large synthetic candidate pool must still be distilled to a compact surrogate of size $K$, and the selection step must address two residual problems shared by all prior text \ac{DD} methods.
  Property~(1), knowledge-informed weighting: naive importance scores (single-checkpoint gradient norms) are systematically biased toward noisy, hard-to-fit examples at convergence, not toward representative samples.
  Property~(2), distributional coverage: greedy or clustering-based selection operates independently on each example, ignoring whether the overall selected set covers the importance-weighted distribution.
  These two problems are orthogonal, and both must be solved to construct a complete method.
\end{foldable}
 \section{Influence Functions for Importance Weighting} \label{sec:27-background-influence}

\S\ref{sec:26-background-datasetdistill} identified property~(1) as a binding constraint: a complete text dataset distillation method needs knowledge-informed sample weighting that reflects each sample's contribution to the downstream objective.
Influence functions address this by quantifying each sample's counterfactual effect on learned parameters, providing a principled importance score that reflects true contribution rather than convergence-stage noise.
The text setting introduces a structural complication, the hard-sample bias, that this section develops before turning to its correction via trajectory integration.

\begin{foldable}
  \textbf{Influence as counterfactual contribution.}
  Influence functions~\cite{cook1982residuals,koh2017understanding} formalize a simple question: how much would removing a training example $z_n = (\mathbf{x}_n, y_n)$ from the dataset change the learned parameters?
  This counterfactual contribution is a direct measure of example importance for the learning objective, not a proxy based on geometry (which examples are central), frequency (which examples appear often), or loss (which examples are hard to fit).
  The influence of $z_n$ on the converged model $\hat\theta$ is approximated via implicit differentiation, treating the converged parameters as an implicit function of the training set:
  \begin{equation}
    \mathcal{I}(z_n) \approx
    -\nabla_\theta \mathcal{L}(z_n, \hat\theta)^\top
    H_{\hat\theta}^{-1}
    \nabla_\theta \mathcal{L}(z_n, \hat\theta),
    \label{eq:influence}
  \end{equation}
  where $\hat\theta$ is the converged model and $H_{\hat\theta} = \frac{1}{N}\sum_n \nabla^2_\theta \mathcal{L}(z_n, \hat\theta)$ is the Hessian of the training loss.
  Intuitively, the influence combines two factors: the gradient $\nabla_\theta \mathcal{L}(z_n, \hat\theta)$ (how far is $z_n$ from being fit) and $H_{\hat\theta}^{-1}$ (how sensitive are the converged parameters to that misfitting).
  A large positive influence indicates that $z_n$ is informative for learning, removing it would degrade the model significantly.
  Whereas a large negative influence indicates that $z_n$ is either mislabeled or an outlier and is harmful to generalization.
  Influence is defined relative to convergence, not relative to a training loss threshold; it captures counterfactual impact on the final learned \DIFdelbegin \DIFdel{behaviour}\DIFdelend \DIFaddbegin \DIFadd{behavior}\DIFaddend .
\end{foldable}

\begin{foldable}
  \textbf{Tractable approximations.}
  The practical limitation of exact influence functions is computational: Hessian inversion is $O(p^2)$ in space and $O(p^3)$ in time for dense matrices, intractable for modern scale.
  For a model the size of BERT-base ($p \approx 110$\,M parameters), storing a dense $p \times p$ Hessian is infeasible on any current hardware.
  Tractable approximations replace Hessian inversion with cheaper operations that capture the same intuition.
  TracIn~\cite{pruthi2020estimating} (Influence via TracIn) replaces $H^{-1}$ with a sum of gradient dot-products across training checkpoints, based on the implicit function theorem:
  \begin{equation}
    \text{TracIn}(z_n) = \sum_{t=1}^{T} \eta_t \,
    \nabla_\theta \mathcal{L}(z_n, \theta_t)^\top \,
    \nabla_\theta \mathcal{L}(z_n, \theta_t),
    \label{eq:tracin}
  \end{equation}
  where $\theta_t$ is the model checkpoint at step $t$ and $\eta_t$ is the learning rate at that step.
  This requires only storing gradients (a vector per checkpoint), not Hessians.
  TracIn is naturally trajectory-aware: it integrates influence contributions across the full training arc (all $T$ steps), not just the convergence point.
  Early-stage gradients capture which examples drive initial learning; late-stage gradients capture refinement and overfitting.
  Their sum provides a balanced importance estimate.
  TRAK~\cite{park2023trak} (Improving Tracin) scales further by projecting gradients onto low-rank subspaces, reducing both memory and computation to $O(p \cdot r)$ where $r \ll p$ is the rank.
  DataInf~\cite{kwon2023datainf} pushes the cost down further by replacing $H^{-1}$ with a closed-form approximation derived from a block-diagonal Hessian assumption, avoiding both Hessian inversion and trajectory storage.
  The simplest proxy, gradient norm $\|\nabla_\theta \mathcal{L}(z_n, \hat\theta)\|$ at convergence, trades theoretical accuracy for practical speed, requiring only a single forward-backward pass.
  This single-checkpoint approximation is widely used because it has near-zero computational overhead and often works well empirically when combined with other mechanisms.
\end{foldable}

\begin{foldable}
  \textbf{From test attribution to corpus-level scoring.}
  TracIn, TRAK, and DataInf were all designed as \textit{test-conditioned} attribution methods: they answer the question ``which training examples most influenced this specific test prediction?''.
  Dataset distillation needs the dual question: which training examples are most important to the corpus as a whole, independent of any one test point.
  The standard adaptation is \textit{self-influence}, $\mathcal{I}_\text{self}(z_n) = \mathcal{I}(z_n; z_n)$, obtained by setting the query point equal to the training point itself.
  Self-influence is then evaluated at the converged parameters $\hat\theta$, which is exactly where the next foldable's bias originates.
\end{foldable}

\begin{foldable}
  \textbf{Hard-sample bias: the structural failure.}
  Single-checkpoint influence estimates, whether computed exactly or approximately as gradient norms, suffer from a systematic bias at convergence.
  After sufficient training, a well-fitted model assigns near-zero loss to clean, easy examples; their gradients $\nabla_\theta \mathcal{L}(z_n, \hat\theta)$ vanish.
  Noisy, mislabeled, or atypical examples resist convergence: the model cannot simultaneously fit them and the clean majority to perfection, so their gradients remain large.
  A selector \DIFdelbegin \DIFdel{maximising }\DIFdelend \DIFaddbegin \DIFadd{maximizing }\DIFaddend single-checkpoint gradient norms $\|\nabla_\theta \mathcal{L}(z_n, \hat\theta)\|$ therefore preferentially selects the hardest, noisiest examples, exactly the examples most harmful to a student trained on the resulting surrogate.
  This creates a perverse inversion: examples that are most influential by counterfactual impact (hard-to-fit examples that the model has compromised on) are conflated with examples that are most representative of the true distribution (central, representative examples that the model fits perfectly).
  The selector picks outliers.
  This is not a tuning problem that can be fixed by rescaling, threshold adjustment, or different gradient aggregation; it is structural.
  The gradient signal itself at convergence is misaligned with the learning objective: a large gradient indicates the example is hard to fit, not that it is representative.
\end{foldable}

\begin{foldable}
  Empirical evidence for hard-sample bias appears in multiple literatures.
  \cite{arpit2017closer} show that neural networks fit clean patterns before noisy ones during training, so any score that averages over the late training phase is dominated by examples the network failed to fit.
  Dataset cartography~\cite{swayamdipta2020dataset} traces training dynamics epoch-by-epoch and defines a confidence-variability plane: each example receives a per-epoch confidence (model's probability on the true label) and a variability (standard deviation of confidence across epochs).
  Examples cluster into four regions: easy (high confidence, low variability), hard (low confidence, high variability), ambiguous (high variability), and memorized (high confidence, low variability).
  The hard region is enriched in mislabeled and noisy examples, examples that hurt generalization.
  Co-teaching~\cite{han2018co} studies noise-robust training through the lens of which examples large-loss samplers select: at late training stages, examples with large losses are predominantly noisy or mislabeled, not informative.
  Both findings confirm that gradient norms (which correlate with loss magnitude) at convergence are poor importance proxies.
\end{foldable}

\begin{foldable}
  \textbf{Trajectory integration as correction.}
  If single-checkpoint scores are biased by late-stage noise dominance, the solution is to integrate importance signals across multiple training stages rather than relying on convergence.
  Consider the training trajectory: at early epochs, all examples have large gradients because the model's initial random \DIFdelbegin \DIFdel{initialisation }\DIFdelend \DIFaddbegin \DIFadd{initialization }\DIFaddend is far from the true target.
  As training progresses, clean examples' gradients shrink (the model learns to fit them) while noisy examples' gradients plateau or grow (the model struggles to fit them without hurting clean examples).
  By late epochs, gradient magnitude is almost entirely determined by noise/hardness, not importance.
  Averaging gradient norms across $T$ checkpoints from different training stages $(t=1, 2, \ldots, T)$ dilutes this late-stage bias.
  The average example (clean or noisy) receives calibrated importance that reflects its true learning contribution across the full trajectory.
  Trajectory-integrated importance scores (as in TracIn or checkpoint averaging) reduce hard-sample bias substantially, producing importance estimates that correlate with representation quality rather than with noise.
  Theoretically, under distributional assumptions (\eg, noise is light-tailed), trajectory integration produces unbiased influence estimates; empirically, even simple checkpoint averaging provides substantial improvement.
\end{foldable}

\begin{foldable}
  \textbf{Corrected weights are necessary but insufficient.}
  Trajectory-integrated importance scores address property~(1): they quantify each sample's true contribution to learning.
  However, the selection rule applied to these corrected scores determines whether the chosen subset covers the importance-weighted distribution or concentrates on the highest-scoring subset of it.
  Property~(2) remains open: no matter how accurate the per-sample weights, ranking and selecting independently cannot guarantee that the chosen subset covers the weighted distribution.
\end{foldable}
 \section{Optimal Transport for Distributional Selection} \label{sec:28-background-ot}

\S\ref{sec:27-background-influence} corrected the per-sample weights but left property~(2) open: even with accurate weights, ranking and selecting independently cannot guarantee that the chosen subset covers the weighted distribution.
This section develops optimal transport as the tool that closes that gap, formalizing selection as a distributional problem rather than a ranking one.

\begin{foldable}
  Optimal transport has been widely applied across \ac{NLP} and machine learning, from Word Mover's Distance~\cite{kusner2015word} and Wasserstein autoencoders~\cite{tolstikhin2018wasserstein} to domain adaptation~\cite{courty2016optimal} and Wasserstein \acp{GAN}~\cite{arjovsky2017wasserstein}, before we turn to its role for distributional coverage in dataset distillation.
\end{foldable}

\begin{foldable}
  \textbf{The distributional coverage problem.}
  Property~(2) requires that the selected subset jointly represent the importance-weighted distribution, not just consist of high-importance examples.
  At extreme compression, when $K \ll N$, a natural objective emerges: select a compact subset that \DIFdelbegin \DIFdel{minimises }\DIFdelend \DIFaddbegin \DIFadd{minimizes }\DIFaddend the distance between the reweighted source distribution and the selected subset.
  This is a \textit{distributional} problem: it cannot be solved by ranking examples independently.
\end{foldable}

\begin{foldable}
  \textbf{Why independent selection fails architecturally.}
  Any rank-and-select strategy decomposes as a sum of independent per-sample decisions, invariant to relationships between selected examples.
  The selection of one example does not constrain or inform the selection of another; the algorithm ignores pairwise distances and redundancy in feature space.
  By contrast, a distributional objective penalizes redundancy and rewards coverage: if two selected examples are similar and both have high importance, a distributional selector can free capacity by deselecting one and using that slot for a distinct uncovered region.
  Architecturally, any greedy or clustering-based method that optimizes per-sample scores independently cannot solve a joint distributional problem.
\end{foldable}

\begin{foldable}
  \textbf{Alternative divergences and their limits.}
  Before naming a specific divergence, two natural candidates from the wider literature are worth ruling out for the regime here.
  Maximum mean discrepancy~\cite{gretton2012kernel} measures distance between distributions by comparing their mean embeddings under a kernel, but the comparison is symmetric and aggregate.
  It does not produce a sample-level assignment from source to surrogate, so selection on top of it requires an auxiliary heuristic that re-introduces the independent-selection problem of the previous foldable.
  Kullback--Leibler and Jensen--Shannon divergences require estimating the densities of source and surrogate, which is unreliable in the low-data regime~\cite{paninski2003estimation}; on top of that they become infinite when the supports are disjoint, which is the rule rather than the exception when $K \ll N$.
\end{foldable}

\begin{foldable}
  \textbf{Wasserstein distance and optimal transport.}
  Optimal transport formalizes the distributional coverage problem geometrically as a transportation problem.
  Imagine mass is distributed over the training set according to importance weights $w_n$; we want to transport this mass to a sparse surrogate of size $K$.
  The cost of moving one unit of mass from source example $z_n$ to surrogate example $\tilde{z}_c$ is the feature-space distance $d(z_n, \tilde{z}_c)$ (\eg, Euclidean distance in embedding space).
  The optimal transport problem finds a coupling (transport plan) $\gamma$ that moves all importance mass from source to surrogate at minimum total cost.
  The Wasserstein-$p$ distance~\cite{villani2008optimal} quantifies this minimum cost:
  \begin{equation}
    \label{eq:wasserstein}
    W_p(\mu, \nu) = \left(\inf_{\gamma \in \Gamma(\mu,\nu)} \int d(z, z')^p \, \text{d}\gamma(z, z') \right)^{1/p},
  \end{equation}
  where $\Gamma(\mu, \nu)$ is the set of all couplings (joint distributions) with marginals $\mu$ (source) and $\nu$ (surrogate).
  Intuitively, a small Wasserstein distance means the surrogate can cover the importance-weighted source cheaply: few examples have to travel far.
  Geometrically, $W_p(\mu, \nu)$ is the minimum cost of transporting probability mass from $\mu$ to $\nu$.
  Where \ac{KL} and Jensen-Shannon blow up under disjoint supports, Wasserstein keeps producing a meaningful distance, which matters once $K \ll N$.
  This makes Wasserstein the natural choice for measuring distributional coverage in dataset distillation.
\end{foldable}

\begin{foldable}
  \textbf{Discrete optimal transport formulation.}
  In dataset distillation, both source and target are discrete, finite sets and can be formulated as an \ac{OT} problem.
  Let the source be $\{(z_n, w_n)\}_{n=1}^{N}$ where $w_n$ are normalized importance weights (each $w_n \geq 0$, $\sum_n w_n = 1$).
  The candidate pool (synthetic examples or a large real example pool) is $\{(\tilde{z}_c, 1/C)\}_{c=1}^{C}$ with uniform weights (each candidate has equal potential budget $1/C$).
  The discrete optimal transport problem is a distributional selection problem: given weighted source examples and an unweighted candidate pool, find a transport plan that moves all source mass to the candidate pool at minimum cost.
  Formally, seek a transport plan $\gamma^* \in \mathbb{R}^{N \times C}_{\geq 0}$ \DIFdelbegin \DIFdel{minimising }\DIFdelend \DIFaddbegin \DIFadd{minimizing }\DIFaddend total transport cost subject to marginal constraints:
  \begin{equation}
    \gamma^* = \arg\min_{\gamma \geq 0} \langle C_{\text{dist}}, \gamma \rangle
    \quad \text{s.t.} \quad \gamma \mathbf{1}_C = \mathbf{w}, \quad
    \gamma^\top \mathbf{1}_N = \tfrac{1}{C}\mathbf{1}_C,
    \label{eq:discrete-ot}
  \end{equation}
  where $C_{\text{dist},nc} = d(z_n, \tilde{z}_c)^2$ is the squared cost of transporting mass from source example $z_n$ to candidate example $\tilde{z}_c$ (typically Euclidean distance in embedding space).
  The constraints enforce the transportation semantics: all importance mass leaves the source ($\gamma \mathbf{1}_C = \mathbf{w}$) and equal weight arrives at each candidate ($\gamma^\top \mathbf{1}_N = \tfrac{1}{C}$).
  The entry $\gamma^*_{nc}$ is a fractional transport weight encoding how much importance mass from source $z_n$ is optimally routed to candidate $\tilde{z}_c$.
  After solving, the $K$ candidates with the largest total assigned mass $m_c = \sum_n \gamma^*_{nc}$ are selected as the surrogate.
  Finally, \ac{OT} selects jointly to \DIFdelbegin \DIFdel{minimise }\DIFdelend \DIFaddbegin \DIFadd{minimize }\DIFaddend distributional divergence, not independently to \DIFdelbegin \DIFdel{maximise }\DIFdelend \DIFaddbegin \DIFadd{maximize }\DIFaddend per-example scores.
\end{foldable}

\begin{foldable}
  \textbf{Sinkhorn algorithm for tractable computation.}
  Solving discrete \ac{OT} exactly via linear programming is expensive for large pools (quadratic or worse in $N + C$).
  The Sinkhorn algorithm~\cite{cuturi2013sinkhorn} reformulates the problem with entropic regularization: adding a penalty term $\varepsilon \, \text{KL}(\gamma \| \mathbf{w} \otimes \tfrac{1}{C}\mathbf{1}_C)$ to the objective.
  This transforms the problem from a hard-constrained LP into a strongly convex one with a simple closed-form update: alternating row and column normalization.
  The update is simple: normalize rows of $\gamma$ to satisfy the source marginal, then normalize columns to satisfy the target marginal, repeat until convergence.
  Each iteration costs $O(N \cdot C)$, and convergence occurs in tens of iterations, making the algorithm practical for large pools.
  The \DIFdelbegin \DIFdel{regularised }\DIFdelend \DIFaddbegin \DIFadd{regularized }\DIFaddend plan $\gamma_\varepsilon^*$ converges to the exact \DIFdelbegin \DIFdel{unregularised }\DIFdelend \DIFaddbegin \DIFadd{unregularized }\DIFaddend \ac{OT} optimum $\gamma^*$ as $\varepsilon \to 0$, providing a principled approximation-accuracy tradeoff.
  The regularization parameter $\varepsilon$ is a tuning knob: large $\varepsilon$ spreads mass diffusely (smooth, entropy-maximizing plan), small $\varepsilon$ produces concentrated plans closer to exact \ac{OT}.
\end{foldable}

\begin{foldable}
  \textbf{Concrete example: \ac{OT} vs. greedy and clustering.}
  Suppose a training set has six examples in two semantic regions: Region~A has four examples each with importance $0.25$ (total $1.0$); Region~B has two examples each with importance $0.05$ (total $0.10$).
  From a candidate pool of four examples (two near each region), select $K=2$ for the surrogate.
\end{foldable}

\begin{foldable}
  \textbf{$k$-means clustering} partitions candidates by geometric proximity without consulting importance weights.
  If Region~B candidates are geometrically cohesive, they form one cluster; Region~A candidates form the other.
  $k$-means selects one centroid per cluster: one Region~A candidate and one Region~B candidate.
  The surrogate allocates $50\%$ capacity to Region~B despite Region~B carrying only $9\%$ of the importance mass, a factor-of-five mismatch.
\end{foldable}

\begin{foldable}
  \textbf{Greedy $k$\DIFdelbegin \DIFdel{-centre }\DIFdelend \DIFaddbegin \DIFadd{-center }\DIFaddend selection} \DIFdelbegin \DIFdel{minimises }\DIFdelend \DIFaddbegin \DIFadd{minimizes }\DIFaddend worst-case distance: it chooses the two geometric extremes, again one Region~A and one Region~B candidate.
  Same outcome: both methods conflate geometric coverage with distributional fidelity, and when the two decouple, as in text where low-importance examples cluster in distinctive surface-form niches, geometry actively misleads.
\end{foldable}

\begin{foldable}
  \textbf{Optimal transport} with weights $\mathbf{w} = (0.25, 0.25, 0.25, 0.25, 0.05, 0.05)$ constructs a plan that moves Region~A's mass (1.0 units) and Region~B's mass (0.10 units) to the candidate pool.
  The Sinkhorn plan assigns $\approx 0.50$ units to each Region~A candidate and $\approx 0.05$ units to each Region~B candidate.
  Selecting the two candidates with largest total assigned mass returns both Region~A candidates and neither Region~B candidate.
  The surrogate omits Region~B entirely, exactly the outcome that \DIFdelbegin \DIFdel{minimises }\DIFdelend \DIFaddbegin \DIFadd{minimizes }\DIFaddend Wasserstein distance from the importance-reweighted source.
  The advantage of \ac{OT} is not geometric coverage but \textit{importance-weighted} coverage: it concentrates surrogate capacity on the regions that matter most to the learning objective.
\end{foldable}
 \section{Summary and Research Gaps} \label{sec:29-background-summary}
\begin{foldable}
  This chapter has built the technical foundation for the thesis along two parallel tracks.
\end{foldable}

\begin{foldable}
  \textbf{Model Track (Sections~\ref{sec:23-background-modeldistill}--\ref{sec:24-background-kdrestricted}).}
  Image knowledge distillation under restricted access: the standard \ac{KD} framework and soft-target signal; the increasingly restricted settings, few-shot and data-free, that arise in practice; and the diagnosis that distributional coverage is the binding constraint.
  Section~\ref{sec:25-background-distribution} established generative and contrastive mechanisms for enforcing coverage when standard assumptions fail.
\end{foldable}

\begin{foldable}
  \textbf{Dataset Track (Sections~\ref{sec:26-background-datasetdistill}--\ref{sec:28-background-ot}).}
  Text dataset distillation: the language \DIFdelbegin \DIFdel{modelling }\DIFdelend \DIFaddbegin \DIFadd{modeling }\DIFaddend foundations and \acl{NLP} task structure; the dataset distillation objective and why image synthesis methods do not transfer to discrete text; prior text \ac{DD} methods and their limitations; and the diagnosis of distributional coverage as the binding constraint.
  Sections~\ref{sec:27-background-influence} and~\ref{sec:28-background-ot} then established the two tools the dataset track relies on: influence-based importance weighting and discrete optimal transport for distributional selection when independent ranking fails.
\end{foldable}

\begin{foldable}
  \textbf{The unifying thread: distributional fidelity.}
  Whether distillation operates on the model (Model Track: Chapters~\ref{ch:30-research01}--\ref{ch:40-research02}) or the dataset (Dataset Track: Chapter~\ref{ch:50-research03}), the binding constraint is the same: the compact \DIFdelbegin \DIFdel{artefact }\DIFdelend \DIFaddbegin \DIFadd{artifact }\DIFaddend must cover the distribution of its source.
  Coverage failures manifest differently, mode collapse in GAN-based synthesis and coverage gaps in clustering-based corpus selection, but stem from the same root cause: the optimization objectives used to produce the \DIFdelbegin \DIFdel{artefact }\DIFdelend \DIFaddbegin \DIFadd{artifact }\DIFaddend do not directly \DIFdelbegin \DIFdel{penalize distributional divergence}\DIFdelend \DIFaddbegin \DIFadd{promote distributional coverage}\DIFaddend .
  The three contributions in Chapters~\ref{ch:30-research01}--\ref{ch:50-research03} \DIFdelbegin \DIFdel{each make direct divergence minimization tractable in its respective context}\DIFdelend \DIFaddbegin \DIFadd{address this in their respective contexts: Chapters~\ref{ch:30-research01} and~\ref{ch:40-research02} actively promote distributional coverage, and Chapter~\ref{ch:50-research03} minimizes a transport cost to the source distribution}\DIFaddend .
\end{foldable}

\begin{table}[ht]
  \centering
  \begin{threeparttable}
    \caption{Research gaps addressed by the three thesis contributions.}
    \label{tab:research-gaps}
    \begin{tabular}{l l l l}
      \toprule
      \textbf{Research}  & \textbf{Setting} & \textbf{Motivation} & \textbf{Contributions}  \\
      \midrule
      DivBFKD & \makecell[l]{Black-box teacher \\ Limited real data} & \makecell[l]{Enrich diversity} & \makecell[l]{High-confidence images \\ Adaptive adversarial training} \\
      \midrule
      DIPKD & \makecell[l]{Black-box teacher \\ No real data} & \makecell[l]{Domain coverage} & \makecell[l]{Diverse image priors \\ Contrastive optimization \\ Primer student} \\
      \midrule
      TAKE & \makecell[l]{Text dataset \\ Extreme compression} & \makecell[l]{Preserve fidelity} & \makecell[l]{Influence-based scoring \\ \ac{LLM}-based generation \\ Discrete OT distillation} \\
      \bottomrule
    \end{tabular}
  \end{threeparttable}
\end{table}

\begin{foldable}
  Three open problems follow directly from this analysis.
\end{foldable}

\begin{foldable}
  \textbf{Gap 1 (Chapter~\ref{ch:30-research01}): Few-shot black-box knowledge distillation.}
  In black-box few-shot \ac{KD}, the student trains on a small \DIFaddbegin \DIFadd{labeled }\DIFaddend image set and receives only \DIFdelbegin \DIFdel{top-1 predictions }\DIFdelend \DIFaddbegin \DIFadd{class probabilities }\DIFaddend from a black-box teacher.
  Existing methods expand coverage via synthetic images but lack explicit diversity mechanisms.
  BBKD uses MixUp (linear interpolation within the few-shot support), adding little novel signal.
  FS-BBT employs a conditional VAE for out-of-distribution synthesis, but optimizes pixel reconstruction, not diversity, causing blurring and mode collapse.
  Both methods apply standard GAN training without steering toward unexplored regions.
  The binding constraint is diversity: synthetic samples must venture beyond the few-shot set's inherent modes to provide \DIFdelbegin \DIFdel{decision-boundary coverage }\DIFdelend \DIFaddbegin \DIFadd{broader coverage of the teacher's input space}\DIFaddend .
  Chapter~\ref{ch:30-research01} (DivBFKD) introduces a teacher-grounded feedback loop: adaptively filtering high-confidence synthetic images (with class-specific, dataset-agnostic thresholds), then feeding them back into WGAN training on-the-fly, allowing the teacher itself to guide which regions merit generation.
\end{foldable}

\begin{foldable}
  \textbf{Gap 2 (Chapter~\ref{ch:40-research02}): Black-box data-free knowledge distillation.}
  In black-box data-free \ac{KD}, no real training data exists and the teacher returns only top-1 labels, the narrowest information interface.
  Existing methods use synthetic data but face two independent bottlenecks.
  First, domain coverage: synthetic images fail to replicate the hierarchical structures and semantic variation of natural environments, resulting in distribution mismatch with the teacher's expertise.
  Second, information loss: hard top-1 labels withhold the inter-class relationships (dark knowledge) that distinguish \ac{KD} from standard supervised learning.
  Prior methods (ZSDB3, IDEAL, DFHL-RS) address one or the other but not both.
  ZSDB3 explores decision boundaries but ignores information; IDEAL uses student feedback but inherits mode collapse; DFHL-RS adds stochasticity but does not recover soft targets.
  Chapter~\ref{ch:40-research02} (\ac{DIPKD}) tackles both simultaneously via a three-phase pipeline: (1) Synthesis of diverse image priors using multi-scale noise; (2) Contrast via a primer student (white-box mediator) that enforces distinction via contrastive learning; (3) Distillation where the primer student \DIFdelbegin \DIFdel{recovers soft probabilistic signals, enabling dark-knowledge transfer despite }\DIFdelend \DIFaddbegin \DIFadd{supplies surrogate soft supervision, standing in for the dark knowledge that }\DIFaddend hard-label teacher access \DIFaddbegin \DIFadd{withholds}\DIFaddend .
\end{foldable}

\begin{foldable}
  \textbf{Gap 3 (Chapter~\ref{ch:50-research03}): Text dataset distillation.}
  Text dataset distillation requires three properties: (1) knowledge-informed weighting, (2) distributional coverage over the importance-weighted corpus, (3) human-readable output.
  No prior method achieves all three: soft-embedding methods are architecture-locked and unreadable; clustering methods assign uniform weight per cluster, ignoring importance; importance-based methods optimize per-sample scores without joint coverage.
  The binding constraint is dual: importance estimation suffers from what Chapter~\ref{ch:50-research03} names and formalizes as the \textit{hard-sample bias}, where single-checkpoint gradient norms favor noisy examples over representative ones.
  Even with corrected scores, independent ranking cannot guarantee that the chosen set jointly covers the full importance-weighted distribution.
  These failures are orthogonal and both must be solved.
  Chapter~\ref{ch:50-research03} (TAKE) is the first method to address all three: (1) correcting hard-sample bias via trajectory-integrated influence scores; (2) leveraging \ac{LLM} fine-tuning to generate human-readable candidates; (3) applying discrete optimal transport to distill candidates into a compact surrogate \DIFdelbegin \DIFdel{minimising }\DIFdelend \DIFaddbegin \DIFadd{minimizing }\DIFaddend importance-weighted Wasserstein distance to the full corpus.
\end{foldable}

\begin{foldable}
  Chapters~\ref{ch:30-research01}, \ref{ch:40-research02}, and~\ref{ch:50-research03} address Gaps~1, 2, and~3 respectively, each \DIFdelbegin \DIFdel{operationalising }\DIFdelend \DIFaddbegin \DIFadd{operationalizing }\DIFaddend distributional fidelity under its specific constraint.
  Despite their different modalities (images vs.\ text) and mechanisms (generation vs.\ distillation, soft targets vs.\ hard labels, continuous vs.\ discrete), all three contributions share a unifying design principle: distributional coverage must be actively enforced during training, not left as a byproduct of optimizing traditional objectives (generator loss, importance scores, gradient matching).
  This principle, introduced in \S\ref{sec:25-background-distribution}, is the strategic through-line that unifies this thesis.
\end{foldable}
 \clearpage{}
\clearpage{}\acresetall

\chapter[Diversity in Black-box Few-shot Knowledge Distillation]{Diversity in Black-box Few-shot \\ Knowledge Distillation} \label{ch:30-research01}

\begin{foldable}
  The background section on model distillation under restricted access (\S\ref{sec:24-background-kdrestricted}) identifies a fundamental bottleneck: standard augmentation and generative augmentation both operate inside the information horizon defined by the few-shot training set.
  Addressing the coverage deficit requires a signal that comes from outside this few-shot set---one that can anchor synthesis to regions of the teacher's distribution that the few-shot examples do not represent.
  This chapter answers the open problem by treating the teacher as precisely this external signal: high-confidence synthetic images, selected under the teacher's supervision, inject information from outside the few-shot support into the generator's training loop.
\end{foldable}

\section{Introduction} \label{sec:31-research01-introduction}
\begin{foldable}
  The fruitful performance of \ac{KD} typically relies on the availability of an extensive training set to train the student.
  Traditional \ac{KD} methods often assume that the student is distilled with the same data used for training the teacher~\cite{hinton2015distilling,kim2018paraphrasing,ahn2019variational,tian2020contrastive}.
  However, in real-world scenarios, the distillation often happens at an external party side, where one can only afford a limited-size training set.
  For example, to distill Google's FaceNet (trained using approximately 200 million non-public face images~\cite{schroff2015facenet}) into an in-house model without access to a large-scale dataset like Google's, an external party may only be able to gather a few thousand images or even less.
  This scenario poses the \textit{few-shot} setting and makes it more challenging than standard \ac{KD}.
\end{foldable}

\begin{figure}[ht]
  \centering
  \includegraphics[width=0.99\columnwidth]{./figures/30-research01/31-research01-introduction-standardkd-vs-bbfskd.pdf}
  \caption[Common KD methods and black-box few-shot KD.]{
    Common KD methods and black-box few-shot KD.
    Common KD methods (a) assume no constraints whereas black-box few-shot KD, and (b) only has access to a few images and the teacher's predictive probabilities.
  }
  \label{fig:01-introduction-standard_kd_vs_bbfskd}
\end{figure}

\begin{foldable}
  Another assumption made by traditional \ac{KD} methods is the requirement of a \textit{white-box} teacher \ie, we have complete access to its parameters, activations, or gradients, \etc~\cite{romero2015fitnets,yim2017gift,kimura2018few,chen2019data}.
  However, the teacher model is often \textit{black-box}, \ie, only its final predictions are accessible.
  Understandably, companies often avoid full disclosure of the model parameters as it is their competitive edge and intellectual property.
  For instance, even with a paid subscription to ChatGPT~\cite{openai2022introducing}, only text outputs and their prediction confidence can be retrieved.
  The realistic constraints in the \textit{black-box few-shot} \ac{KD} setting (Fig.~\ref{fig:01-introduction-standard_kd_vs_bbfskd}) render most \ac{KD} methods inapplicable and demand novel approaches.
\end{foldable}

\begin{foldable}
  Recently, some methods have addressed the few-shot \ac{KD} problem, but most of them still require a white-box teacher~\cite{kimura2018few,kong2020learning}.
  To the best of our knowledge, only two methods apply a black-box teacher to few-shot \ac{KD}, namely BBKD~\cite{wang2020neural} and FS-BBT~\cite{nguyen2022black}.
  Both these methods use MixUp~\cite{zhang2018mixup} to generate synthetic images.
  However, as MixUp linearly interpolates between two original images, the synthetic images are either very similar to the original images or do not have realistic semantics, adding little value to the student's training.
  To overcome this problem, FS-BBT employs a conditional variational autoencoder (CVAE)~\cite{sohn2015learning} to generate out-of-distribution synthetic images.
  However, while the pixel-based loss in CVAE is beneficial for reconstructing images, it is not intended for promoting image diversity and have side-effect blurring artifacts~\cite{bond2021deep}.
\end{foldable}

\begin{foldable}
  We address the problem of black-box few-shot \ac{KD} \ie, distilling the student with only a few images and a black-box teacher.
  To boost student performance, we expand the distillation set with diverse synthetic images.
  To achieve diversity in the synthetic images, we propose a novel training scheme for Wasserstein Generative Adversarial Network (WGAN)~\cite{arjovsky2017wasserstein}.
  Firstly, we specifically define \textit{high-confidence} images, which are synthetic images that the teacher can predict confidently w.r.t.\ a threshold, and we use the teacher to compute \textit{adaptive thresholds} in a class-specific and dataset-agnostic manner to avoid biased and abnormal images.
  Then, during the WGAN training loop, we select high-confidence images from the latest iteration on-the-fly and introduce them in the role of additional real images, allowing our WGAN to learn and generate diverse synthetic images \DIFdelbegin \DIFdel{close to the teacher 's training images}\DIFdelend \DIFaddbegin \DIFadd{that resemble what the teacher recognizes confidently}\DIFaddend .
  Finally, we augment our few-shot images with synthetic images to improve the training of the student, leading to a significantly better performance.
\end{foldable}

\begin{foldable}
  We summarize our contributions as follows:
  \begin{enumerate}
    \item \textbf{DivBFKD}: a method for black-box few-shot \ac{KD} that directly targets diversity of synthetic images as the binding constraint.
    \item A WGAN training scheme that uses the teacher as an external signal, introducing high-confidence synthetic images on-the-fly to anchor generation outside the few-shot information horizon.
    \item Adaptive per-class confidence thresholds that calibrate selection to each class's \DIFdelbegin \DIFdel{decision boundary}\DIFdelend \DIFaddbegin \DIFadd{confidence distribution}\DIFaddend , preventing bias toward easy \DIFdelbegin \DIFdel{cases}\DIFdelend \DIFaddbegin \DIFadd{classes}\DIFaddend .
    \item Comprehensive evaluation across seven datasets, with ablation confirming that class-wise calibration is load-bearing.
  \end{enumerate}
\end{foldable}

\begin{foldable}
  The chapter is organized as follows.
  \S\ref{sec:32-research01-framework} details the WGAN training scheme and the two mechanisms that anchor synthesis outside the few-shot support: high-confidence image selection and adaptive per-class thresholds.
  \S\ref{sec:33-research01-experiments} evaluates DivBFKD across seven image classification benchmarks and includes ablation studies confirming that class-wise calibration is load-bearing.
  \S\ref{sec:34-research01-conclusion} discusses limitations and connects this contribution to the constraints that motivate Chapter~\ref{ch:40-research02}.
\end{foldable}
 \section{Framework} \label{sec:32-research01-framework}
\subsection{Problem statement}
\begin{foldable}
  Given a black-box pre-trained teacher network $T$ and a small set of labeled images $\mathcal{D}=\{x_{i},y_{i}\}_{i=1}^{N}$, our goal is to train a student network $S$ on $\mathcal{D}$ such that $S$ can approximate the performance of $T$.
\end{foldable}

\begin{foldable}
  A direct solution is to use standard \ac{KD}~\cite{hinton2015distilling}, which trains $S$ to match both the ground-truth labels $y_{i}$ and the teacher's predictions $T(x_{i})$ via minimizing:
  \begin{equation}
    \mathcal{L}_\text{KD} = \mathbb{E}_{(x_{i}, y_{i}) \sim \mathcal{D}} \Bigl[{
        (1 - \lambda) \mathcal{L}_{\text{CE}}(S(x_{i}), y_{i}) +
        \lambda \mathcal{L}_\text{KL}(S(x_{i}), T(x_{i}))
    }\Bigr]
    \label{eq:StandardKD}
  \end{equation}
  where $\mathcal{L}_{\text{CE}}$ and $\mathcal{L}_\text{KL}$ are the cross-entropy and the Kullback-Leibler divergence losses, $S(x_{i})$ and $T(x_{i})$ are the student's and teacher's predictions (\ie, class probabilities), and $\lambda$ is a coefficient to balance the losses.
  Note that the \textit{temperature} factor in the original paper is not used as it requires access to the teacher's logits, which violates our black-box teacher assumption.
  Moreover, with a small distillation set $\mathcal{D}$, common \ac{KD} approaches lead to a \textit{sub-optimal} model as they require lots of training images~\cite{hinton2015distilling,kim2018paraphrasing,ahn2019variational,tian2020contrastive}.
\end{foldable}

\begin{foldable}
  \textbf{Proposed Method.}
  To perform black-box few-shot \ac{KD}, we propose our novel method \textbf{DivBFKD}.
  Our method has two phases: Generation and Distillation.
  In the Generation phase, we train a WGAN to produce synthetic images.
  In the Distillation phase, we train the student with the combined real and synthetic images.
\end{foldable}

\subsection{Generation phase.}
\begin{foldable}
  We train a WGAN model~\cite{gulrajani2017improved} with several key changes to generate diverse synthetic images.
  Like standard WGAN, our WGAN consists of a generator $G$ and a discriminator $D$.
  Iteratively, $G$ produces a synthetic image $\tilde{x} = G(z)$ from a random latent vector $z \sim \mathcal{N}(0, \mathbf{I})$.
  Meanwhile, $D$ predicts a scalar score indicating the likelihood of an input image being real.
  Jointly trained with opposite objectives, $G$ tries to fool $D$ by producing realistic images while $D$ tries to best distinguish between real images $x$ and synthetic images $\tilde{x}$.
  These objectives can be formulated into the generator loss $\mathcal{L}_G$ and the discriminator loss $\mathcal{L}_{D}$:
  \begin{equation}
    \mathcal{L}_G = -\mathbb{E}_{z \sim \mathcal{N}(0, \mathbf{I})} [D(G(z))],\label{eq:loss-G-wgan}
  \end{equation}
  \begin{equation}
    \mathcal{L}_{D} = \mathbb{E}_{z \sim \mathcal{N}(0, \mathbf{I})} [D(G(z))] - \mathbb{E}_{x \sim \mathcal{D}}[D(x)].\label{eq:loss-D-wgan}
  \end{equation}
\end{foldable}

\begin{foldable}
  However, if we only use the few-shot set $\mathcal{D}$ to train the WGAN, $G$ will generate synthetic images similar to the images in $\mathcal{D}$~\cite{karras2020training,yang2021data}, which are not diverse enough to facilitate effective distillation.
  To improve diversity of images generated by $G$, we propose the use of synthetic images that the teacher $T$ can predict confidently.
  We refer to these images as \textit{high-confidence} images.
\end{foldable}

\begin{foldable}
  \textbf{High-confidence images.}
  Given a synthetically generated (fake) image $\tilde{x}$, we define its \textit{confidence-score} $c_{\tilde{x}} = \max{k \in \{0, \ldots, K-1\}}T(\tilde{x})$, by taking the maximum over $T$'s predictive probability vector.
  For example, given three classes $\{0, 1, 2\}$ and class probabilities $T(\tilde{x})=[0.1, 0.7, 0.2]$, then $c_{\tilde{x}}=0.7$.
  We define $\tilde{x}$ a high-confidence image if $c_{\tilde{x}}$ is above a confidence-threshold $\tau\in[0,1]$.
  Intuitively, when $c_{\tilde{x}}$ is high, $\tilde{x}$ is likely close to the teacher's training images as it can be confidently classified.
\end{foldable}

\begin{foldable}
  \textbf{Adaptive thresholds.} The critical factor in determining a high-confidence image is the confidence-threshold $\tau$.
  We observe that the teacher has a class-specific bias, \ie, it predicts different classes with different confidence levels.
  Take the dataset FMNIST~\cite{xiao2017fashion} for example, the averaged confidence-score is $>0.99$ for classes `trouser' and `bag', but is only 0.75 for class `shirt' (Fig.~\ref{fig:03-framework-pre_clamp_tau}).
  In this case, fixing a single value $\tau=0.95$ may lead to high-confidence images being too lenient for `trouser' and `bags', but too strict for the class `shirt'.
  Thus, a single $\tau$ for all classes will make high-confidence images imbalanced.
  Further, considering any single class, the confidence-scores have a long-tailed distribution, where they concentrate near 1 and decay with lower values.
  Even though $\tau$ can be set to the averaged confidence-score, it would provide little flexibility and interpretability and also make high-confidence images biased \DIFdelbegin \DIFdel{towards }\DIFdelend \DIFaddbegin \DIFadd{toward }\DIFaddend \textit{outliers}.
\end{foldable}

\begin{figure}[ht]
  \centering
  \includegraphics[viewport=0bp 20bp 455bp 174bp,clip,width=0.75\columnwidth]{./figures/30-research01/33-research01-framework-preclamptau.pdf}
  \caption[DivBFKD confidence-score distributions on FMNIST.]{
    DivBFKD confidence-score distributions across classes on FMNIST.
    They concentrate around 1 and have a long-tail.
    Red bars show adaptive confidence-thresholds $\tau$.
  }
  \label{fig:03-framework-pre_clamp_tau}
\end{figure}

\begin{foldable}
  To address these problems, we use adaptive thresholds $\{\tau^{k}\}_{k=0}^{K-1}$ that is class-specific based on the $q$-quantile of the confidence-scores of real images.
  Specifically, for each of the $K$ classes:
  \begin{equation}
    \tau^{k} = \text{quantile}_{q}(\{c_{x_{i}} \mid(x_{i}, y_{i}) \in \mathcal{D} \text{ and } y_{i} = k\}_{i=1}^{N}), q\in[0,1].
    \label{eq:adaptive-threshold}
  \end{equation}
  We choose $q = 0.1$, \ie, high-confidence images must have teacher's confidence higher than at least 10\% of those for the real ones, as perceived by the teacher $T$.
  This way, $\tau$ adaptively adjusts to any dataset and any class-wise bias from the teacher.
  Fig.~\ref{fig:03-framework-pre_clamp_tau} illustrates $\tau$ as red bars using $q = 0.1$ on FMNIST.
\end{foldable}

\begin{foldable}
  \textbf{WGAN architecture.} We employ a deep convolutional GAN where the generator and discriminator have symmetrical architectures.
  The generator accepts latent vectors of dimension 100 for MNIST, 200 for FMNIST, and 256 for other datasets.
  The latent inputs are projected to 256 base feature maps of dimension 8$\times$8 (except for Imagenette, 7$\times$7), followed by batch normalization.
  These feature maps pass through upscaling convolutional blocks, each consisting of a nearest neighbor (2$\times$) upsample, a convolutional layer with half the channels, batch normalization, and leaky ReLU activation.
  Once the feature maps match the training image dimensions, they are compressed to grayscale or RGB via a final convolutional layer with sigmoid activation, mapping pixel values to $[0, 1]$.

  The discriminator processes input images through convolutional blocks, each with (0.5$\times$) downscaling convolution, batch normalization (except the first block), and leaky ReLU activation.
  The final block contains 256 base feature maps of dimension 8$\times$8 (except Imagenette, 7$\times$7), with each previous block having half the channels.
  These feature maps are passed through a linear layer to produce a scalar score.
\end{foldable}

\begin{foldable}
  \textbf{Our WGAN training scheme.}
  We aim to use the teacher's guidance to promote diversity of synthetic images.
  Fundamentally, synthetic images would be more diverse if the WGAN can learn on images that resemble the teacher's \textit{unknown} training images $\mathcal{D}^{*}$, whose diversity is superior to few-shot images $\mathcal{D}$.
  While there is no direct access to $\mathcal{D}^{*}$, high-confidence images serve as a satisfactory alternative for both quantity (new high-confidence images are generated at every WGAN training step) and quality (they likely reside close to $\mathcal{D}^{*}$, as we previously discussed).
  In other words, high-confidence images implicitly allow our WGAN to learn from $\mathcal{D}^{*}$.
\end{foldable}

\begin{figure}
  \centering
  \includegraphics[width=0.99\columnwidth]{./figures/30-research01/33-research01-framework-teaser.pdf}
  \caption[Illustration of our method DivBFKD.]{
    Illustration of our method DivBFKD.
    \textbf{(a)} \textbf{Generation:} We train our WGAN with the losses $\mathcal{L}_G$ and $\mathcal{L}_{D}^{\text{new}}$.
    When optimizing $G$, we construct the high-confidence set $\mathcal{H}_{n}=\{\tilde{x} = G(z) \mid c_{\tilde{x}} \geq \tau^{\tilde{y}}\}$ from synthetic images $\tilde{x}$, with their confidence-score $c_{\tilde{x}}$, pseudo-label $\tilde{y}$, and adaptive threshold $\tau^{\tilde{y}}$.
    When optimizing $D$, we sample $x$ from the combined real and high-confidence images set $\mathcal{D} \cup \mathcal{H}_{n-1}$.
    \textbf{(b) Distillation:} We augment few images in $\mathcal{D}$ with synthetic images in $\mathcal{D}_G$ to train the student $S$ on $\mathcal{D} \cup \mathcal{D}_G$.
  }
  \label{fig:03-framework-teaser}
\end{figure}

\begin{foldable}
  By this motivation, we introduce high-confidence images in the role of real images to the adversarial learning (Fig.~\ref{fig:03-framework-teaser}(a)), exposing our WGAN to a wider variety of training images.
  At each $n$-th training step:
  \begin{itemize}
    \item {
        The generator $G$ is trained with Eq~\ref{eq:loss-G-wgan}.
        For a fake image $\tilde{x}$, we compute its corresponding confidence-score $c_{\tilde{x}}$ and pseudo-label $\tilde{y} = \argmax_{k \in \{0, \ldots, K-1\}} T(\tilde{x})$.
        Using the adaptive thresholds $\{\tau^{k}\}_{k=0}^{K-1}$ computed with Eq~\ref{eq:adaptive-threshold}, we add qualified $\tilde{x}$ to the \textit{high-confidence} set $\mathcal{H}_{n} = \{\tilde{x} \mid c_{\tilde{x}} \geq \tau^{\tilde{y}}\}$.
        Then, $\mathcal{H}_{n}$ is stored to train the discriminator $D$ in the next step.
      }
    \item {
        Training the discriminator $D$ requires both real and fake images.
        However, our `real' images include real images that actually come from $\mathcal{D}$ and high-confidence images that come from the high-confidence set $\mathcal{H}_{n-1}$.
        Note that $\mathcal{H}_{n-1}$ is constructed in the previous step, and $\mathcal{H}_{0}=\emptyset$.
        With these changes, we train our discriminator with the new loss as:
      }
  \end{itemize}
  \begin{equation}
    \mathcal{L}_{D}^{\text{new}} = {
      \mathbb{E}_{z \sim \mathcal{N}(0, \mathbf{I})}[D(G(z))] -
      \mathbb{E}_{x \sim (\mathcal{D} \cup \mathcal{H}_{n-1})}[D(x)]
    }.
    \label{eq:loss-D-divbfkd}
  \end{equation}
\end{foldable}

\begin{foldable}
  Viewing as an adversarial game, when $D$ learns to appreciate high-confidence synthetic images as realistic, it also encourages $G$ to create more synthetic images like those.
  This interaction forms a positive feedback loop that helps $G$ to be more diverse.
  Furthermore, we emphasize that forming $\mathcal{H}_{n}$ is a \textit{cheap} operator and requires small memory.
  These modifications not only effectively improve the diversity of synthetic images, but also incur only trivial computation overhead.
\end{foldable}

\begin{foldable}
  \textbf{Summary of our novelty in WGAN.}
  We introduce a new training scheme for WGAN using high-confidence images based on the teacher-supervised adaptive thresholds.
  Since we aim to generate images close to the teacher's training images, our synthetic images are more \textit{diverse} and better cover \textit{unseen images}, significantly improving the student's generalization, as shown in \S\ref{sec:33-research01-experiments}.
\end{foldable}

\begin{algorithm}[ht]
  \SetKwInOut{KwIn}{Input}
  \SetKwInOut{KwOut}{Output}
  \SetKwComment{Comment}{}{}

  \caption{Diverse Black-box Few-Shot Distillation (\textbf{DivBFKD}).}
  \label{alg:algorithm-divbfkd}

  \KwIn{teacher $T$; labeled set $\mathcal{D} = \{x_{i}, y_{i}\}_{i=1}^{N}$; quantile $q$; synthetic budget $M$; batch size $B$}
  \KwOut{student network $S$}

  \vspace{-0.5\baselineskip}
  \Comment{\noindent\hrulefill}

  \vspace{0.5\baselineskip}
  Initialize generator $G$, discriminator $D$, student $S$, and high-confidence set $\mathcal{H}_0 \leftarrow \emptyset$ \\
  \Comment{$\triangleright$ \textbf{\textit{Stage 1: Generation}}}
  Infer $T$ on $\mathcal{D}$ to compute adaptive thresholds $\{\tau^k\}_{k=1}^K$ with $q$ \hfill Eq.~\ref{eq:adaptive-threshold} \\
  \While{$G$ not converged}{
    Sample latent vectors $\{z_i\}_{i=1}^B \sim \mathcal{N}(0,I)$ \\
    Generate synthetic images $\{\tilde{x}_i\}_{i=1}^B$ with $G$ \\
    Sample $\{x_i\}_{i=1}^B \sim (\mathcal{D} \cup \mathcal{H}_n)$ as real images \\
    Update $G$ on $\left\{\tilde{x}_i\right\}_{i=1}^B$ via $\mathcal{L}_G$ \hfill Eq.~\ref{eq:loss-G-wgan} \\
    Update $D$ on $\left\{x_i\right\}_{i=1}^B$ and $\left\{\tilde{x}_i\right\}_{i=1}^B$ via $\mathcal{L}_{D}^{new}$ \hfill Eq.~\ref{eq:loss-D-divbfkd} \\
    Compute pseudo-labels $\{\tilde{y}_i\}$ and confidence-scores $\{c_{\tilde{x}_i}\}$ for $\{\tilde{x}_i\}$ \\
    Construct $\mathcal{H}_n=\left\{\tilde{x_{i}} \mid c_{\tilde{x}_i}\ge\tau^{\tilde{y}_i}\right\}_{i=1}^B$ \\
  }

  \vspace{0.5\baselineskip}
  \Comment{$\triangleright$ \textbf{\textit{Stage 2: Distillation}}}
  Construct $\mathcal{D}_\text{KD}$ with $N$ real and $M$ synthetic images \hfill Eq.~\ref{eq:distill-set} \\
  \While{$S$ not converged}{
    Update $S$ on $\mathcal{D}_\text{KD}$ via $\mathcal{L}_\text{KD}$ \hfill Eq.~\ref{eq:loss-S}
  }

  \vspace{0.5\baselineskip}
  \Return{$S$}
\end{algorithm}

\subsection{Distillation phase}
\begin{foldable}
  After training WGAN in the Generation phase, we use it to generate synthetic images.
  We employ our generator $G$ to build a set of synthetic images $\mathcal{D}_G$, and their pseudo-labels (\ie, class probabilities) are obtained via the teacher $T$.
  Following BBKD~\cite{wang2020neural} and FS-BBT~\cite{nguyen2022black}, we construct a distillation set $\mathcal{D}_\text{KD}$ comprising of $N$ real images from $\mathcal{D}$ and $M$ synthetic images from $\mathcal{D}_G$:
  \begin{equation}
    \mathcal{D}_\text{KD} = {
      \{x_{i} \in \mathcal{D}\}_{i=1}^{N} \cup
      \{G(z_{j}) | z_{j} \sim \mathcal{N}(0, \mathbf{I})\}_{j=1}^{M}
    },
    \label{eq:distill-set}
  \end{equation}
  where $x_{i} \in \mathcal{D}$ and $\tilde{x}_{j} \in \mathcal{D}_G$ are real and synthetic images.
  Like other few-shot \ac{KD} methods~\cite{wang2020neural,nguyen2022black}, we train the student $S$ with $\mathcal{D}_\text{KD}$ using a cross-entropy loss:
  \begin{equation}
    \mathcal{L}_\text{KD} = \mathbb{E}_{x_{i} \sim \mathcal{D}_\text{KD}} \mathcal{L}_\text{CE}(S(x_{i}), T(x_{i})).
    \label{eq:loss-S}
  \end{equation}
\end{foldable}

\begin{foldable}
  \textbf{Class balancing.} Inspired by~\cite{chen2019data,zhang2022ideal}, we generate a balanced set of synthetic images.
  Recalling that $M$ is the budget of synthetic images for \ac{KD} and $K$ is the number of classes, each class would have $M / K$ synthetic images.
  We ensure class balance through \textit{rejection sampling} until we achieve the sufficient amount.
\end{foldable}

We summarize our method in Algorithm \ref{alg:algorithm-divbfkd}.
 \section{Experiments} \label{sec:33-research01-experiments}
\begin{foldable}
  We conduct a wide range of experiments to show our improvements over current few-shot \ac{KD} baselines.
  We also include cross-architecture \ac{KD}, visualization, ablation studies, and comparison with data-free \ac{KD} methods.
\end{foldable}

\subsection{Baselines}
\begin{foldable}
  We compare our method DivBFKD with the following baselines:
  \begin{itemize}
    \item \textit{Student-Full}: The student is trained on the whole teacher's training set.
    \item \textit{Student-Alone}: The student is trained on the few-shot set $\mathcal{D}$.
    \item {
        \textit{Standard-KD}~\cite{hinton2015distilling}: The student is trained with the standard \ac{KD} loss (Eq.~\ref{eq:StandardKD}) on $\mathcal{D}$.
        We select temperature $\lambda=0.9$, a standard value used in \ac{KD} methods~\cite{hinton2015distilling,tian2020contrastive,yuan2020revisiting,ma2021undistillable}.
      }
    \item \textit{Few-shot \ac{KD}} methods: We compare with two white-box methods, FSKD~\cite{kimura2018few} and WaGe~\cite{kong2020learning}, and two black-box methods, BBKD~\cite{wang2020neural} and FS-BBT~\cite{nguyen2022black}.
  \end{itemize}
\end{foldable}

\begin{foldable}
  For fair comparisons, we use the same teacher-student network architecture and the same numbers of real and synthetic images $N$ and $M$ as in other few-shot \ac{KD} methods~\cite{kimura2018few,wang2020neural,nguyen2022black}.
  Their accuracy numbers are obtained from~\cite{nguyen2022black}\footnote{This is possible because we use benchmark datasets with fixed train and test splits.}.
  For our method, we set the quantile $q = 0.1$ to select high-confidence images across all experiments.
  We repeat each experiment three times with random seeds and report the averaged accuracy with standard deviation on the hold-out test set.
\end{foldable}

\subsection{Implementation Details}

\subsubsection{Datasets}
\begin{table}[ht]
  \centering
  \begin{threeparttable}
    \caption[Dataset details in DivBFKD.]{
      Dataset details in DivBFKD.
      $N$ ($\times 10^3$) denotes the number of images in thousands and $K$ denotes the number of classes.
      Dimension (images and transforms) are in pixels (px).
    }
    \label{tab:Dataset-details}
    \begin{tabular}{c c c c c}
      \toprule
      \textbf{Dataset} & $N~(\times 10^3)$ & $K$ & \textbf{Dimension} & \textbf{Transforms} \\
      \midrule
      MNIST         & 60  & 10   & $1\times 28 \times 28$          & Resize to 32$\times$32 \\
      FMNIST       & 60  & 10   & $1\times 28 \times 28$           & Resize to 32$\times$32 \\
      SVHN          & 73  & 10   & $3\times 32 \times 32$          & $\Bigl[$ \makecell{Zero-padding (4 px/side) \\ Random crop to $32\times32$} $\Bigr]$\\
      CIFAR10       & 50  & 10   & $3\times 32 \times 32$          & $\Biggl[$ \makecell{Zero-padding (4 px/side) \\ Random horizontal flip \\ Random crop to $32\times32$} $\Biggr]$\\
      CIFAR100      & 50  & 100  & $3\times 32 \times 32$          & $\Biggl[$ \makecell{Zero-padding (4 px/side) \\ Random horizontal flip \\ Random crop to $32\times32$} $\Biggr]$\\
      Tiny-ImageNet & 100 & 200  & $3\times 64 \times 64$          & $\Biggl[$ \makecell{Zero-padding (8 px/side) \\ Random horizontal flip \\ Random crop to $64\times64$} $\Biggr]$\\
      Imagenette    & 9.5 & 10   & $3\times 408 \times 472 ^{(*)}$ & $\Biggl[$ \makecell{Resize to $256\times256$ \\ Random horizontal flip \\ Random crop to $224\times224$} $\Biggr]$\\
      \bottomrule
    \end{tabular}
    \begin{tablenotes}[flushleft] \footnotesize
    \item {[$*$] Imagenette dimensions are reported as mean values}
    \end{tablenotes}
  \end{threeparttable}
\end{table}

\begin{foldable}
  We evaluate our method across seven benchmark image datasets: MNIST, FMNIST, SVHN, CIFAR10, CIFAR100, Tiny-ImageNet, and Imagenette~\cite{lecun1998gradient,xiao2017fashion,netzer2011reading,le2015tiny,krizhevsky2009learning,howard2020imagenette}.
  These datasets cover different difficulty levels, including: (1) \textit{Simple} datasets with MNIST, FMNIST, SVHN, and CIFAR10 (10 classes) and (2) \textit{Complex} datasets with CIFAR100 (100 classes), Tiny-ImageNet (200 classes), and Imagenette (high-resolution at 224\texttimes 224).
\end{foldable}

\begin{foldable}
  Following~\cite{wang2020neural,nguyen2022black}, to construct the few-shot dataset $\mathcal{D}$, we used $N = 2000$ images for MNIST, FMNIST, SVHN, CIFAR10, and Imagenette, $N=5000$ images for CIFAR100, and $N=10000$ images for Tiny-ImageNet.
  They were randomly drawn from the original training set such that each class had $N/K$ images ($K$ is the number of classes).
\end{foldable}

\subsubsection{Teacher and Student}
\begin{foldable}
  We consider several popular network architectures for distillation, namely LeNet5~\cite{lecun1998gradient}, AlexNet~\cite{krizhevsky2012imagenet}, and ResNet~\cite{he2016deep}, and VGG~\cite{simonyan2014very}.
  These architectures are widely used to evaluate previous few-shot \ac{KD} methods~\cite{kimura2018few,wang2020neural,nguyen2022black}.
\end{foldable}

\begin{foldable}
  \textbf{Teacher training.} We trained the teacher network with a stochastic gradient descent (SGD) optimizer using 0.9 momentum, $5\times10^{-4}$ weight decay, and 128 batch size.
  We fixed the learning rate to $10^{-2}$ on MNIST and FMNIST.
  On other datasets, we scheduled the learning rate to $10^{-1}$ at the beginning and decreased it by ten times at 50\% and 75\% progress.
  We trained for 100 epochs on MNIST and FMNIST, and 200 epochs on other datasets.
\end{foldable}

\begin{foldable}
  \textbf{Student training.} For a fair comparison, we used the same number of synthetic images for distillation as in other few-shot \ac{KD} methods~\cite{kimura2018few,wang2020neural,nguyen2022black}, which is 24\,K for MNIST, 48\,K for FMNIST, 40\,K for SVHN, CIFAR10, CIFAR100, and Imagenette, and 50\,K for Tiny-ImageNet.
  We trained the Student with an SGD optimizer using 0.9 momentum and $5 \times 10^{-4}$ weight decay.
  We fixed the learning rate to $10^{-2}$ on MNIST and FMNIST.
  On other datasets, we scheduled the learning rate to $10^{-1}$ at the beginning and decrease it by ten times at 50\% and 75\% distillation progress.
  We used a batch size of 250 and the number of epochs of 100 for MNIST, 200 for FMNIST and SVHN, 400 for CIFAR10 and Imagenette, and 1000 for CIFAR100 and Tiny-ImageNet.
\end{foldable}

\subsubsection{WGAN}
\begin{foldable}
  \textbf{Architecture.} We employed a deep convolutional GAN as in~\cite{radford2015unsupervised}, where the generator and discriminator had symmetrical architectures.
  The generator accepts latent vectors of dimension 100 for MNIST, 200 for FMNIST, and 256 for other datasets.
  The latent inputs are then projected to 256 base feature maps of dimension 8$\times$8 (except for Imagenette, 7$\times$7), which then undergo batch normalization.
  These feature maps are passed through upscaling convolutional blocks, each consisting of a nearest neighbor (2$\times$) upsample layer, a convolutional layer with half the channels, a batch normalization, and a leaky ReLU activation.
  Finally, once the feature maps have identical dimension to that of the training images, they are compressed to either grayscale or RGB with a final convolutional layer and a sigmoid activation to map the output pixel values to range $[0, 1]$.

  The discriminator feeds forward input images through convolutional blocks, each consisting of a (0.5$\times$) downscaling convolutional layer, a batch normalization (except for the first block), and a leaky ReLU activation.
  While the final block has 256 base feature maps of dimension 8$\times$8 (except for Imagenette, 7$\times$7), each previously adjacent block has half the channels.
  These base feature maps are eventually passed through a linear layer to yield a scalar score.
\end{foldable}

\begin{foldable}
  \textbf{Training.} We followed the implementation in~\cite{gulrajani2017improved} to train our WGAN~\cite{arjovsky2017wasserstein}.
  As recommended, we enforced the Lipschitz constraint on the discriminator with gradient penalty, whose coefficient with the Wasserstein loss was set to 10:1.
  The number of discriminator updates per generator updates was set to 5.
  However, we alternatively used the RMSProp optimizer with learning rate $5 \times 10^{-5}$ and batch size 250 for both models.
  We used the number of epochs of 500 for MNIST, 1000 for FMNIST and SVHN, 2000 for CIFAR10 and Imagenette, and 5000 for CIFAR100 and Tiny-ImageNet.
\end{foldable}

\subsection{Distillation performance}

\subsubsection{Simple datasets} \label{sec:experiments-simple-datasets}
\begin{table}[ht]
  \centering
  \begin{threeparttable}
    \caption[DivBFKD accuracy on Simple datasets.]{
      DivBFKD distillation accuracy (\%) on Simple datasets.
      $N$ and $M$ show the budget of real and synthetic images.
    }
    \label{tab:04-experiments-results_few_classes}
    \begin{tabular}{l l l l l}
      \toprule
      \textbf{Dataset} & \textbf{MNIST} & \textbf{FMNIST} & \textbf{SVHN} & \textbf{CIFAR10} \\
      $(N,M)$ & (2\,K, 24\,K) & (2\,K, 48\,K) & (2\,K, 40\,K) & (2\,K, 40\,K) \\
      \midrule
      Teacher                   & 99.28                   & 90.90                   & 96.16                   & 90.05 \\
      Student-Full              & \meanstd{98.91}{0.08}   & \meanstd{88.68}{0.53}   & \meanstd{95.61}{0.07}   & \meanstd{87.58}{0.36} \\
      \midrule

      \multicolumn{5}{c}{\textbf{Smaller Architecture}} \\
      \midrule
      Student-Alone             & \meanstd{95.14}{0.19}   & \meanstd{79.24}{1.05}   & \meanstd{85.68}{0.46}   & \meanstd{59.02}{0.54} \\
      Standard-KD\blackbox      & \meanstd{95.36}{0.11}   & \meanstd{80.76}{0.37}   & \meanstd{89.11}{0.22}   & \meanstd{59.77}{0.96} \\
      FSKD\whitebox             & 80.43\dg                & 68.64\dg                & $-$                     & 40.58\dg \\
      WaGe\whitebox             & $-$                     & 85.18\dg                & $-$                     & 73.08\dg \\
      FS-BBT\blackbox           & 98.42\dg                & 84.73\dg                & $-$                     & 74.10\dg \\
      \textbf{DivBFKD}\blackbox & \meanstdbf{98.58}{0.06} & \meanstdbf{86.50}{0.18} & \meanstdbf{94.47}{0.12} & \meanstdbf{76.97}{0.56} \\
      \midrule

      \multicolumn{5}{c}{\textbf{Same Architecture}} \\
      \midrule
      BBKD\blackbox             & 98.74\dg                & 80.90\dg                & $-$                     & 74.60\dg \\
      FS-BBT\blackbox           & 98.91\dg                & 86.53\dg                & $-$                     & 76.17\dg \\
      \textbf{DivBFKD}\blackbox & \meanstdbf{98.97}{0.04} & \meanstdbf{87.65}{0.37} & \meanstdbf{94.61}{0.08} & \meanstdbf{78.07}{0.26} \\
      \bottomrule
    \end{tabular}
    \begin{tablenotes}[flushleft] \footnotesize
    \item {
        [(W)] white-box; [(B)] black-box; [$-$] not reported; [$\dagger$] reported from respective paper; [\textbf{bold}] best results.
      }
    \end{tablenotes}
  \end{threeparttable}
\end{table}

\begin{foldable}
  Following~\cite{nguyen2022black}, we conduct distillation in two settings:
  \begin{enumerate}
    \item \textit{Smaller} architecture: The student has a smaller architecture than the teacher: LeNet5 to LeNet5-Half on MNIST, FMNIST; and AlexNet to AlexNet-Half on SVHN, CIFAR10.\footnote{The \textit{-Half} version has half the channels in each layer.}
    \item \textit{Same} architecture: The student has an identical architecture to the teacher.
  \end{enumerate}
\end{foldable}

\begin{foldable}
  We report the results in Table~\ref{tab:04-experiments-results_few_classes} and make the following observations:
  \begin{itemize}
    \item Our method DivBFKD improves by around 3-6\% over Standard-KD on MNIST, FMNIST, and SVHN and by significantly 17\% on CIFAR10.
    \item {
        Compared to other few-shot \ac{KD} methods, our DivBFKD also excels in both smaller and same architecture settings.
        It improves upon second-best FS-BBT by around 2\% on FMNIST and CIFAR10.
      }
    \item {
        The performance gap between our DivBFKD and its upper bound Student-Full is very close on MNIST (0.33\%), FMNIST (2.18\%), and SVHN (1.14\%).
        We emphasize that Student-Full uses the teacher's training set ($\ge$ 60\,K real images), much larger than the 2\,K real images we started with.
      }
  \end{itemize}
\end{foldable}

\subsubsection{Complex datasets} \label{sec:experiments-complex-datasets}
\begin{foldable}
  Similar to~\cite{wang2020neural,nguyen2022black,kong2020learning}, we distill from ResNet32 to ResNet20 on CIFAR100 and Tiny-ImageNet (Table~\ref{tab:ResultsHundredClasses}).
  With limited data, Student-Alone loses roughly half of its accuracy compared to Student-Full.
  As a naive solution, Standard-KD improves over Student-Alone 10\% on CIFAR100 and 14\% on Tiny-ImageNet.
  Ultimately, our method DivBFKD is the best method, exceeding the second-best FS-BBT by 3\% on CIFAR100 and 1\% on Tiny-ImageNet.
\end{foldable}

\begin{foldable}
  We further evaluate our method on high-resolution images with Imagenette.
  On 2\,K real images, Student-Alone achieves 76.80\% while Standard-KD marginally improves to 76.94\%.
  DivBFKD further improves upon this result by 10\%, approximating Student-Full.
\end{foldable}

\begin{table}[ht]
  \centering
  \begin{threeparttable}
    \caption[DivBFKD accuracy on Complex datasets.]{
      DivBFKD distillation accuracy (\%) on Complex datasets.
      $N$ and $M$ show the budget of real and synthetic images.
    }
    \label{tab:ResultsHundredClasses}
    \begin{tabular}{l l l l}
      \toprule
      \textbf{Dataset} & \textbf{CIFAR100} & \textbf{Tiny-ImageNet} & \textbf{Imagenette} \\
      $(N,M)$                   & (5\,K, 40\,K)           & (10\,K, 50\,K)        & (2\,K, 40\,K) \\
      \midrule
      Teacher                   & 71.41                   & 56.27                 & 90.47 \\
      Student-Full              & \meanstd{68.28}{0.20}   & \meanstd{51.51}{0.31} & \meanstd{89.89}{0.41} \\
      \midrule
      Student-Alone             & \meanstd{35.71}{0.39}   & \meanstd{26.59}{0.25} & \meanstd{76.80}{0.78} \\
      Standard-KD\blackbox      & \meanstd{45.76}{0.48}   & \meanstd{40.42}{0.19} & \meanstd{76.94}{0.14} \\
      WaGe\whitebox             & 20.32\dg                & $-$                   & $-$ \\
      BBKD\blackbox             & 53.41\dg                & 40.01\dg              & $-$ \\
      FS-BBT\blackbox           & 56.28\dg                & 43.29\dg              & $-$ \\
      \textbf{DivBFKD}\blackbox & \meanstdbf{59.91}{0.31} & \meanstdbf{44.25}{0.38} & \meanstdbf{86.48}{0.72} \\
      \bottomrule
    \end{tabular}
    \begin{tablenotes}[flushleft] \footnotesize
    \item[] {
        [(W)] white-box; [(B)] black-box; [$-$] not reported; [$\dagger$] reported from respective paper; [\textbf{bold}] best results.
      }
    \end{tablenotes}
  \end{threeparttable}
\end{table}

\subsubsection{Cross-architecture distillation} \label{sec:experiments-crossarch}
\begin{foldable}
  In previous subsections, we distill from the teacher to a student networks of similar architectures to have a consistent setting with other few-shot \ac{KD} methods (\eg, LeNet5 to LeNet5-Half, AlexNet to AlexNet-Half, and ResNet32 to ResNet20).
  This section further explores the cases where the teacher and student architectures belong to different families: ResNet~\cite{he2016deep}, VGG~\cite{simonyan2014very}, and AlexNet~\cite{krizhevsky2012imagenet}.
  We adopt the same setting from the standard experiment on CIFAR10 from \S\ref{sec:experiments-complex-datasets}.
  Table~\ref{tab:Cross-architecture} shows the performance of our method DivBFKD with three teachers ResNet32, VGG16, and AlexNet combined with three students ResNet20, VGG11, and AlexNet-Half on CIFAR10.
\end{foldable}

\begin{table}[ht]
  \centering
  \begin{threeparttable}
    \caption[DivBFKD cross-architecture accuracy on CIFAR10.]{
      DivBFKD cross-architecture distillation accuracy (\%) on CIFAR10.
    }
    \label{tab:Cross-architecture}
    \begin{tabular}{c c|c c c c c}
      \toprule
      \multicolumn{2}{c|}{\textbf{Architecture}} & \multicolumn{5}{c}{\textbf{Accuracy}}\\
      Teacher & Student & Teacher & S-Full & S-Alone & Standard-KD & DivBFKD \\
      \midrule
      ResNet32 & ResNet20     & 93.34 & & & \meanstd{65.86}{1.90} & \meanstd{84.05}{0.10} \\
      VGG16    & ResNet20     & 91.30 & \meanstd{92.34}{0.06} & \meanstd{64.52}{1.60} & \meanstd{64.76}{0.89} & \meanstd{80.69}{0.48} \\
      AlexNet  & ResNet20     & 90.05 & & & \meanstd{65.29}{0.89} & \meanstd{79.57}{0.16} \\
      \midrule
      ResNet32 & VGG11        & 93.34 & & & \meanstd{53.78}{0.78} & \meanstd{69.42}{0.83} \\
      VGG16    & VGG11        & 91.30 & \meanstd{89.12}{0.17} & \meanstd{53.34}{0.57} & \meanstd{53.80}{1.21} & \meanstd{75.41}{0.06} \\
      AlexNet  & VGG11        & 90.05 & & & \meanstd{53.51}{0.14} & \meanstd{74.60}{0.29} \\
      \midrule
      ResNet32 & AlexNet-Half & 93.34 & & & \meanstd{60.36}{0.82} & \meanstd{74.70}{0.31} \\
      VGG16    & AlexNet-Half & 91.30 & \meanstd{87.58}{0.36} & \meanstd{59.02}{0.54} & \meanstd{60.47}{0.10} & \meanstd{76.16}{0.57} \\
      AlexNet  & AlexNet-Half & 90.05 & & & \meanstd{59.77}{0.96} & \meanstd{76.97}{0.56} \\
      \bottomrule
    \end{tabular}
    \begin{tablenotes}[flushleft] \footnotesize
    \item {Same teacher-student pairs share the same Student-Full and Student-Alone accuracy.}
    \end{tablenotes}
  \end{threeparttable}
\end{table}

\begin{foldable}
  We have the following observations:
  \begin{itemize}
    \item Among three architecture families, ResNet is the best one with a 93.34\% Teacher, 92.34\% Student-Full, and 64.52\% Student-Alone.
    \item Standard-KD is helpful and it is always better than Student-Alone in all teacher-student combinations.
    \item {
        VGG is better than AlexNet when the teacher and the student are trained on the whole training set.
        However, VGG student is worse than AlexNet student when it is trained on the few-shot dataset.
        Namely, Student-Alone of VGG11 has 53.34\% accuracy while Student-Alone of AlexNet-Half has 59.02\% accuracy.
      }
    \item Our method DivBFKD remarkably improves accuracy over Standard-KD by 14--22\% on all teacher-student combinations.
    \item Distillation is often most effective when the teacher's and student's architecture belong to the same family.
  \end{itemize}
  Overall, the robust performance across multiple data scales and architectures suggests that our method successfully tackles the black-box few-shot \ac{KD} setting.
\end{foldable}

\subsection{Diversity of our synthetic images}

\subsubsection{Comparison with the teacher's training images.}

\begin{foldable}
  As we mentioned earlier, our goal is to generate synthetic images close to the teacher's training images.
  To prove this argument, we embed and visualize the teacher's training images (they are \textit{unknown} and not used in the \ac{KD} process) and our synthetic images of CIFAR10 in Fig.~\ref{fig:Distributions-of-CIFAR10}.
  It shows that our synthetic images belong to the distribution of the teacher's training images.
  With the teacher's supervision, our synthetic images show their strong resemblance to the teacher's training images.
\end{foldable}

\begin{foldable}
  For quantitative results, we compute the \textit{Coverage} metric~\cite{alaa2022faithful} to measure the fraction of teacher's training samples whose neighborhoods contain at least one synthetic sample.
  Our synthetic images achieve 0.42, compared to 0.18 of the few-shot images and 0.29 of the synthetic images generated from the standard WGAN.
  This shows that our synthetic images resemble many variability of the teacher's training images, thus improving image diversity.
\end{foldable}

\begin{figure}[ht]
  \centering
  \includegraphics[width=1\columnwidth]{./figures/30-research01/34-research01-experiments-viz-embed2d.png}
  \caption[DivBFKD image embeddings on CIFAR10.]{
    DivBFKD image embeddings of teacher's training images, few-shot images, and synthetic images on CIFAR10.
  }
  \label{fig:Distributions-of-CIFAR10}
\end{figure}

\begin{figure}[ht]
  \centering
  \caption[DivBFKD visualization of real and synthetic images on CIFAR10.]{
    DivBFKD visualization of real and synthetic images along with their predictive labels and confidence-scores provided by the teacher on CIFAR10.
  }
  \label{fig:Visualization}
  \begin{subfigure}[t]{0.475\columnwidth}
    \centering
    \includegraphics[width=\linewidth]{./figures/30-research01/34-research01-experiments-viz-real.pdf}
    \caption{Real images}
  \end{subfigure}
  \hfill
  \begin{subfigure}[t]{0.475\columnwidth}
    \centering
    \includegraphics[width=\linewidth]{./figures/30-research01/34-research01-experiments-viz-synth.pdf}
    \caption{Synthetic images}
  \end{subfigure}
\end{figure}

\subsubsection{Quality and quantitative measurements}
\begin{foldable}
  We randomly visualize 24 real images and 24 synthetic images generated by our WGAN in Fig.~\ref{fig:Visualization}.
  Primarily, our goal is not to generate visually aesthetic images but to generate diverse images that the teacher confidently predicts their labels.
  As a side note, image quality is also not a strict requirement in \ac{KD} methods relied on synthetic images~\cite{chen2019data,wang2020neural,wang2021zero}.
  However, our synthetic images in Fig.~\ref{fig:Visualization}(b) are able to show convincing properties of their real counterparts.
\end{foldable}

\begin{foldable}
  \textbf{Inception score (IS) and Fr\'{e}chet inception distance (FID).} To quantify the diversity and quality of synthetic images, we compute the IS~\cite{salimans2016improved} (higher is better) and FID~\cite{heusel2017gans} (lower is better) of our synthetic images in Table~\ref{tab:IS-FID}.
  We compare with the standard WGAN and FS-BBT (the second-best baseline), where our method achieves the best scores.
  These results confirm that our method generate synthetic images with high diversity and quality, and explain their effectiveness for the distillation step.
\end{foldable}

\begin{table}[ht]
  \centering
  \begin{threeparttable}
    \caption{DivBFKD synthetic image quality on CIFAR10.}
    \label{tab:IS-FID}
    \begin{tabular}{l c c}
      \toprule
      \textbf{Method} & \textbf{IS ($\uparrow$)} & \textbf{FID ($\downarrow$)} \\
      \midrule
      CVAE (from~\cite{nguyen2022black}) & 3.51          & 19.91 \\
      Standard WGAN                      & 3.99          & 14.79 \\
      WGAN with AT (Ours)                & \textbf{4.42} & \textbf{13.96} \\
      \bottomrule
    \end{tabular}
    \begin{tablenotes}[flushleft] \footnotesize
    \item {[AT] adaptive thresholds.}
    \end{tablenotes}
  \end{threeparttable}
\end{table}

\subsection{Ablation studies}

\subsubsection{Impacts of different components}
\begin{foldable}
  As described in \S\ref{sec:32-research01-framework}, the core components of our method include WGAN and adaptive thresholds in the generation phase, and class balancing in the distillation phase.
  Adding both latter components to WGAN completes our method.
  Table~\ref{tab:Ablation-study-Components} breaks down the contribution of each component, where both adaptive thresholds (AT) and class balancing (CB) improve our accuracy when applied either independently or jointly.
\end{foldable}

\begin{table}[ht]
  \centering
  \begin{threeparttable}
    \caption[DivBFKD ablation study: incremental components.]{
      DivBFKD ablation study: student accuracy (\%) with incremental components.
    }
    \label{tab:Ablation-study-Components}
    \begin{tabular}{l c c}
      \toprule
      \textbf{Components}                 & \textbf{CIFAR10}        & \textbf{Tiny-ImageNet} \\
      \midrule
      WGAN~\cite{arjovsky2017wasserstein} & \meanstd{75.59}{0.69}   & \meanstd{42.70}{0.25} \\
      WGAN with AT                        & \meanstd{76.35}{0.21}   & \meanstd{43.25}{0.71} \\
      WGAN with CB                        & \meanstd{76.88}{0.45}   & \meanstd{43.75}{0.40} \\
      \textbf{DivBFKD}                    & \meanstdbf{76.97}{0.56} & \meanstdbf{44.25}{0.38} \\
      \bottomrule
    \end{tabular}
    \begin{tablenotes}[flushleft] \footnotesize
    \item {[AT] adaptive thresholds; [CB] class balancing.}
    \end{tablenotes}

  \end{threeparttable}
\end{table}

\subsubsection{Impacts of the numbers of real and synthetic images}
\begin{foldable}
  We investigate our performance with varying amounts of data on CIFAR10.
  Fig.~\ref{fig:34-research01-experiments-abla_data}(a) shows that both Standard-KD and our DivBFKD improve student accuracy with more real images, as expected.
  At $N$ = 5\,K, our DivBFKD achieves 79.93\%, while Standard-KD only achieves 72.66\%.
  Conversely, at $N$ = 250, Standard-KD drops abruptly to 33.33\%, while our method still has an acceptable 66.76\%.
  This trend proves that our method can still distill robustly in extremely limited data cases.
  Besides, we also expect that our performance should improve with more synthetic images.
  This agrees with Fig.~\ref{fig:34-research01-experiments-abla_data}(b) where our DivBFKD clearly outperforms Standard-KD and improves with larger $M$.
\end{foldable}

\begin{figure}[ht]
  \centering
  \includegraphics[width=0.485\columnwidth]{./figures/30-research01/34-research01-experiments-ablation-N.pdf}
  \hspace{0.2cm}
  \includegraphics[width=0.485\columnwidth]{./figures/30-research01/34-research01-experiments-ablation-M.pdf}
  (a)\hspace{0.49\columnwidth}(b)
  \caption[DivBFKD accuracy versus image budget on CIFAR10.]{
    DivBFKD accuracy versus the budget of real ($N$) and synthetic images ($M$) on CIFAR10.
  }
  \label{fig:34-research01-experiments-abla_data}
\end{figure}

\subsubsection{Impact of the quantile}
\begin{foldable}
  In our method, the quantile $q$ is principally used to compute adaptive thresholds to determine high-confidence images.
  Contrarily, while too small $q$ might introduce counter-productive noise, too large $q$ does not allow enough high-confidence images to improve diversity.
  We achieve a relatively stable accuracy of 76.22--76.97\% with $q\in[0.03,0.10]$ (Fig.~\ref{fig:Results-Hyperparameter-tuning}) and recommend this range for a good trade-off between image diversity and training stability.
\end{foldable}

\begin{figure}[ht]
  \centering
  \includegraphics[width=0.75\columnwidth]{./figures/30-research01/34-research01-experiments-hyperparam-q.pdf}
  \caption[DivBFKD accuracy versus quantile on CIFAR10.]{
    DivBFKD accuracy versus the quantile $q$ on CIFAR10.
  }
  \label{fig:Results-Hyperparameter-tuning}
\end{figure}

\subsection{Comparison with data-free KD methods} \label{ssec:04-experiments-comparison_zero_shot}
\begin{foldable}
  We also compare our method with several popular data-free (zero-shot) \ac{KD} methods, including four white-box methods: Meta-KD~\cite{lopes2017data}, ZSKD~\cite{nayak2019zero}, DAFL~\cite{chen2019data}, and DFKD~\cite{wang2021data}, and two black-box methods ZSDB3KD~\cite{wang2021zero} and IDEAL~\cite{zhang2022ideal}.
  From Table~\ref{tab:04-experiments-comparison_data_free}, our method DivBFKD can successfully take advantage of a small amount of data (2\,K images) and perform much better than data-free methods on FMNIST and CIFAR10, while it is comparable on MNIST.
\end{foldable}

\begin{table}[ht]
  \centering
  \begin{threeparttable}
    \caption[DivBFKD accuracy compared with zero-shot KD methods.]{
      DivBFKD distillation accuracy (\%) compared with zero-shot KD methods.
    }
    \label{tab:04-experiments-comparison_data_free}
    \begin{tabular}{l|c|c|c}
      \hline
      \textbf{Dataset} & \textbf{MNIST} & \textbf{FMNIST} & \textbf{CIFAR10}  \\
      \hline
      Meta-KD\whitebox          & 92.47          & $-$            & $-$ \\
      ZSKD\whitebox             & 98.77          & 79.62          & 69.56 \\
      DAFL\whitebox             & 98.20          & $-$            & 66.38 \\
      DFKD\whitebox             & \textbf{99.08} & $-$            & 73.91 \\
      ZSDB3KD\blackbox          & 96.54          & 72.31          & 59.46 \\
      IDEAL\blackbox            & 96.32          & 83.92          & 65.61 \\
      \textbf{DivBFKD}\blackbox & 98.97          & \textbf{87.65} & \textbf{78.07} \\
      \hline
    \end{tabular}
    \begin{tablenotes}[flushleft] \footnotesize
    \item {[(W)] white-box; [(B)] black-box; [$-$] not reported; [\textbf{bold}] best results.}
    \end{tablenotes}
  \end{threeparttable}
\end{table}
 \section{Conclusion} \label{sec:34-research01-conclusion}

\begin{foldable}
  While a large training set and internal access to the teacher are the prerequisites to most knowledge distillation (KD) approaches, they are rarely available due to various restrictions.
  To address these practical challenges, we propose DivBFKD, a novel method to bridge the gap for this black-box few-shot \ac{KD} problem.
  We investigate the diversity of training data, a crucial but inadequately addressed factor for \ac{KD}.
  Specifically, we propose the use of high-confidence images---ones that the teacher can classify confidently.
  We also define adaptive thresholds, the teacher-supervised criteria to negate teacher's class-wise bias during the selection of these images.
  Under the teacher's supervision, we introduce high-confidence images to the WGAN training loop on-the-fly to boost the diversity of image generation process with negligible cost.
  Through comprehensive evaluation, we have shown that our DivBFKD could use diversity to achieve \DIFdelbegin \DIFdel{state-of-the-art }%DIFDELCMD < \ac{KD} %%%
\DIFdel{performance among other }\DIFdelend \DIFaddbegin \DIFadd{the highest accuracy among the compared }\DIFaddend few-shot \ac{KD} methods \DIFaddbegin \DIFadd{(BBKD and FS-BBT) }\DIFaddend on multiple image datasets and architectures at different scales.
  As a \DIFaddbegin \DIFadd{limitation, the adversarial training of the WGAN remains sensitive to the balance between generator and discriminator, and we have not characterized this at larger resolutions.
  As a }\DIFaddend future work, it would be interesting to study the robustness aspects of the distilled model relative to the original model to have the assurance~\cite{gopakumar2018algorithmic} for successful deployment in the real world.
\end{foldable}
 \clearpage{}
\clearpage{}\acresetall

\chapter[Diverse Image Priors for Black-box Data-free Knowledge Distillation]
{Diverse Image Priors \\ for Black-box Data-free Knowledge Distillation} \label{ch:40-research02}

\begin{foldable}
  Section~\ref{sec:24-background-kdrestricted-bbdatafree} identifies two structural deficits in black-box data-free \ac{KD}: the domain-coverage deficit (without real images to anchor the generator, the synthesized distribution is not constrained to cover the teacher's training distribution) and the information-bottleneck deficit (hard top-1 labels discard the inter-class relational structure that soft targets encode).
  Prior methods~\cite{wang2021zero, zhang2022ideal, yuan2024data} resolve at most one deficit independently; closing the gap requires a unified approach that enforces diversity during synthesis while recovering soft signals during distillation.
  This chapter answers the open problem posed at the end of \S\ref{sec:24-background-kdrestricted-bbdatafree} by proposing \acl{DIPKD}: a three-phase pipeline where hierarchical noise and contrastive learning enforce diverse image synthesis, and a primer student supplies soft surrogate logits to restore the dark-knowledge signal lost to black-box access.
\end{foldable}

\section{Introduction} \label{sec:41-research02-introduction}

\begin{foldable}
  The dual challenge of domain-coverage deficit and information-bottleneck deficit is qualitatively different from the few-shot setting addressed in Chapter~\ref{ch:30-research01}.
  Few-shot black-box \ac{KD} at least provides a distributional anchor: real training images from the teacher's domain, even if limited in diversity.
  Data-free black-box \ac{KD} removes this anchor entirely.
  The generator must reconstruct the image manifold from hard-label queries alone, with no real samples to guide synthesis and no soft targets to preserve inter-class structure.
\end{foldable}

\begin{foldable} Recent methods in BBDFKD~\cite{wang2021zero,zhang2022ideal,yuan2024data} use synthetic data, but they still encounter two primary bottlenecks.
  First, they suffer from a lack of domain coverage: synthetic samples fail to replicate the hierarchical structures and semantic variation of natural environments, resulting in a distribution mismatch with the teacher's expertise.
  Second, the distillation signals are suboptimal.
  By relying exclusively on hard top-1 indices, students are denied the \textit{dark} knowledge of inter-class relationships---a cornerstone of \ac{KD} essential for emulating expert logic.
\end{foldable}

\begin{foldable} \textbf{Our method.} We propose \textbf{\acl{DIPKD}} (\acs{DIPKD}) to tackle the BBDFKD problem, \ie, training a student from a black-box teacher with only top-1 signal and without access to any real data.
  To improve \ac{KD}, we incorporate \textit{diversity}, \ie, how broadly the samples distribute, into a collaborative three-phase pipeline:{
    (1)~\textit{Synthesis}: crafts synthetic \textit{image priors} with diverse \DIFdelbegin \DIFdel{patterns and semantics of }\DIFdelend \DIFaddbegin \DIFadd{structural patterns resembling }\DIFaddend natural objects;
    (2)~\textit{Contrast}: employs a \textit{primer student} to act as a white-box mediator and ensure synthetic samples are informationally distinct via contrastive learning;
    and (3)~\textit{Distillation}: distills the final student optimally with hard labels from teacher and soft probabilistic labels from the primer student.
  }
  Synthesis and Contrast together address the domain-coverage deficit; the primer student's soft labels in Distillation partially restore the information-bottleneck deficit.
\end{foldable}

\begin{foldable} \textbf{Contribution.} In summary, we highlight our contributions as follows:
  \begin{enumerate}
    \item We introduce \textit{image priors}, a novel class of synthetic images exhibiting diversity across hierarchical structures, nonlinearity, and \DIFdelbegin \DIFdel{semantics}\DIFdelend \DIFaddbegin \DIFadd{spatial structure}\DIFaddend .
    \item We propose the novel use of a \textit{primer student}, enabling contrastive optimization and \DIFdelbegin \DIFdel{extraction of soft probabilistic signals originally restricted in }\DIFdelend \DIFaddbegin \DIFadd{surrogate soft supervision where the }\DIFaddend black-box \DIFdelbegin %DIFDELCMD < \ac{KD}%%%
\DIFdelend \DIFaddbegin \DIFadd{teacher provides only hard labels}\DIFaddend .
    \item Extensive analysis of our framework \DIFdelbegin \DIFdel{demonstrates that it effectively expands the semantic }\DIFdelend \DIFaddbegin \DIFadd{shows that it expands the }\DIFaddend coverage of synthetic domains, leading to \DIFdelbegin \DIFdel{our }\DIFdelend state-of-the-art performance \DIFdelbegin \DIFdel{across }\DIFdelend \DIFaddbegin \DIFadd{on 11 of }\DIFaddend 12 benchmarks.
  \end{enumerate}
\end{foldable}

\begin{foldable}
  The remainder of this chapter develops the framework (\S\ref{sec:43-research02-framework}), reports extensive experiments and ablations (\S\ref{sec:44-research02-experiments}), and concludes with discussion (\S\ref{sec:45-research02-conclusion}).
\end{foldable}
 \section{Proposed Framework} \label{sec:43-research02-framework}
\begin{figure}[ht] \centering
  \includegraphics[width=0.99\textwidth]{./figures/40-research02/43-research02-framework-teaser.pdf}
  \caption[Illustration of our method DIPKD of three phases.]{
    Illustration of our method \textbf{DIPKD} of three phases.
    (\textbf{1})~\textbf{Synthesis}:{
      We craft image priors $\mathcal{X}_\mathcal{P}$ with three sub-components: hierarchical noise, nonlinear transformation, and semantic cutmixing.
      Synthetic images $x \sim \mathcal{X}_\mathcal{P}$ have diverse and distinct patterns.
    }
    (\textbf{2})~\textbf{Contrast}:{
      We construct the dataset $\mathcal{D} = \{x | x \sim \mathcal{X}_\mathcal{P}\}$ and train a primer student $S_0$ as a feature extractor.
      Then, we enforce an instance-discriminator $R$ to distinguish embeddings of images $x \sim \mathcal{D}$ to be similar to its positive view $x^{+}$, and dissimilar to its negative view $x^{-}\ne x$.
      Thus, $\mathcal{D}$ is optimized to be even more diverse.
    }
    (\textbf{3})~\textbf{Distillation}:{
      We distill the student $S$ by matching with both hard labels from teacher $T$ via $\mathcal{L}_\text{Hard-KD}$ and soft labels from primer student $S_0$ via $\mathcal{L}_\text{Soft-KD}$, providing extra knowledge.
    }
  }
  \label{fig:43-research02-framework-teaser}
\end{figure}

\subsection{Problem setup} \label{sec:43-research02-framework-setup}
\begin{foldable} \textbf{Problem statement.}
  Given solely a \textit{black-box} pre-trained teacher network $T$ that returns only top-1 hard labels (no internal features, logits, or gradients) and \textit{no real training data}, our goal is to train a student $S$ to approximate $T$.
\end{foldable}

\begin{foldable} \textbf{Naive solution.}
  Lacking a direct solution, one might adapt standard \ac{KD}~\cite{hinton2015distilling} as a naive baseline.
  Without real data, this necessitates using uniform noise images, \eg, $x \sim \mathcal{X}_\mathcal{U}$, where $\mathcal{X}_\mathcal{U} = \mathcal{U}[0,1]^{C\times H\times W}$.
  In this restricted interface, only the hard label $\hat{y}_{T}^\text{hard} = T(x) \in \{1, \ldots, K\}$ is available to guide the student's probabilistic outputs $\hat{y}_{S}^\text{soft} = S(x)$.
  The objective is defined as:
  \begin{equation} \label{eq:loss-naive-kd}
    \mathcal{L}_\text{NaiveKD} = \mathbb{E}_{x \sim \mathcal{X}_\mathcal{U}} \left[\mathcal{L}_\text{CE} \left(\hat{y}_{S}^\text{soft}, \hat{y}_{T}^\text{hard} \right)\right],
  \end{equation}
  where $\mathcal{L}_\text{CE}$ is the cross-entropy loss.
  While this can handle trivial cases, it uses semantically non-diverse noise, and fails to use inter-class knowledge---the primary advantage of \ac{KD}.
\end{foldable}

\begin{foldable} \textbf{Proposal.}
  To improve BBDFKD, we propose \textbf{\acl{DIPKD}} (\acs{DIPKD}), structured into three collaborative phases:
  \begin{enumerate}
    \item {
        In the \textbf{Synthesis} phase, we design a pipeline to create \textit{image priors} $\mathcal{X}_\mathcal{P}$ that simulate the hierarchical structures, nonlinearity, and \DIFdelbegin \DIFdel{meaningful semantics }\DIFdelend \DIFaddbegin \DIFadd{spatial structure }\DIFaddend in natural scenes and objects.
        This is achieved via three sub-components: hierarchical noise, nonlinear transformation, and semantic cutmixing.
      }
    \item {
        Then, in the \textbf{Contrast} phase, we sample a dataset $\mathcal{D} = \{x | x \sim \mathcal{X}_\mathcal{P}\}$.
        We further train a \textit{primer student} $S_0$ to optimize $\mathcal{D}$ via a contrastive learning objective to make every sample informationally \textit{distinguishable}.
      }
    \item {
        Finally, in the \textbf{Distillation} phase, we use the optimized dataset $\mathcal{D}$ to distill the student $S$.
        Crucially, we extract \textit{soft} knowledge from the primer student $S_0$ \DIFdelbegin \DIFdel{to restore }\DIFdelend \DIFaddbegin \DIFadd{as surrogate supervision for }\DIFaddend the inter-class signals absent in previous BBDFKD methodologies.
      }
  \end{enumerate}
\end{foldable}

\subsection{Synthesis: Reconstructing the knowledge domain} \label{sec:43-research02-framework-synthesis}
\DIFdelbegin %DIFDELCMD < \begin{figure}[ht]
%DIFDELCMD <   \centering
%DIFDELCMD <   \includegraphics[width=0.99\textwidth]{./figures/40-research02/43-research02-framework-synthesis.pdf}
%DIFDELCMD <   \caption[Overview of DIPKD Synthesis pipeline.]{
%DIFDELCMD <     Overview of DIPKD Synthesis pipeline.
%DIFDELCMD <     In the \textbf{Synthesis} phase, we generate image priors $\mathcal{X}_\mathcal{P}$ with 3 sub-components.
%DIFDELCMD <     (\textbf{1})~\textbf{Hierarchical noise}: We sample $x_\text{hn}$ from multi-scale noise images to capture hierarchical structures.
%DIFDELCMD <     (\textbf{2})~\textbf{Nonlinear transform}: We transform $x_\text{hn}$ to $x_\text{nt}$ with nonlinear filters to capture nonlinearity.
%DIFDELCMD <     and (\textbf{3})~\textbf{Semantic cutmixing}: We apply cutmixing on $x_\text{nt}$ with a semantic mask $\hat{m}$ to construct the final images $x$, simulating diverse semantics.
%DIFDELCMD <     Our image priors $\mathcal{X}_\mathcal{P}$ are visually much more diverse than noise images sampled from random Gaussian or random uniform distributions $\mathcal{X}_{\mathcal{N}}, \mathcal{X}_\mathcal{U}$.
%DIFDELCMD <   }
%DIFDELCMD <   \label{fig:43-research02-framework-viz-synthesis}
%DIFDELCMD < \end{figure}
%DIFDELCMD < %%%
\DIFdelend \DIFaddbegin \begin{figure}[ht]
  \centering
  \includegraphics[width=0.99\textwidth]{./figures/40-research02/43-research02-framework-synthesis.pdf}
  \caption[Overview of DIPKD Synthesis pipeline.]{
    Overview of DIPKD Synthesis pipeline.
    In the \textbf{Synthesis} phase, we generate image priors $\mathcal{X}_\mathcal{P}$ with 3 sub-components.
    (\textbf{1})~\textbf{Hierarchical noise}: We sample $x_\text{hn}$ from multi-scale noise images to capture hierarchical structures.
    (\textbf{2})~\textbf{Nonlinear transform}: We transform $x_\text{hn}$ to $x_\text{nt}$ with nonlinear filters to capture nonlinearity.
    and (\textbf{3})~\textbf{Semantic cutmixing}: We apply cutmixing on $x_\text{nt}$ with a semantic mask $\hat{m}$ to construct the final images $x$, simulating diverse spatial structures.
    Our image priors $\mathcal{X}_\mathcal{P}$ are visually much more diverse than noise images sampled from random Gaussian or random uniform distributions $\mathcal{X}_{\mathcal{N}}, \mathcal{X}_\mathcal{U}$.
  }
  \label{fig:43-research02-framework-viz-synthesis}
\end{figure}
\DIFaddend 

\begin{foldable} Natural images are characterized by strong local pixel structures and complex, nonlinear semantics~\cite{lecun1998gradient}.
  To capture these properties, we generate image priors $\mathcal{X}_\mathcal{P}$ through a pipeline of three stages: \textit{hierarchical noise}, \textit{nonlinear transformation}, and \textit{semantic cutmixing}, visualized in Fig.~\ref{fig:43-research02-framework-viz-synthesis}.
  This pipeline serves to capture the structural \DIFdelbegin \DIFdel{and semantic }\DIFdelend breadth of the expert's original knowledge domain through enhanced diversity.
\end{foldable}

\subsubsection{Hierarchical noise} \label{sec:43-research02-framework-synthesis-hiernoise}
\begin{foldable} Hierarchy is a universal trait in natural scenes and objects.
  Based on this observation, we design a sampling technique that captures image patterns at both local and global scales to mimic the hierarchical structures of visual signals.
\end{foldable}

\begin{foldable} In a standard pixel-wise noise image $x \sim X_\mathcal{U}$, patterns are primarily local and independent.
  Conversely, a monochromatic (single-color) image has dominant global correlation.
  To balance these patterns, we design a sampler that combines noise at multiple scales.
  For an image of dimension $C \times H \times W$, we sample noise images with increasing scales:
  \begin{equation} \label{eq:synthesis-hiernoise-multiscale}
    \left\{ x_{d} |x_{d} \sim \mathcal{U}[0,1]^{C\times2^{d}\times2^{d}} \right\}_{d=0}^{d_\text{max}},
  \end{equation}
  where $d_\text{max} = \text{ceil}(\log_{2} \max(H, W))$ is the first scale to sufficiently cover $H\times W$.
  The dimensions are derived deterministically, \eg, a $32 \times 32$ image is generated from scales $\{2^0, 2^1, \ldots, 2^5\}$.
  To combine local-to-global noise into a single image, we sample coefficients $\alpha_d$ from a normal distribution and apply softmax to get weights $\bar{\alpha}_{d}$.
  We then mix the noise images $\{x_d\}$ to obtain the \textit{hierarchical noise} image:
  \begin{equation} \label{eq:synthesis-hiernoise-combine}
    x_\text{hn} = \sum_{d=0}^{d_\text{max}} \bar{\alpha}_{d} \cdot f_\text{upscale}\left(x_{d};2^{d_\text{max}}\right),
  \end{equation}
  where $f_\text{upscale}$ upscales all noise images to the same dimension $2^{d_\text{max}}$ for proper mixing.
  This approach provides a multi-scale structural foundation that allows the student to learn complex visual hierarchies in the absence of real-world data.
\end{foldable}

\subsubsection{Nonlinear transformation} \label{sec:43-research02-framework-synthesis-nonlintrans}
\begin{foldable} To enrich nonlinear patterns of synthetic images, we apply random rotation $f_\text{rot}$ of $[-45^{\circ}, 45^{\circ}]$, random elastic transform $f_\text{elas}$, and random cropping $f_\text{crop}$ to the specified size $H \times W$.
  These operations are known to be nonlinear and effective in data augmentation~\cite{simard2003best}, and in our case, boosting data diversity.
  Accordingly, we create a \textit{nonlinear-transformed} image as:
  \begin{equation} \label{eq:synthesis-nonlintrans}
    x_\text{nt} = f_\text{crop}^{H\times W} \circ f_\text{elas} \circ f_\text{rot} (x_\text{hn}).
  \end{equation}
\end{foldable}

\subsubsection{Semantic cutmixing} \label{sec:43-research02-framework-synthesis-semcutmix}
\begin{foldable} We aim to augment synthetic images with the cohesive structures often found in natural objects.
  However, previous techniques like CutMix~\cite{yun2019cutmix} rely on simple rectangular occlusions that lack semantic meaning.
  To address this, we extend CutMix by improving the masking process to nonlinearly blend image pairs into \DIFdelbegin \textit{\DIFdel{semantically coherent}} %DIFAUXCMD
\DIFdel{structures.
}\DIFdelend \DIFaddbegin \DIFadd{irregular, spatially coherent structures, which we call }\textit{\DIFadd{semantic cutmix}}\DIFadd{. These masks are structural: they are not object-level segments, and a teacher label on a blended image does not make it a meaningful class instance.
}\DIFaddend \end{foldable}

\begin{foldable} We achieve these patterns through refinement of a semantic mask $\hat{m}$, initialized as $m_{0} \sim \mathcal{U}[0,1]^{1 \times H \times W}$.
  To balance spatial continuity with regional partitioning, we respectively use a blurring filter $f_\text{blur}$ and our \textit{diverging filter} $f_\text{divg}$.
  Specifically, $f_\text{divg}$ is a piecewise quadratic operator that is continuously differentiable at its inflection point (0.5), driving pixel values toward binary states (0 or 1), producing sharp, well-defined region boundaries.
  \begin{equation} \label{eq:synthesis-diverging-filter}
    f_\text{divg}(m) =
    \begin{cases}
      2 m^2           & , m \ge 0.5 \\
      1 - 2 (m - 1)^2 & , m < 0.5
    \end{cases}.
  \end{equation}
\end{foldable}

\begin{foldable} We obtain the final $\hat{m}$ by iteratively applying $m_{i+1} = f_\text{divg} \circ f_\text{blur}(m_i)$ until the mean absolute pixel-wise update reaches a steady-state threshold $\epsilon < 10^{-2}$, empirically around 10 iterations.
  This mask is used for pairwise cutmixing $\{x_\text{nt}\}$ against themselves or random monochromatic images $\mathcal{X}_\text{mc}$ to form our final \textit{image priors} $\mathcal{X}_\mathcal{P} = \{x\}$, where:
  \begin{equation} \label{eq:synthesis-semcutmix}
    x= \text{CutMix}(x_i, x_j; \hat{m}), \text{\quad s.t. \quad} x_i,x_j \sim \{x_\text{nt}\} \cup \mathcal{X}_\text{mc} \text{ and } x_i \ne x_j.
  \end{equation}
  This approach \DIFdelbegin \DIFdel{ensures }\DIFdelend \DIFaddbegin \DIFadd{encourages }\DIFaddend the synthetic domain of $\mathcal{X}_\mathcal{P}$ \DIFdelbegin \DIFdel{captures diverse semantics required for effective }\DIFdelend \DIFaddbegin \DIFadd{to capture diverse structural patterns useful for }\DIFaddend black-box knowledge acquisition.
\end{foldable}

\subsection{Contrast: Collaborative diversity optimization} \label{sec:43-research02-framework-contrast}
\begin{foldable} Standard contrastive learning methods require internal feature extraction, which is prohibited by the black-box nature of our teacher.
  We propose the novel use of a \textit{primer student} $S_{0}$ as a white-box mediator to bridge this informational gap.
  We first train $S_{0}$ on an image priors dataset $\mathcal{D} = \{x | x \sim \mathcal{X}_\mathcal{P}\}$ following the Hard-KD objective:
  \begin{equation} \label{eq:loss-hard-kd}
    \mathcal{L}_\text{Hard-KD} = \mathbb{E}_{x \sim \mathcal{D}} \left[\mathcal{L}_\text{CE} \left(\hat{y}_{S_0}^\text{soft}, \hat{y}_{T}^\text{hard} \right)\right],
  \end{equation}
  where $\hat{y}_{S_0}^\text{soft} = S_0(x)$ is its probability prediction, and $\hat{y}_{T}^\text{hard} = T(x)$ is the teacher's hard label.
  Equipped with this white-box proxy, we can access its internal representations for explicit diversity optimization.
\end{foldable}

\begin{foldable} Following~\cite{fang2021contrastive}, diversity correlates with how \textit{``distinguishable''} samples are, under an \textit{explicitly optimizable} objective.
  To maximize diversity, we use the primer student's backbone $S_{0}^\mathcal{H}$ to extract image embeddings and append an instance-discriminator $R$ (a shallow multi-layer perceptron) to recognize each embedding as a unique class.
  We argue that even if $S_0$ is yet to match the teacher's expertise, its backbone $S_{0}^\mathcal{H}$ is capable of capturing the underlying semantics required for contrastive learning.
  This is supported by prior works in contrastive learning~\cite{chen2020simple,he2020momentum,zbontar2021barlow}, which do not depend on a high-accuracy classifier, but often train from scratch.
  For any two images $(x_i, x_j)$, we compute their cosine similarity as:
  \begin{equation}
    \text{Sim}\left(x_i, x_j\right)=
    \frac{ \left\langle R \circ S_{0}^\mathcal{H}(x_i), R \circ S_{0}^\mathcal{H}(x_j)\right\rangle }{ \left\Vert R \circ S_{0}^\mathcal{H}(x_i) \right\Vert \cdot \left\Vert R \circ S_{0}^\mathcal{H}(x_j) \right\Vert }.
  \end{equation}
\end{foldable}

\begin{foldable} For diversity optimization, from an image $x \sim \mathcal{D}$, we generate positive views $x^{+}$ via augmentation and sample other images $x^{-} \sim \mathcal{D}, x^{-} \neq x$ as negative views.
  The contrastive objective maximizes the similarity between $(x, x^{+})$ and minimizes that between $(x, x^{-})$:
  \begin{equation} \label{eq:loss-contrastive}
    \mathcal{L}_\text{CT} = \mathbb{E}_{x \sim \mathcal{D}} \left[
      -\log \frac{ \exp\left( \text{Sim}(x, x^{+}) \right) }{ \sum_{x^{-} \ne x} \exp\left( \text{Sim}(x, x^{-}) \right) }
    \right].
  \end{equation}
  By backpropagating $\nabla\mathcal{L}_\text{CT}$ to update both the instance-discriminator $R$ and the samples $x \in \mathcal{D}$, we explicitly reinforce the semantic diversity of $\mathcal{D}$.
\end{foldable}

\subsection{Distillation: Restoring informative signals} \label{sec:43-research02-framework-distillation}
\begin{foldable}
  Once the synthetic dataset $\mathcal{D}$ has been optimized for diversity, we use it for the final training of the student $S$.
  Our objective is to \DIFdelbegin \DIFdel{restore }\DIFdelend \DIFaddbegin \DIFadd{approximate }\DIFaddend the high-bandwidth signals typical of standard \ac{KD}~\cite{hinton2015distilling}, \ie, incorporating both \textit{hard} and \textit{soft} knowledge, which were \textit{limited or entirely absent} in recent works~\cite{wang2021zero, zhang2022ideal, yuan2024data}.
  Having privileged access to the primer network $S_0$, we form \DIFdelbegin \DIFdel{an optimal }\DIFdelend \DIFaddbegin \DIFadd{a }\DIFaddend distillation objective by jointly matching the student's logits $S^\mathcal{Z}(x)$ with soft labels $S_{0}^\mathcal{Z}(x)$ and hard labels $\hat{y}_{T}^\text{hard}$:
  \begin{equation} \label{eq:loss-kd-total}
    \mathcal{L}_{S}
    = \mathcal{L}_\text{Hard-KD} + \mathcal{L}_\text{Soft-KD} \\
    = \mathbb{E}_{x \sim \mathcal{D}}\Big[
      \mathcal{L}_\text{CE}\left( \hat{y}_{S}^\text{soft}, \hat{y}_{T}^\text{hard} \right) +
      \mathcal{L}_\text{L1}\left( S^\mathcal{Z}(x), S_{0}^\mathcal{Z}(x) \right)
    \Big],
  \end{equation}
  where $\mathcal{L}_\text{L1}$ is the L1 divergence loss, applied directly to logits to preserve raw inter-class relationships.
  By reconciling the teacher's authoritative expertise with the primer student's relational hints, the student $S$ overcomes the signal scarcity of the black-box interface.
  This collaborative training allows for a deeper understanding of the teacher's decision-making logic, resulting in superior performance across restricted \ac{AI} ecosystems.
\end{foldable}
 \section{Experiments} \label{sec:44-research02-experiments}
\begin{foldable} We evaluate \ac{DIPKD} on extensive BBDFKD benchmarks, supplemented by ablation studies to validate its mechanics and practical effectiveness.
\end{foldable}

\subsection{Experimental setup} \label{sec:44-research02-experiments-setup}
\begin{foldable} We evaluate our method across 12 benchmarks, including 8 general-purpose datasets (USPS, MNIST, SVHN, FMNIST, CIFAR10, CIFAR100, Tiny-ImageNet, Imagenette)~\cite{lecun1998gradient,hull1994database,netzer2011reading,xiao2017fashion,krizhevsky2009learning,le2015tiny,howard2020imagenette} and 4 subsets from the medical dataset MedMNIST~\cite{yang2021medmnist}.
  We investigate three teacher-student architectures: LeNet5~\cite{lecun1998gradient}, AlexNet~\cite{krizhevsky2012imagenet}, and ResNet~\cite{he2016deep}.
  These datasets and architectures are the de facto evaluation in BBDFKD methods~\cite{wang2021zero, zhang2022ideal, yuan2024data}.
\end{foldable}

\subsection{Baselines} \label{sec:44-research02-experiments-baselines}
\begin{foldable} We compare our method \ac{DIPKD} with the baselines:
  \begin{itemize}
    \item \textbf{NaiveKD}: The student is trained on uniform noise $x \sim \mathcal{X}_\mathcal{U}$, where $\mathcal{X}_\mathcal{U}=\mathcal{U}[0,1]^{C \times H \times W}$ and the teacher's hard labels $\hat{y}_{T}^\text{hard} = T(x) \in \{1, \ldots, K\}$.
    \item \textbf{BBDFKD methods}: We reproduce three state-of-the-art methods: ZSDB3~\cite{wang2021zero}, IDEAL~\cite{zhang2022ideal}, and DFHL-RS~\cite{yuan2024data}.
  \end{itemize}
  For fair comparison, we use the same teacher-student network architectures and synthetic image budget for all methods.
  We report the \textit{mean} $\pm$ \textit{standard error} (SE) of distillation accuracy, which is the student accuracy on the hold-out test set, across five runs.
  The accuracies of baseline BBDFKD methods are reproduced with their official implementations\footnote{ZSDB3 at \url{https://github.com/zwang84/zsdb3kd}}\,\footnote{IDEAL at \url{https://github.com/SonyResearch/IDEAL}}\,\footnote{DFHL-RS at \url{https://github.com/LetheSec/DFHL-RS-Attack}}.
  We ensured that all baseline implementations were tuned ethically and fairly following the recommended settings.
\end{foldable}

\subsection{Standard experiments} \label{sec:44-research02-experiments-standard}
\begin{foldable} We categorize the benchmarks as \textit{simple}, \textit{complex}, and \textit{domain-specific} datasets.
  Four simple datasets---MNIST, USPS, SVHN, FMNIST---all have 10 classes of simple objects, at a $32\times32$ or smaller resolution.
  The complex datasets have more complicated objects or more classes: CIFAR10 and CIFAR100 (10 and 100 classes, $32\times32$), Tiny-ImageNet (200 classes, $64\times64$), and Imagenette (10 classes, but high-resolution).
  Finally, we evaluate the domain-specific performance on 4 subsets of MedMNIST, including BreastMNIST, OrganAMNIST, DermaMNIST, and BloodMNIST, which include $32\times32$ images of varying medical modalities.
  We report the distillation accuracy of all methods in Table~\ref{tab:standard-all}, where our method \textbf{\ac{DIPKD}} consistently outperforms all baselines across 11/12 datasets, with significant margins on complex and domain-specific ones.
\end{foldable}

\begin{table}[ht]
  \centering
  \begin{threeparttable}
    \caption[DIPKD accuracy on simple, complex, and domain-specific datasets]{
      DIPKD distillation accuracy (\%) on simple, complex, and domain-specific datasets.
    }
    \label{tab:standard-all}
    \small
    \begin{tabular}{l c c c c c c c}
      \toprule
      \boldtight{Dataset} & \boldtight{Setup} & \boldtight{Teacher} & \boldtight{NaiveKD} & \boldtight{ZSDB3} & \boldtight{IDEAL} & \boldtight{DFHL-RS} & \boldtight{DIPKD} \\
      \midrule

      \multicolumn{8}{c}{\textbf{Simple datasets}} \\
      \midrule
      MNIST    & L, 20K  & 99.28 & 96.54$_{\pm 0.60}$           & 96.66$_{\pm 0.39}$           & 96.96$_{\pm 0.34}$ & 98.53$_{\pm 0.10}$          & \textbf{98.59}$_{\pm 0.08}$ \\
      USPS     & L, 50K  & 95.47 & 42.05$_{\pm 0.08}$           & 65.50$_{\pm 0.16}$           & 92.66$_{\pm 0.57}$ & 93.97$_{\pm 0.60}$          & \textbf{94.20}$_{\pm 0.21}$ \\
      SVHN     & A, 50K  & 96.16 & 83.35$_{\pm 0.57}$           & 85.38$_{\pm 1.69}$           & 87.86$_{\pm 0.93}$ & 95.41$_{\pm 0.13}$          & \textbf{95.94}$_{\pm 0.05}$ \\
      FMNIST   & L, 50K  & 90.90 & 55.11$_{\pm 1.23}$           & 71.60$_{\pm 2.35}$           & 82.84$_{\pm 1.22}$ & 83.89$_{\pm 1.35}$          & \textbf{83.98}$_{\pm 0.65}$ \\
      \midrule

      \multicolumn{8}{c}{\textbf{Complex datasets}} \\
      \midrule
      CIFAR10  & R, 50K  & 95.21 & 13.99$_{\pm 0.23}$           & 56.74$_{\pm 1.98}$           & 67.58$_{\pm 1.29}$ & 76.50$_{\pm 0.84}$          & \textbf{83.41}$_{\pm 0.46}$ \\
      CIFAR100 & R, 100K & 78.36 & \phantom{0}2.68$_{\pm 0.13}$ & \phantom{0}9.19$_{\pm 0.39}$ & 33.64$_{\pm 0.91}$ & 51.97$_{\pm 0.11}$          & \textbf{52.85}$_{\pm 0.30}$ \\
      TinyImgN & R, 500K & 66.08 & \phantom{0}0.75$_{\pm 0.07}$ & \phantom{0}3.29$_{\pm 0.16}$ & 23.76$_{\pm 1.03}$ & \textbf{30.52}$_{\pm 1.43}$ & 28.03$_{\pm 0.22}$ \\
      Imagenette & R, 50K  & 91.29 & 33.21$_{\pm 0.39}$         & 33.43$_{\pm 0.85}$           & 36.90$_{\pm 0.38}$ & 40.53$_{\pm 2.36}$          & \textbf{61.84}$_{\pm 0.77}$ \\
      \midrule

      \multicolumn{8}{c}{\textbf{Domain-specific datasets}} \\
      \midrule
      BreastM  & L, 50K  & 79.49 & 61.78$_{\pm 1.60}$           & 63.37$_{\pm 1.48}$           & 67.54$_{\pm 1.27}$ & 73.81$_{\pm 1.80}$          & \textbf{77.95}$_{\pm 0.77}$ \\
      OrganAM  & L, 50K  & 87.60 & 24.09$_{\pm 0.44}$           & 51.81$_{\pm 1.21}$           & 69.09$_{\pm 1.92}$ & 75.12$_{\pm 2.89}$          & \textbf{82.32}$_{\pm 1.08}$ \\
      DermaM   & R, 50K  & 79.60 & 56.82$_{\pm 1.50}$           & 60.15$_{\pm 2.93}$           & 62.76$_{\pm 2.76}$ & 64.71$_{\pm 3.02}$          & \textbf{71.29}$_{\pm 2.67}$ \\
      BloodM   & R, 50K  & 98.19 & \phantom{0}9.09$_{\pm 0.13}$ & 39.62$_{\pm 1.67}$           & 55.28$_{\pm 2.18}$ & 61.67$_{\pm 3.11}$          & \textbf{73.98}$_{\pm 1.77}$ \\
      \bottomrule
    \end{tabular}
    \begin{tablenotes}[flushleft] \footnotesize
    \item {Setup denotes architecture \{\textbf{\underline{L}}eNet5, \textbf{\underline{A}}lexNet, \textbf{\underline{R}}esNet18\} and synthetic dataset size; [\textbf{bold}] best results.}
    \end{tablenotes}
  \end{threeparttable}
\end{table}

\begin{foldable}
  While performance saturates across all baselines on simple datasets, the accuracy gap widens as domain complexity scales.
  NaiveKD is first to collapse to near random-guessing levels in complex datasets, implying that unstructured noise cannot elicit the deep semantic knowledge embedded in black-box experts.
  Notably, on high-resolution images like Imagenette, our \ac{DIPKD} achieves a +21.31\% absolute gain over the strongest baseline, DFHL-RS.
  On domain-specific medical tasks, \ac{DIPKD} outperforms baselines by at least +4\% to +12\%.
  We also notice that abundant classes represent an interesting boundary condition, where \ac{DIPKD} remains highly competitive on CIFAR100 and approaches state-of-the-art on Tiny-ImageNet.
  These results suggest that data diversity is a critical factor for effective knowledge acquisition, particularly as task complexity and domain specificity increase.
\end{foldable}

\begin{foldable}
  In terms of robustness, we observe that \ac{DIPKD} has consistently low standard errors on most datasets, indicating strong stability across runs.
  We attribute this robustness to our integrated focus on diversity and signal restoration, which provide the student with softer training signals than noise-heavy or hard-label baselines.
\end{foldable}

\subsection{Ablation studies} \label{sec:44-research02-experiments-abla}
\begin{foldable}
  In this section, we provide deeper insights into our framework through component dissection and practical application.
  We first investigate the contribution of individual components, the effect of synthetic data budget, and the semantics of our image priors.
  Then, we further extend our evaluation to real-world scenarios, including cross-architecture settings and aggressive model compression.
\end{foldable}

\subsubsection{Contribution of components} \label{sec:44-research02-experiments-abla-components}

\begin{table}[ht]
  \centering
  \begin{threeparttable}
    \caption[DIPKD ablation study: incremental components.]{
      DIPKD ablation study: distillation accuracy (\%) with incremental components.
    }
    \label{tab:abla-components}
    \begin{tabular}{l l l}
      \toprule
      \makecell{\textbf{Method}} & \makecell{\textbf{MNIST}} & \makecell{\textbf{CIFAR10}} \\
      \midrule
      NaiveKD (noise $x \sim \mathcal{X}_\mathcal{U}$, Hard-KD) & \phantom{$+$ }\meanstd{96.02}{.60} & \phantom{$+$ }\meanstd{13.99}{.23} \\
      $+$ \textbf{Synthesis:}                                   & $+$ \phantom{0}2.27                & $+$ 64.18 \\
      \phantom{12} \bltrg Hierarchical noise                    & \phantom{$+$ 0}1.79\bltrg          & \phantom{$+$} 59.26\bltrg  \\
      \phantom{12} \bltrg Nonlinear transform                   & \phantom{$+$ 0}0.18\bltrg          & \phantom{$+$ 0}2.38\bltrg  \\
      \phantom{12} \bltrg Semantic cutmixing                    & \phantom{$+$ 0}0.30\bltrg          & \phantom{$+$ 0}2.54\bltrg  \\
      $+$ \textbf{Contrast:} Update $\mathcal{D}$ via $\mathcal{L}_\text{CT}$ & $+$ \phantom{0}0.16  & $+$ \phantom{0}3.14 \\
      $+$ \textbf{Distillation:} Add Soft-KD to $\mathcal{L}_S$ & $+$ \phantom{0}0.14                & $+$ \phantom{0}2.10 \\
      \midrule
      Complete \textbf{DIPKD}                                  & \phantom{$+$ }98.59$_{\pm .08}$    & \phantom{$+$ }83.41$_{\pm .46}$ \\
      \bottomrule
    \end{tabular}
    \begin{tablenotes}[flushleft] \footnotesize
    \item {[$\blacktriangle$] denotes individual contribution of each Synthesis sub-component.}
    \end{tablenotes}
  \end{threeparttable}
\end{table}

\begin{foldable}
  We evaluate each component on MNIST and CIFAR10 to represent simple and complex benchmarks.
  Starting from NaiveKD, we sequentially add components and report student accuracy in Table~\ref{tab:abla-components}.
  Synthesis provides the largest gain ($+$2.27\% on MNIST, $+$64.18\% on CIFAR10), primarily driven by hierarchical noise.
  This is equivalent to the primer student $S_0$ trained on $\mathcal{X}_\mathcal{P}$, achieving 98\% and 78\% on MNIST and CIFAR10 (which already outperforms the baselines from Sec. \ref{sec:44-research02-experiments-standard}).
  This confirms that $S_0$ is sufficiently competent to perform contrastive optimization and soft distillation.
  For instance, on CIFAR10, we achieve $+$3.14\% via Contrast and $+$2.10\% via Distillation.
  These results confirm that the Synthesis pipeline plays a vital role in generating diverse data, while the novel use of the primer student for Contrast and Distillation phases significantly boosts \ac{KD} performance beyond prior works.
  \DIFaddbegin \DIFadd{In our component ablation study, the Contrast and Distillation phases add little on simple datasets but clearly help on complex ones, so the extra training beyond reusing the primer as the final student is most justified there.
}\DIFaddend \end{foldable}

\subsubsection{Embeddings analysis --- Synthesis} \label{sec:44-research02-experiments-abla-embeddings-generation}
\begin{figure}[ht] \centering
  \includegraphics[trim={120 0 0 0},clip,width=0.90\columnwidth]{./figures/40-research02/44-research02-experiments-embeddings-objects.png}
  \caption[DIPKD feature embeddings of complex datasets, noise, and image priors.]{
    DIPKD feature embeddings of complex datasets, noise ($\mathcal{X}_\mathcal{U}$, $\mathcal{X}_\mathcal{N}$), and image priors $\mathcal{X}_\mathcal{P}$.
  }
  \label{fig:embeddings-universal}
\end{figure}

\begin{foldable}
  To assess the diversity and semantics of our image priors, we sample synthetic images and extract embeddings from a ResNet18 teacher backbone pre-trained on CIFAR10.\footnote{This is the 95.21\% ResNet18 teacher on CIFAR10, being reused as a held-out ``oracle''.}
  These 512-dimensional features are reduced to two dimensions using UMAP~\cite{mcinnes2018umap} with default configuration.
  We explore the remaining complex datasets CIFAR100, Tiny-ImageNet, and Imagenette, as well as uniform noise $\mathcal{X}_\mathcal{U}$, Gaussian noise $\mathcal{X}_\mathcal{N}$, and our image priors $\mathcal{X}_\mathcal{P}$.
  We declare that the use of real data and the teacher's backbone here is solely for evaluation purposes.
  We visualize these embeddings in Fig.~\ref{fig:embeddings-universal}.
\end{foldable}

\begin{table}[ht]
  \centering
  \begin{threeparttable}
    \caption[DIPKD distribution metrics on complex datasets.]{
      DIPKD distribution metrics (\%) on complex datasets.
    }
    \label{tab:abla-embeddings-coverage-objects}
    \small
    \begin{tabular}{c c c c c c}
      \toprule
      \textbf{Metric} & \textbf{Source} & \textbf{CIFAR10} & \textbf{CIFAR100} & \textbf{Tiny-ImageNet} & \textbf{Imagenette} \\
      \midrule
      & $\mathcal{X}_\mathcal{U}$ & \phantom{0}20.00 &           100.65 &           108.65 &           101.20 \\
      \textbf{Density}   & $\mathcal{X}_\mathcal{N}$ & \phantom{0}20.00 &           107.98 & \phantom{0}90.87 &           103.64 \\
      & $\mathcal{X}_\mathcal{P}$ & \phantom{0}49.14 & \phantom{0}96.20 & \phantom{0}96.00 & \phantom{0}92.14 \\
      \midrule
      & $\mathcal{X}_\mathcal{U}$ & \phantom{000}.00 & \phantom{000}.10 & \phantom{000}.06 & \phantom{000}.05 \\
      \textbf{Coverage}  & $\mathcal{X}_\mathcal{N}$ & \phantom{000}.00 & \phantom{000}.10 & \phantom{000}.06 & \phantom{000}.05 \\
      & $\mathcal{X}_\mathcal{P}$ & \phantom{00}7.33 & \phantom{0}60.79 & \phantom{0}63.88 & \phantom{0}78.48 \\
      \midrule
      & $\mathcal{X}_\mathcal{U}$ &           100.00 &           100.10 &           100.00 &           100.00 \\
      \textbf{Precision} & $\mathcal{X}_\mathcal{N}$ &           100.00 &           100.10 &           100.00 &           100.00 \\
      & $\mathcal{X}_\mathcal{P}$ & \phantom{0}80.83 & \phantom{0}97.76 & \phantom{0}99.29 & \phantom{0}99.10 \\
      \midrule
      & $\mathcal{X}_\mathcal{U}$ & \phantom{000}.00 & \phantom{000}.09 & \phantom{000}.06 & \phantom{000}.04 \\
      \textbf{Recall}    & $\mathcal{X}_\mathcal{N}$ & \phantom{000}.00 & \phantom{000}.10 & \phantom{000}.06 & \phantom{000}.05 \\
      & $\mathcal{X}_\mathcal{P}$ & \phantom{0}66.92 & \phantom{0}96.75 & \phantom{0}96.86 & \phantom{0}95.07 \\
      \bottomrule
    \end{tabular}
  \end{threeparttable}
\end{table}

\begin{foldable}
  As shown, $\mathcal{X}_\mathcal{U}$ and $\mathcal{X}_\mathcal{N}$ form a compact, separate cluster far from real images, indicating that their semantics are limited and dissimilar to natural objects.
  In contrast, $\mathcal{X}_\mathcal{P}$ overlaps with real image regions, suggesting that they share similar semantics.
  For quantitative analysis, we evaluate four distribution metrics (precision, recall, density, coverage) from~\cite{naeem2020reliable}.
  The statistics in Table~\ref{tab:abla-embeddings-coverage-objects} show that our image priors $\mathcal{X}_\mathcal{P}$ achieve significantly higher density, coverage, and recall at a small trade-off in precision, compared to noise sources $\mathcal{X}_\mathcal{U}$ and $\mathcal{X}_\mathcal{N}$.
  These results further confirm that $\mathcal{X}_\mathcal{P}$ has superior fidelity and diversity, and explain why $\mathcal{X}_\mathcal{P}$ excels in BBDFKD.
\end{foldable}

\subsubsection{Embeddings analysis --- Contrast} \label{sec:44-research02-experiments-abla-embeddings-contrast}
\begin{foldable}
  To analyze the Contrast phase's impact on diversity, we evaluate 50K image prior embeddings in Fig.~\ref{fig:embeddings-cifar10}.
  The backbone embeds real images into ten clusters (equivalent CIFAR10 classes; gray).
  While pre-Contrast priors $\mathcal{X}_\mathcal{P}$ (green) already scatter around real images, the Contrast phase (orange) expands coverage into neighboring regions.
  This increased semantic diversity lets the student capture broader teacher knowledge.
  The distribution metrics in Table~\ref{tab:abla-embeddings-coverage-cifar10} confirm this with a significant increase in coverage and recall, with a trivial trade-off in density.
  Overall, these results show that through fulfilling the contrastive objective, the Contrast phase improves data diversity semantically and empirically.
\end{foldable}

\begin{figure}[ht] \centering
  \includegraphics[trim={120 0 0 0},clip,width=0.90\columnwidth]{./figures/40-research02/44-research02-experiments-embeddings-cifar10.png}
  \caption[DIPKD feature embeddings of image priors]{
    DIPKD feature embeddings of image priors $\mathcal{X}_\mathcal{P}$ before and after the Contrast phase.
  }
  \label{fig:embeddings-cifar10}
\end{figure}

\begin{table}[ht]
  \centering
  \begin{threeparttable}
    \caption[DIPKD distribution metrics during the Contrast phase]{
      DIPKD distribution metrics (\%) during the Contrast phase.
    }
    \label{tab:abla-embeddings-coverage-cifar10}
    \begin{tabular}{l c c c c}
      \toprule
      & \textbf{Density} & \textbf{Coverage} & \textbf{Precision} & \textbf{Recall} \\
      \midrule
      Before Contrast & 70.46            & \phantom{0}6.13   & 91.69              & 82.18 \\
      After Contrast  & 66.93            &           14.66   & 91.68              & 94.00 \\
      \bottomrule
    \end{tabular}
  \end{threeparttable}
\end{table}

\subsubsection{Increasing synthetic dataset size} \label{sec:44-research02-experiments-abla-nsamples}
\begin{figure}[H] \centering
  \includegraphics[width=0.80\columnwidth]{./figures/40-research02/44-research02-experiments-abla-nsamples.pdf}
  \caption[DIPKD accuracy by image priors budget]{
    DIPKD distillation accuracy (\%) by image priors budget $N$.
  }
  \label{fig:abla-nsamples}
\end{figure}

\begin{foldable}
  The relationship between synthetic data budget $N$ and \ac{KD} performance mirrors the emergence of collective intelligence.
  While it is expected that \ac{KD} performance improves with large $N$ (Fig.~\ref{fig:abla-nsamples}), we are interested in the optimal cut-off point, which indicates the \textit{minimum quorum} required for the student to replicate the teacher's expertise.
  This threshold scales directly with the structural complexity of the environment.
  For digits datasets (MNIST, USPS, SVHN), minimal data ($N$ = 10K) is required to reach good results, implying informational redundancy.
  In large-scale settings (CIFAR100, Tiny-ImageNet, Imagenette), the student continues to improve around $N$ = 100K, suggesting that the information density has not yet saturated and the collective wisdom of the teacher's priors remains beneficial.
\end{foldable}

\subsubsection{Cross-architecture KD} \label{sec:44-research02-experiments-abla-crossarch}
\begin{table}[ht]
  \centering
  \begin{threeparttable}
    \caption[DIPKD cross-architecture accuracy]{
      DIPKD cross-architecture distillation accuracy (\%) on FMNIST, SVHN, and CIFAR10.
    }
    \label{tab:abla-crossarch}
    \begin{tabular}{l l l l}
      \toprule
      & \textbf{FMNIST} & \textbf{SVHN} & \textbf{CIFAR10} \\
      \midrule
      Teacher  & LeNet5 & AlexNet & ResNet18 \\
      & 90.90 & 96.16 & 95.21 \\
      \midrule
      LeNet5   & \meanstdbf{83.98}{0.65} & \meanstd{64.44}{1.83} & \meanstd{32.01}{3.18} \\
      AlexNet  & \meanstd{78.98}{0.92} & \meanstdbf{95.94}{0.05} & \meanstd{61.30}{0.69} \\
      ResNet18 & \meanstd{81.69}{1.58} & \meanstd{95.91}{0.04} & \meanstdbf{83.41}{0.46} \\
      \bottomrule
    \end{tabular}
    \begin{tablenotes}[flushleft] \footnotesize
    \item {[\textbf{bold}] best results.}
    \end{tablenotes}
  \end{threeparttable}
\end{table}

\begin{foldable}
  An under-investigated setting in BBDFKD methods is when the teacher and student have different architectures.
  We evaluate \ac{KD} with heterogeneous student architectures (LeNet5, AlexNet, ResNet18) on FMNIST, SVHN, and CIFAR10.
  Table~\ref{tab:abla-crossarch} shows that students perform best when matching the teacher's architecture (diagonal), with performance degrading as the capacity gap widens as expected.
  Thus, DIPKD is practical in real-world scenarios with unknown teachers, provided the student has sufficient capacity.
\end{foldable}

\subsubsection{Students of compressed capacity} \label{sec:44-research02-experiments-abla-compress}
\begin{figure}[ht] \centering
  \includegraphics[width=0.80\columnwidth]{./figures/40-research02/44-research02-experiments-abla-compression.pdf}
  \caption[DIPKD accuracy by compression ratio]{
    DIPKD distillation accuracy (\%) by compression ratio (CR).
  }
  \label{fig:abla-compression}
\end{figure}

\begin{foldable}
  Beyond same-size configurations, we evaluate DIPKD under aggressive compression, a scenario unaddressed in prior methods~\cite{wang2021zero,zhang2022ideal,yuan2024data}.
  Specifically, we distill student networks with 25\%, 4\%, and 1\% compression ratio (CR) by reducing channel counts by 2$\times$, 5$\times$, 10$\times$, respectively.
  Results indicate strong empirical robustness: performance remains stable at 25\% CR, with only tolerable decline at 4\% CR.
  Significant degradation occurs only at 1\% CR, where the student's extreme sparsity prevents it from capturing the complex inter-class signals.
  These results demonstrate that DIPKD is ideal for \textit{resource-constrained} mobile or edge environments.
\end{foldable}
 \section{Conclusion} \label{sec:45-research02-conclusion}

\begin{foldable}
  This work addresses information flow bottlenecks in decentralized \ac{AI} systems, where privacy and proprietary barriers create a black-box data-free environment.
  We introduce \textbf{DIPKD}, a framework framing \ac{KD} as the emergence of intelligence from a collective of synthetic image priors.
  By coordinating a three-phase pipeline: Synthesis, Contrast, and Distillation, we enable students to emulate expert logic through highly restricted interfaces.
  We achieve state-of-the-art performance \DIFdelbegin \DIFdel{across }\DIFdelend \DIFaddbegin \DIFadd{on 11 of }\DIFaddend 12 benchmarks, with notable improvements in complex and domain-specific settings.
  Beyond establishing a new empirical baseline, this study identifies data diversity as a key factor for effective knowledge acquisition in restricted environments.
  We hope our work can provide a robust foundation for future research into resilient and decentralized collective intelligence.
\end{foldable}
 \clearpage{}
\clearpage{}\acresetall

\chapter[Trajectory-Aware Knowledge Estimation for Text Dataset Distillation]{Trajectory-Aware Knowledge Estimation \\ for Text Dataset Distillation} \label{ch:50-research03}

\begin{foldable}
  The background on text \ac{DD} (\S\ref{sec:26-background-datasetdistill-textdd}) lays out what a complete method has to do, and \S\ref{sec:27-background-influence} and \S\ref{sec:28-background-ot} pin down the two parts no prior method has solved together: at extreme compression, importance scores read off a single checkpoint are systematically pulled toward noisy examples, and even with corrected scores, independent ranking cannot guarantee that the selected subset covers the weighted distribution.
  This chapter develops \textbf{\acl{TAKE}} (\textbf{\acs{TAKE}}), which scores each sample by integrating its influence over the entire training trajectory rather than at a single point, and then uses those scores as the source measure in a discrete Optimal Transport selection over an \ac{LLM}-generated candidate pool.
  The trajectory-aware score corrects the bias that the single-checkpoint baseline introduces, and the \ac{OT} formulation turns sample-level scoring into a corpus-level selection objective whose output remains human-readable.
\end{foldable}

\section{Introduction} \label{sec:51-research03-introduction}

\begin{foldable} The systemic costs that motivate dataset compression sharpen at the corpus scale: at extreme compression, the residual budget has to be spent on the samples that contribute the most to learning.
  Yet, adapting image \ac{DD} techniques~\cite{wang2018dataset,zhao2021dataset,zhao2023dataset} to text faces two concrete obstacles: non-differentiable token decoding blocks gradient-based synthesis, and embedding spaces are tightly coupled to specific architectures.
\end{foldable}

\begin{foldable} Existing text \ac{DD} methods~\cite{sucholutsky2021soft,li2021data,maekawa2023dataset,maekawa2024dilm,tao2024textual} fall short in at least one of the following: at extreme compression, uniform weighting exhausts the budget on uninformative samples; their objectives lack task alignment or global optimization; and their outputs are embeddings, neither auditable nor transferable.
  We survey this landscape in \S\ref{sec:26-background-datasetdistill-textdd} and show that no prior method simultaneously addresses all of the above.
  These gaps motivate three desiderata for TAKE:{
    (1)~\textit{knowledge-based weighting} --- weighting samples by their task-aligned downstream contribution, directing the budget toward the most informative ones;
    (2)~a \textit{dataset-level distillation objective} --- optimizing over the full training distribution rather than batch-by-batch, ensuring global coverage;
    and (3)~\textit{human-readable output} --- a distilled corpus that is directly inspectable, auditable, and transferable across architectures.
  }
\end{foldable}

\begin{table}[ht]
  \centering
  \begin{threeparttable}
    \caption{Landscape of text dataset distillation methods.}
    \label{tab:landscape-textdd-methods}
    \setlength{\tabcolsep}{8pt}
    \begin{tabular}{l c c c c c c}
      \toprule
      \textbf{Gap}         & SLDD         & DDTC         & DDAL         & DiLM         & DaLLME       & \textbf{TAKE} \\
      \midrule
      Task alignment       & $\checkmark$ & $\checkmark$ & $\triangle$  & $\triangle$  & $-$          & $\checkmark$ \\
      Global optimization  & $\triangle$  & $\triangle$  & $\triangle$  & $\triangle$  & $\triangle$  & $\checkmark$ \\
      Interpretability     & $\triangle$  & $-$          & $-$          & $\checkmark$ & $\checkmark$ & $\checkmark$ \\
      Importance weighting & $-$          & $-$          & $-$          & $-$          & $-$          & $\checkmark$ \\
      \bottomrule
      \multicolumn{7}{l}{$-$: absent, $\triangle$: partial, $\checkmark$: addressed}
    \end{tabular}
  \end{threeparttable}
\end{table}

\begin{foldable} The \textit{importance weighting} row in Table~\ref{tab:landscape-textdd-methods} stands out: every prior method leaves it open.
  The natural fix is influence functions~\cite{koh2017understanding}, which quantify each sample's counterfactual effect on the downstream loss.
  Scoring at a single checkpoint, however, introduces the \textit{\acl{HSB}} (\ac{HSB}): at convergence, clean and moderate samples have near-zero gradient by definition, so noisy or hard samples dominate the influence scores, exactly inverting the ranking a robust distilled corpus needs.
  In standard supervised training the bias can often be absorbed by regularization, but in the extreme low-data regime of \ac{DD} it amplifies instead: with only a handful of examples per class to allocate, biased selection actively suppresses the informative samples the surrogate depends on.
  To our knowledge, no prior text \ac{DD} method has identified or corrected for the \ac{HSB}, which is why the table's last row remains uniformly empty.
\end{foldable}

\begin{foldable} We propose \textbf{Trajectory-Aware Knowledge Estimation} (\textbf{TAKE}), the first text \ac{DD} method to satisfy all three desiderata.
  TAKE corrects the \ac{HSB} by integrating influence scores along the full training trajectory into a per-sample knowledge score that captures clean and moderate samples while down-weighting noisy ones (\S\ref{sec:52-research03-theory}).
  Then, TAKE fine-tunes a \ac{LLM} to generate a pool of human-readable candidate instances.
  The knowledge scores and synthetic candidates jointly feed a discrete \ac{OT} objective that globally selects the distilled corpus.
  Although \ac{OT} has been widely applied in \ac{NLP} and machine learning, to our knowledge it has not been applied to text \ac{DD}, nor combined with a knowledge-reweighted source distribution in any \ac{DD} setting; TAKE fills both gaps.
  Concretely, we contribute:
  \begin{itemize}
    \item {
        \textbf{Theory:} We formalize the \ac{HSB} in the text \ac{DD} setting, prove that single-checkpoint influence scores are biased toward noisy samples, and derive a formal gap bound for knowledge-reweighted distribution matching.
      }
    \item {
        \textbf{Method:} TAKE is the first text \ac{DD} method to jointly address sample weighting and distributional coverage, via trajectory-integrated knowledge scores and a discrete \ac{OT} objective aligned with the \ac{DD} goal.
      }
    \item {
        \textbf{Empirical:} TAKE matches or exceeds prior text \ac{DD} methods across six language benchmarks at extreme compression (0.1\% or 20 instances/class), producing human-readable distilled corpora that generalize across diverse backbone families.
      }
  \end{itemize}
\end{foldable}

\begin{foldable} The chapter follows the order in which the three desiderata are settled.
  \S\ref{sec:52-research03-theory} casts text \ac{DD} as reweighted distribution matching problem and shows where the naive instantiation fails.
  \S\ref{sec:53-research03-method} then builds the TAKE pipeline: gradient-based trajectory influence and its knowledge score, synthetic candidate generation, and distributional prototype selection via discrete \ac{OT}.
  \S\ref{sec:54-research03-experiments} evaluates TAKE on classification and natural language inference benchmarks at extreme compression and confirms that each component is load-bearing.
  \S\ref{sec:55-research03-limitations} discusses limitations and ethical considerations, and \S\ref{sec:56-research03-conclusion} closes the chapter.
\end{foldable}
 \section{Theoretical Framework} \label{sec:52-research03-theory}

\subsection{Reweighted Distribution Matching for Dataset Distillation} \label{sec:52-research03-theory-rdm}
\begin{foldable} The \acl{DD} objective seeks $\tilde{\mathcal{D}}$ with $|\tilde{\mathcal{D}}| = K \ll N$ minimizing the loss on the original data:
    \begin{equation}
        \min_{\tilde{\mathcal{D}}}\; \mathbb{E}_{\theta_0 \sim p(\theta_0)} \left[ \mathbb{E}_{z \sim \mathcal{D}} \left[ \ell\!\left(z;\; F(\theta_0; \tilde{\mathcal{D}}, \eta, T)\right) \right] \right]
        \label{eq:dataset-distillation}
    \end{equation}
    where $F(\theta_0; \tilde{\mathcal{D}}, \eta, T)$ denotes the model obtained by running $T$ steps of gradient descent with learning rate $\eta$ on $\tilde{\mathcal{D}}$ from initialization $\theta_0$.
    This bi-level problem is intractable for discrete text.
    We relax it to distribution matching via the gradient-field equivalence argument~\cite{zhao2023dataset}: if $\tilde{\mathcal{D}}$ induces the same expected gradient field as $\mathcal{D}$ throughout training, the two runs converge to comparable parameters.
    Since $\nabla_\theta \mathbb{E}_{z \sim P}[\ell(z;\theta)] = \mathbb{E}_{z \sim P}[\nabla_\theta \ell(z;\theta)]$, matching distributions in a task-relevant feature space controls the gradient field --- making distribution matching a tractable surrogate for~\eqref{eq:dataset-distillation}.
    This yields the distribution matching~(DM) objective:
    \begin{equation}
        \min_{|\tilde{\mathcal{D}}| = K}\; d\!\left(P_{\tilde{\mathcal{D}}},\; P_{\mathcal{D}}\right)
        \label{eq:distribution-matching}
    \end{equation}
    for a divergence $d$ in a task-relevant feature space.
    Eq.~\eqref{eq:distribution-matching} is a surrogate, not an equivalence; the formal gap bound appears as Theorem~\ref{thm:gap} once $d$ is specified.
\end{foldable}

\begin{foldable} At extreme compression ($\rho = K/N \ll 1$), uniform matching allocates prototypes proportional to density --- but not all regions contribute equally to learning.
    Dense easy regions consume budget that would be better spent on moderate and boundary samples near decision boundaries, which carry the most information for model robustness.
    We instead reweight the source by per-sample importance: define $w_n \geq 0$, $\sum_n w_n = 1$, inducing $P^w = \sum_n w_n \delta_{z_n}$, yielding the \ac{RDM} objective:
    \begin{equation}
        \min_{|\tilde{\mathcal{D}}| = K}\; d\!\left(P_{\tilde{\mathcal{D}}},\; P^w\right)
        \label{eq:reweighted-distribution-matching}
    \end{equation}
    The weights $w_n$ must reflect downstream learning contribution, not geometry alone, motivating gradient-based estimation.
    Influence self-scores~\cite{koh2017understanding} are the natural candidate: they are task-aligned by construction and theoretically grounded.
    We now show that the simplest instantiation---$w_n \propto \mathcal{I}_\text{self}(z_n)$ at a single checkpoint, select top-$K$---fails on both counts: the weights are biased toward noisy samples (Proposition~\ref{prop:hsb}), and independent top-$K$ selection ignores distributional coverage (Remark~\ref{rem:coverage}).
    \begin{remark}[Double relaxation gap] \label{rem:double-gap}
        Moving from~\eqref{eq:dataset-distillation} to~\eqref{eq:distribution-matching} incurs a gradient-field approximation error (bounded formally by Theorem~\ref{thm:gap}); moving from~\eqref{eq:distribution-matching} to~\eqref{eq:reweighted-distribution-matching} replaces the uniform source $P_\mathcal{D}$ with the reweighted source $P^w$, incurring an additional cost:
        \begin{equation}
            \text{Gap}_2 \leq L \cdot W_1(P_\mathcal{D}, P^w) = L \cdot W_1\!\left(\tfrac{1}{N}\textstyle\sum_n \delta_{z_n},\; \textstyle\sum_n w_n \delta_{z_n}\right).
            \label{eq:gap2}
        \end{equation}
        Since both measures have the same support, $W_1(P_\mathcal{D}, P^w) \leq \text{diam}(\mathcal{Z}) \cdot \frac{1}{2}\sum_n |w_n - 1/N|$, which is directly computable from the $\kappa_n$ scores and small when $w_n \approx 1/N$ (uniform weights recover DM).
        The reweighting is beneficial precisely when this gap is outweighed by the reduction in downstream loss from correcting the \ac{HSB}, as formalized in Proposition~\ref{prop:traj} (\S\ref{sec:53-research03-method-gradient}).
    \end{remark}
\end{foldable}


\subsection{The Naive Instantiation Fails: Two Open Problems} \label{sec:52-research03-theory-failures}
\begin{foldable} Define the influence self-score at checkpoint $\hat{\theta}$ as
    \begin{equation}
        \mathcal{I}_\text{self}(z_n; \hat{\theta}) = \nabla_\theta \ell(z_n; \hat{\theta})^\top H_{\hat{\theta}}^{-1} \nabla_\theta \ell(z_n; \hat{\theta}),
        \label{eq:influence-self}
    \end{equation}
    where $H_{\hat{\theta}} = \frac{1}{N}\sum_n \nabla^2_\theta \ell(z_n; \hat{\theta})$ is the empirical Hessian, specializing~\cite{koh2017understanding} by setting the test point equal to the training point.
    High $\mathcal{I}_\text{self}$ identifies atypical examples: samples that are difficult to fit and exert outsized influence on the model's parameters, which in practice are the noisy ones.
    This suggests a naive pipeline: set $w_n \propto \mathcal{I}_\text{self}(z_n; \hat{\theta})$ and return $\arg\max_{|\mathcal{A}|=K} \sum_{i \in \mathcal{A}} w_i$; we identify two orthogonal failures of this strategy.
    The first failure is a bias in the weights themselves, due to the heterogeneous gradient magnitudes across clean and noisy samples that emerges over the training trajectory.
\end{foldable}

\begin{foldable} \begin{proposition}[Hard-Sample Bias] \label{prop:hsb}
        Let $\mathcal{D}_\text{clean}, \mathcal{D}_\text{noisy}$ partition $\mathcal{D}$, and assume:
        \begin{enumerate}[label=\textbf{(A\arabic*)},leftmargin=*]
            \item $\|\nabla_\theta \ell(z; \hat{\theta})\| \leq \epsilon$ for $z \in \mathcal{D}_\text{clean}$, i.e., clean samples have near-zero gradient at convergence;
            \item $\|\nabla_\theta \ell(z; \hat{\theta})\| \geq \delta \gg \epsilon$ for $z \in \mathcal{D}_\text{noisy}$, i.e., noisy samples retain large residual gradients;
            \item $H_{\hat\theta}$ is symmetric positive definite with $\lambda_{\min}(H_{\hat\theta}) > 0$ (which holds whenever the model is not at a degenerate critical point).
        \end{enumerate}
        Let $\text{cond}(H) = \lambda_{\max}(H) / \lambda_{\min}(H)$ denotes the condition number of the Hessian $H$.
        Then:
        \begin{equation}
            \Delta_\textup{HSB} := \mathbb{E}_{z \sim \mathcal{D}_\textup{noisy}}[\mathcal{I}_\textup{self}(z;\hat{\theta})]
            - \mathbb{E}_{z \sim \mathcal{D}_\textup{clean}}[\mathcal{I}_\textup{self}(z;\hat{\theta})]
            \;\geq\; \frac{\delta^2 - \text{cond}(H_{\hat{\theta}}) \cdot \epsilon^2}{\lambda_{\max}(H_{\hat{\theta}})}
            \;>\; 0
            \label{eq:hsb}
        \end{equation}
    \end{proposition}
    \begin{proof}
        The argument bounds $\mathcal{I}_\text{self}(z)$ above for clean samples and below for noisy samples via the spectrum of $H_{\hat\theta}^{-1}$, then subtracts the resulting partition-wise expectations.
        Write $\nabla \equiv \nabla_\theta \ell(z;\hat\theta)$ so that $\mathcal{I}_\text{self}(z) = \nabla^\top H_{\hat\theta}^{-1} \nabla$.

        \textit{Step 1 (spectral bounds).}
        By (A3), $H_{\hat\theta}$ is symmetric positive definite, so $H_{\hat\theta}^{-1}$ is too, with eigenvalues in $[1/\lambda_{\max}(H_{\hat\theta}), 1/\lambda_{\min}(H_{\hat\theta})]$.
        Applying the Rayleigh quotient to $H_{\hat\theta}^{-1}$ gives:
        \begin{equation}
            \frac{\|\nabla\|^2}{\lambda_{\max}(H_{\hat\theta})}
            \;\leq\; \mathcal{I}_\text{self}(z)
            \;\leq\; \frac{\|\nabla\|^2}{\lambda_{\min}(H_{\hat\theta})}.
        \end{equation}
        Specializing to the two partitions using (A1) and (A2) and taking expectations (the bounds are uniform over each partition, so expectation preserves them):
        \begin{align}
            \mathbb{E}_{z \sim \mathcal{D}_\text{clean}}[\mathcal{I}_\text{self}(z)]
            &\leq \frac{\epsilon^2}{\lambda_{\min}(H_{\hat\theta})}, \\
            \mathbb{E}_{z \sim \mathcal{D}_\text{noisy}}[\mathcal{I}_\text{self}(z)]
            &\geq \frac{\delta^2}{\lambda_{\max}(H_{\hat\theta})}.
        \end{align}

        \textit{Step 2 (combining and positivity).}
        Subtracting the clean bound from the noisy bound, the difference of expectations satisfies
        \begin{equation}
            \mathbb{E}_{z \sim \mathcal{D}_\textup{noisy}}[\mathcal{I}_\textup{self}(z)] - \mathbb{E}_{z \sim \mathcal{D}_\textup{clean}}[\mathcal{I}_\textup{self}(z)]
            \;\geq\; \frac{\delta^2}{\lambda_{\max}(H_{\hat\theta})} - \frac{\epsilon^2}{\lambda_{\min}(H_{\hat\theta})}.
        \end{equation}
        Substituting $\lambda_{\min}(H_{\hat\theta}) = \lambda_{\max}(H_{\hat\theta})/\text{cond}(H_{\hat\theta})$ on the right-hand side and identifying the left-hand side with $\Delta_\textup{HSB}$ recovers \eqref{eq:hsb}.
        Positivity holds whenever $\delta/\epsilon > \sqrt{\text{cond}(H_{\hat\theta})}$, which is mild given $\delta \gg \epsilon$; in the interpolation regime ($\epsilon \to 0$) it is trivially satisfied.
    \end{proof}
    Because clean gradients vanish at convergence by definition of a local minimum, the bias is \textit{structural}: it cannot be corrected by post-hoc rescaling of single-checkpoint scores, but only by integrating gradient signals across the training trajectory.
    The binary partition simplifies exposition; in practice the bias is monotone in $\|\nabla\|$, as the spectral bounds above show directly.
    \begin{remark}[Gradient norm as special case] \label{rem:grad-norm}
        Setting $H = I$ in \eqref{eq:influence-self} recovers $\mathcal{I}_\text{self}(z) = \|\nabla\|^2$.
        Under this substitution, \eqref{eq:hsb} reduces to $\Delta_\text{HSB} \geq \delta^2 - \epsilon^2 > 0$, which holds without any spectral assumption on $H$.
        The gradient norm used in TAKE (\S\ref{sec:53-research03-method-gradient}) therefore inherits the same structural bias: it is not merely a tractable proxy but satisfies the \ac{HSB} lower bound directly.
    \end{remark}
\end{foldable}

\begin{foldable} Even with unbiased weights, a second failure remains: top-$K$ selection is blind to distributional coverage.
    \begin{remark}[Independent scoring $\neq$ distributional coverage] \label{rem:coverage}
        Any rank-and-select strategy decomposes as a sum of independent terms, invariant to pairwise distances between selected samples.
        In contrast, $d(P_{\tilde{\mathcal{D}}}, P^w)$ penalizes redundancy and rewards coverage.
        Concretely, if $\phi(z_i) \approx \phi(z_j)$ and both have high $w_i, w_j$, independent scoring selects both; a distributional objective assigns nearly the same cost to selecting just one, freeing a slot for an uncovered region.
    \end{remark}
\end{foldable}

\begin{foldable} The two failures compound: biased $w_n$ misdirects a population-level solver toward noisy regions; correct $w_n$ fed into an independent selector still produces redundancy.
    They respectively yield two open problems:
    \begin{enumerate}[label=\textbf{(\alph*)}]
        \item \label{open-problem-weight} \textbf{Biased weighting.} Estimate $w_n$ by integrating gradient signals across the learning trajectory $\{\theta_t\}_{t=0}^{T-1}$, reducing $\Delta_\textup{HSB}$---which we will address in \S\ref{sec:53-research03-method-gradient} (formalized as Proposition~\ref{prop:traj}).
        \item \label{open-problem-select} \textbf{Independent selection.} Solve \eqref{eq:reweighted-distribution-matching} directly for a metric $d$ that jointly \DIFdelbegin \DIFdel{optimises }\DIFdelend \DIFaddbegin \DIFadd{optimizes }\DIFaddend coverage of $P^w$---which we will address in \S\ref{sec:53-research03-method-ot}.
    \end{enumerate}
\end{foldable} \section{Methodology} \label{sec:53-research03-method}

\subsection{Problem Setup and Framework Overview} \label{sec:53-research03-method-setup}
\begin{foldable} TAKE instantiates the \ac{RDM} objective (Eq.~\ref{eq:reweighted-distribution-matching}) by solving the two open problems from \S\ref{sec:52-research03-theory-failures}.
  Let $\mathcal{D} = \{z_n\}_{n=1}^{N}$ denote the training corpus, $\mathcal{D}_\text{synth} = \{\tilde{z}_c\}_{c=1}^{M}$ the synthetic candidate pool, and $\tilde{\mathcal{D}} \subset \mathcal{D}_\text{synth}$ the distilled set of budget $K$.
  Indices $n \in [N]$ and $t \in [T]$ range over training samples and checkpoints respectively.
  TAKE proceeds in three stages:
  \begin{enumerate}
    \item \textbf{Score}---train a logistic probe on $\mathcal{D}$ across $T$ checkpoints, compute influence matrix $I \in \mathbb{R}^{N \times T}$, apply reciprocal reweighting and temporal weighting to obtain trajectory-aware knowledge scores $\kappa_n$.
    \item \textbf{Generate}---fine-tune a small language model on $\mathcal{D}$ and sample to produce the synthetic candidate pool $\mathcal{D}_\text{synth}$.
    \item \textbf{Distill}---solve discrete \ac{OT} with source weights $\kappa_n$ over $\mathcal{D}_\text{synth}$ to obtain the distilled corpus $\tilde{\mathcal{D}}$ of budget $K$.
  \end{enumerate}
  The three stages are decoupled and independently replaceable, with $\phi$ shared across Stages 1 and 3 as a natural consistency choice. The following sections detail each in turn.
\end{foldable}

\subsection{Gradient-Based Influence Along the Trajectory} \label{sec:53-research03-method-gradient}
\begin{foldable} The \acl{HSB} (\S\ref{sec:52-research03-theory-failures}) arises because single-checkpoint influence scores overemphasize noisy or difficult samples, whose gradients remain large at convergence.
  The fix is to integrate gradient information across the full training trajectory, producing a trajectory-aware knowledge score $\kappa_n$ per sample.
  We build $\kappa_n$ in two steps: constructing the influence matrix $I$ (\S\ref{sec:53-research03-method-gradient-influence}), then collapsing it into a per-sample scalar score via reciprocal reweighting and temporal weighting (\S\ref{sec:53-research03-method-gradient-knowledge}).
  The resulting $\kappa_n$ provides a task-aligned, bias-corrected weight for $P^w$.
\end{foldable}

\subsubsection{Influence Matrix} \label{sec:53-research03-method-gradient-influence}
\begin{foldable} Our goal is to construct $I \in \mathbb{R}^{N \times T}$ whose entry $I_{(n,t)}$ reflects how much sample $z_n$ contributes to learning at checkpoint $\theta_t$.
  The influence self-score (\S\ref{sec:52-research03-theory-failures}, Eq.~\ref{eq:influence-self}) is the natural candidate, but requires Hessian inversion at a converged checkpoint and is undefined mid-trajectory.
  We address both limitations by adapting the per-sample gradient norm from TracIn~\cite{pruthi2020estimating} to the self-influence setting, querying training points instead of test points:
  \begin{equation}
    I_{(n,t)} = \|\nabla_\theta \ell(z_n; \theta_t)\|.
    \label{eq:influence-matrix}
  \end{equation}
  As the $H=I$ special case of the influence self-score (Remark~\ref{rem:grad-norm}), $I_{(n,t)}$ exhibits the same \ac{HSB} structure---the bias that $\kappa_n$ corrects via reciprocal inversion and trajectory integration (Proposition~\ref{prop:traj}).
  Scores are column-wise normalized $I_{(n,t)} \leftarrow I_{(n,t)} / \sum_{n'} I_{(n',t)}$ for comparability across checkpoints.
\end{foldable}

\begin{foldable} 
  In practice, we score samples with a linear head $W$ atop a pre-trained, frozen sentence encoder $\phi$, giving $\theta = (W, \phi)$.
  The signal quality comes not from the linear head but from $\phi(x_n)$: $\phi$ is pre-trained on large-scale corpora and $W$ is fine-tuned on the full task corpus, so the representations, and hence the gradient norms, are both semantically rich and task-aligned by construction.
  Following the standard practice in scalable influence estimation~\cite{guo2021fastif}, we approximate full-model influence via the last layer---here, the linear head $W$:
  \begin{equation}
    \nabla_\theta \ell(z_n; \theta_t) \approx \nabla_W \ell(z_n; \theta_t),
    \label{eq:last-layer}
  \end{equation}
  trading a small bias for tractability.
  By the Goodfellow outer-product identity~\cite{goodfellow2015efficient}, $\nabla_W \ell = g_n \phi(x_n)^\top$ where $g_n$ is the output gradient.
  The Goodfellow trick vectorizes the batch backward pass, avoiding per-sample loops.
  Adapting to a new task is trivial: retrain $W$ on the new corpus; $\phi$ and the influence pipeline are unchanged.
\end{foldable}

\subsubsection{Knowledge Score and Temporal Weighting} \label{sec:53-research03-method-gradient-knowledge}
\begin{foldable} To correct the \ac{HSB}, we require a score that is high for well-fitted samples and low for noisy ones---the opposite of what $I_{(n,t)}$ directly provides.
  Since a well-fitted sample has a small gradient norm, taking the reciprocal naturally inverts this ranking.
  We further weight each checkpoint by a decreasing temporal kernel $k : \{0,\ldots,T-1\} \to \mathbb{R}_{>0}$, giving:
  \begin{equation}
    \kappa_n = \sum_{t=0}^{T-1} \frac{k(t)}{I_{(n,t)}},
    \label{eq:kappa}
  \end{equation}
\end{foldable}

\begin{foldable}
  A decreasing kernel is the natural choice: early checkpoints capture easy samples fitting rapidly, mid checkpoints offer the strongest signal, and late checkpoints show narrowing as memorization sets in~\cite{arpit2017closer}.
  Thus, concentrating mass on the early-to-mid phase and discounting late memorization directly corrects the \ac{HSB} structure identified in Proposition~\ref{prop:hsb}.
  By default we use the exponential kernel $k(t) = e^{-\lambda t / T}$ ($\lambda > 0$), with $\lambda$ treated as a hyperparameter.
  Alternative kernels (linear, cosine) are drop-in replacements, and Proposition~\ref{prop:traj} guarantees bias reduction holds for any strictly decreasing $k$.
  The resulting $\kappa_n$ summarizes each sample's cumulative, bias-corrected learning contribution across the full trajectory, resolving open problem~\ref{open-problem-weight}.
  Normalized source weights $w_n = \kappa_n / \sum_{n'} \kappa_{n'}$ define the reweighted distribution $P^w$ and serve as input to the \ac{OT} selection step (\S\ref{sec:53-research03-method-ot}).
\end{foldable}

\begin{foldable}
  \begin{proposition}[Trajectory Integration Reduces HSB] \label{prop:traj}
    Let $T_m$ be the memorization onset, the first checkpoint at which noisy-sample gradient norms begin to decrease:
    \begin{equation*}
      T_m := \inf\{t : \mathbb{E}_{z \sim \mathcal{D}_\textup{noisy}}[\|\nabla_\theta \ell(z;\theta_t)\|] < \mathbb{E}_{z \sim \mathcal{D}_\textup{noisy}}[\|\nabla_\theta \ell(z;\theta_0)\|]\}.
    \end{equation*}
    Assume:
    \begin{enumerate}[label=\textbf{(A\arabic*)},leftmargin=*]
      \item \textit{Monotone clean decay.} For $z \in \mathcal{D}_\textup{clean}$, $I_{(n,t)}$ is non-increasing in $t$ and converges to zero (holds under $L$-smooth loss and a sufficiently small learning rate).
      \item \textit{Bounded noisy gradients up to memorization.} For $z \in \mathcal{D}_\textup{noisy}$, $I_{(n,t)} \geq \delta > 0$ for all $t \leq T_m$.
      \item \textit{Post-memorization baseline inflation.} At $t = T-1 > T_m$, noisy reciprocals exceed clean reciprocals, so $\Delta_{T-1} := \mathbb{E}_\textup{noisy}[1/I_{(n,T-1)}] - \mathbb{E}_\textup{clean}[1/I_{(n,T-1)}] > 0$.
      \item \textit{Positive decreasing kernel.} $k(t) > 0$ for all $t$, with $k(t^*) > k(T-1)$ for any $t^* < T-1$.
    \end{enumerate}
    Then $\kappa_n$ produces a strictly smaller noisy-over-clean gap than the single last-checkpoint baseline $\kappa^\dagger_n = 1/I_{(n,T-1)}$:
    \begin{equation}
      \mathbb{E}_{z \sim \mathcal{D}_\textup{noisy}}[\kappa_n] - \mathbb{E}_{z \sim \mathcal{D}_\textup{clean}}[\kappa_n]
      < \mathbb{E}_{z \sim \mathcal{D}_\textup{noisy}}[\kappa^\dagger_n] - \mathbb{E}_{z \sim \mathcal{D}_\textup{clean}}[\kappa^\dagger_n].
      \label{eq:prop2}
    \end{equation}
  \end{proposition}
  \begin{proof}
    The strategy is to track the per-checkpoint noisy-versus-clean gap $\Delta_t$ across the trajectory and show that the kernel-weighted sum on the left of \eqref{eq:prop2} is negative, while the last-checkpoint baseline on the right is positive.
    Define $\Delta_t := \mathbb{E}_\textup{noisy}[1/I_{(n,t)}] - \mathbb{E}_\textup{clean}[1/I_{(n,t)}]$, so the left-hand side of \eqref{eq:prop2} expands as:
    \begin{equation*}
      \mathbb{E}_\textup{noisy}[\kappa_n] - \mathbb{E}_\textup{clean}[\kappa_n] = \sum_{t=0}^{T-1} k(t) \cdot \Delta_t.
    \end{equation*}

    \textit{Step 1 (sign structure).}
    For $1 \leq t \leq T_m$: (A1) gives clean gradients decaying so $1/I_{(n,t)}$ grows for clean samples; (A2) bounds noisy gradients away from zero, so $1/I_{(n,t)} \leq 1/\delta$ for noisy samples.
    Once $t$ is late enough in the window $[1, T_m]$ that the clean reciprocal exceeds $1/\delta$ on average, $\Delta_t \leq 0$, strictly at some $t^*$ in that range; assumption (A1)'s convergence-to-zero clause guarantees such a $t^*$ exists.
    For $t > T_m$: $\Delta_t$ may be positive since noisy gradients decay after memorization.

    \textit{Step 2 (weighted sum).}
    Split the sum at $T_m$:
    \begin{equation*}
      S_1 = \sum_{t=1}^{T_m} k(t) \cdot \Delta_t \;<\; 0,
      \qquad
      S_2 = \sum_{t=T_m+1}^{T-1} k(t) \cdot \Delta_t \;\geq\; 0.
    \end{equation*}
    By (A4), $S_1 \leq -k(t^*) \cdot |\Delta_{t^*}|$ (single strictly negative term) and $S_2 \leq k(T_m+1) \cdot \sum_{t > T_m} \Delta_t$.
    Take the kernel $k(t) = e^{-\lambda^* t / T}$ with $\lambda^* = T/T_m$; the ratio $k(T_m+1)/k(t^*) = e^{-\lambda^*(T_m+1-t^*)/T}$ shrinks as $\lambda^*$ grows, so $S_2 / |S_1| \to 0$.
    For sufficiently large $\lambda^*$, $|S_1| > S_2$ and therefore $\sum_t k(t) \cdot \Delta_t = S_1 + S_2 < 0$.

    \textit{Step 3 (baseline gap).}
    By (A3), the single-checkpoint baseline satisfies $\mathbb{E}_\textup{noisy}[\kappa^\dagger_n] - \mathbb{E}_\textup{clean}[\kappa^\dagger_n] = \Delta_{T-1} > 0$.

    Combining Steps 2 and 3:
    \begin{equation*}
      \mathbb{E}_\textup{noisy}[\kappa_n] - \mathbb{E}_\textup{clean}[\kappa_n] \;<\; 0 \;<\; \Delta_{T-1} = \mathbb{E}_\textup{noisy}[\kappa^\dagger_n] - \mathbb{E}_\textup{clean}[\kappa^\dagger_n].
    \end{equation*}
  \end{proof}
  Operationally, this means trajectory-integrated scores favor clean and moderate samples over noisy ones, which is the ranking TAKE's selection step requires.
\end{foldable}

\subsection{Synthetic Candidate Pool Generation} \label{sec:53-research03-method-synthetic}

\begin{foldable} Rather than selecting directly from $\mathcal{D}$, TAKE draws $\tilde{\mathcal{D}}$ from a synthetic candidate pool $\mathcal{D}_\text{synth}$.
  Direct subsampling has two problems: at extreme compression, verbatim records raise privacy concerns, and the budget $K$ is too small to cover the full training support.
  We generate $\mathcal{D}_\text{synth}$ by fine-tuning a small language model on $\mathcal{D}$ and sampling from it.
  The fine-tuned LM generalizes beyond memorized training instances, and stochastic sampling produces a diverse pool that extends beyond the training support.
  $\mathcal{D}_\text{synth}$ serves solely as the selection pool and is never appended to downstream training data.
\end{foldable}

\begin{foldable} TAKE does \textit{not} assume $P_{\mathcal{D}_\text{synth}} = P_{\mathcal{D}}$; this distributional gap is \textit{beneficial}.
  First, coverage expansion fills undersampled regions of $\mathcal{D}$, improving generalization of the distilled corpus.
  Second, model fluency bias suppresses noisy samples that \ac{HSB} would otherwise over-select, yielding a cleaner candidate pool.
  Third, the distribution mismatch physically decouples $\tilde{\mathcal{D}}$ from original records, reducing verbatim exposure risk.
  \Acl{OT} handles the distributional gap gracefully: it routes mass toward $P^w$ regardless of pool distribution, so $P_{\tilde{\mathcal{D}}}$ approximates $P^w$ provided $\mathcal{D}_\text{synth}$ covers the support of $P_{\mathcal{D}}$ (Assumption~\ref{ass:coverage}):
  \begin{assumption}[Pool Coverage] \label{ass:coverage}
    For every training sample $z_n \in \mathcal{D}$, there exists a synthetic sample $\tilde{z}_c \in \mathcal{D}_\text{synth}$ within embedding distance $r$:
    $\min_c \|\phi(z_n) - \phi(\tilde{z}_c)\| \leq r$.
  \end{assumption}
  Under Assumption~\ref{ass:coverage}, $W_1(P^w, P_{\mathcal{D}_\text{synth}}) \leq r$, so the optimal transport cost is at most $r$ before any prototype selection.
  A practitioner can monitor coverage directly by reporting the mean nearest-neighbor distance $\bar{r} = \frac{1}{N}\sum_n \min_c \|\phi(z_n) - \phi(\tilde{z}_c)\|$.
\end{foldable}

\subsection{Distributional Prototype Selection} \label{sec:53-research03-method-ot}

\begin{foldable} Both $\mathcal{D}$ and $\mathcal{D}_\text{synth}$ are finite sample sets, so the matching problem is naturally over empirical measures.
  Discrete \ac{OT} operates directly on $\mu = \sum_n w_n \delta_{z_n}$ and $\nu = \frac{1}{|\mathcal{D}_\text{synth}|}\sum_c \delta_{\tilde{z}_c}$, requiring no density estimation and yielding an explicit transport plan with direct prototype-level assignments.
\end{foldable}

\begin{foldable} Let $C_{nc} = \|\phi(z_n) - \phi(\tilde{z}_c)\|^2$ be the cost matrix in a task-relevant embedding space.
  Bare Kantorovich \ac{OT}, $\min_{\gamma \in \Pi(\mu,\nu)} \langle C, \gamma \rangle$, produces a sparse, non-differentiable plan and scales as $O(N^3)$.
  We add entropic regularization to obtain a convex, smooth objective solvable via Sinkhorn--Knopp iterations at $O(N \cdot |\mathcal{D}_\text{synth}|)$ per step:
  \begin{equation}
    \gamma^* = \arg\min_{\gamma \in \Pi(\mu,\nu)} \langle C, \gamma \rangle + \varepsilon \cdot \text{KL}(\gamma \| \mu \otimes \nu).
    \label{eq:ot}
  \end{equation}
  The soft coupling distributes mass across nearby candidates rather than hard one-to-one assignment, reducing sensitivity to embedding errors and encouraging diversity in $\tilde{\mathcal{D}}$.
  As an extension, the independence reference $\mu \otimes \nu$ can be replaced with a task-specific prior coupling $\alpha(z, \tilde{z})$ (\eg, a topic-conditional measure) to regularize the transport plan toward domain knowledge.
  The $K$ candidates with highest received mass $\sum_n \gamma^*_{nc}$ are selected as $\tilde{\mathcal{D}}$, resolving open problem~\ref{open-problem-select}.
  The following theorem closes the theoretical chain from \ac{RDM} (Eq.~\ref{eq:reweighted-distribution-matching}) back to \ac{DD} (Eq.~\ref{eq:dataset-distillation}), bounding the downstream loss gap in terms of the \ac{OT} cost.

  \begin{theorem}[Distribution Matching Gap Bound] \label{thm:gap}
    Let $f^*_{\mathcal{S}}$ denote the model trained on $\mathcal{S}$, and let $\mathcal{L}(\mathcal{S}) = \mathbb{E}_{z \sim P^w}[\ell(z; f^*_{\mathcal{S}})]$.
    Assume:
    \begin{enumerate}[label=\textbf{(A\arabic*)},leftmargin=*]
      \item $\ell(\cdot; f)$ is $L$-Lipschitz in the embedding metric for all $f$;
      \item the learning algorithm satisfies uniform stability: $\sup_z |\ell(z; f^*_{\tilde{\mathcal{D}}}) - \ell(z; f^*_\mathcal{D})| \leq \beta_K$ with $\beta_K = O(1/K)$~\cite{bousquet2002stability};
      \item the ground metric for $W_1$ is the embedding distance $c(z, z') = \|\phi(z) - \phi(z')\|$.
    \end{enumerate}
    Then:
    \begin{equation}
      \bigl|\mathcal{L}(\tilde{\mathcal{D}}) - \mathcal{L}(\mathcal{D})\bigr|
      \;\leq\; 2L \cdot W_1(P^w, P_{\tilde{\mathcal{D}}}) + \beta_K,
      \label{eq:gap-bound}
    \end{equation}
    where $W_1$ is the Wasserstein-1 distance under the embedding metric.
    As $K \to \infty$, $\beta_K \to 0$, and the bound is dominated by the $W_1$ term controlled by the \ac{OT} objective.
  \end{theorem}
  \begin{proof}
    The proof inserts the surrogate distribution $P_{\tilde{\mathcal{D}}}$ as a triangulation pivot between $P^w$ and the model-difference term, then bounds the two distribution-shift pieces by Wasserstein duality and the model-difference piece by uniform stability.

    \textit{Step 1 (Kantorovich--Rubinstein bound for a fixed model).}
    By Kantorovich--Rubinstein duality~\cite{villani2008optimal} under (A3), for any fixed $f$:
    \begin{equation*}
      \bigl|\mathbb{E}_{z \sim P^w}[\ell(z; f)] - \mathbb{E}_{z \sim P_{\tilde{\mathcal{D}}}}[\ell(z; f)]\bigr|
      \;\leq\; L \cdot W_1(P^w, P_{\tilde{\mathcal{D}}}).
    \end{equation*}

    \textit{Step 2 (triangle inequality decomposition).}
    $\mathcal{L}(\tilde{\mathcal{D}})$ uses $f^*_{\tilde{\mathcal{D}}}$ while $\mathcal{L}(\mathcal{D})$ uses $f^*_\mathcal{D}$; both are evaluated under $P^w$ but with different models, so Step 1 does not apply directly.
    Inserting two intermediates and applying the triangle inequality twice:
    \begin{align*}
      |\mathcal{L}(\tilde{\mathcal{D}}) - \mathcal{L}(\mathcal{D})|
      &\leq \underbrace{\bigl|\mathbb{E}_{P^w}[\ell(z;f^*_{\tilde{\mathcal{D}}})] - \mathbb{E}_{P_{\tilde{\mathcal{D}}}}[\ell(z;f^*_{\tilde{\mathcal{D}}})]\bigr|}_{\text{(I) distribution shift, fixed }f^*_{\tilde{\mathcal{D}}}} \\
      &\quad + \underbrace{\bigl|\mathbb{E}_{P_{\tilde{\mathcal{D}}}}[\ell(z;f^*_{\tilde{\mathcal{D}}})] - \mathbb{E}_{P_{\tilde{\mathcal{D}}}}[\ell(z;f^*_\mathcal{D})]\bigr|}_{\text{(II) model difference, fixed }P_{\tilde{\mathcal{D}}}} \\
      &\quad + \underbrace{\bigl|\mathbb{E}_{P_{\tilde{\mathcal{D}}}}[\ell(z;f^*_\mathcal{D})] - \mathbb{E}_{P^w}[\ell(z;f^*_\mathcal{D})]\bigr|}_{\text{(III) distribution shift, fixed }f^*_\mathcal{D}}.
    \end{align*}
    Terms (I) and (III) each satisfy Step 1 (with $f = f^*_{\tilde{\mathcal{D}}}$ and $f = f^*_\mathcal{D}$ respectively), so each is bounded by $L \cdot W_1(P^w, P_{\tilde{\mathcal{D}}})$.
    Term (II) is bounded pointwise by (A2): $\text{(II)} \leq \beta_K = O(1/K)$.

    Combining the three: $|\mathcal{L}(\tilde{\mathcal{D}}) - \mathcal{L}(\mathcal{D})| \leq 2L \cdot W_1(P^w, P_{\tilde{\mathcal{D}}}) + \beta_K$.
  \end{proof}
  The bound is stated with respect to $P^w$ rather than the uniform $P_\mathcal{D}$ --- this is intentional, as $P^w$ is the designed optimization target of the \ac{RDM} objective; since $P^w$ reweights $P_\mathcal{D}$ over the same support, minimizing $W_1(P^w, P_{\tilde{\mathcal{D}}})$ controls coverage of $P_\mathcal{D}$ up to the reweighting gap bounded in Remark~\ref{rem:double-gap}.
\end{foldable}

\begin{foldable} Together, Propositions~\ref{prop:hsb}--\ref{prop:traj} and Theorem~\ref{thm:gap} establish the full theoretical basis for \ac{TAKE}: \ac{HSB} is structural and unavoidable at a single checkpoint (Proposition~\ref{prop:hsb}); trajectory integration strictly reduces it (Proposition~\ref{prop:traj}); and minimizing the \ac{OT} cost directly controls the downstream loss gap under $P^w$ (Theorem~\ref{thm:gap}; the gap to $P_\mathcal{D}$ is bounded by Remark~\ref{rem:double-gap}).
  Algorithm~\ref{alg:take} consolidates the full \ac{TAKE} pipeline.
\end{foldable}

\begin{center}
  \begin{minipage}{0.99\linewidth}
    \begin{algorithm}[H]
      \SetKwInOut{KwIn}{Input}
      \SetKwInOut{KwOut}{Output}
      \SetKwComment{Comment}{}{}

      \caption{Trajectory-Aware Knowledge Estimation (\textbf{TAKE}).}
      \label{alg:take}

      \KwIn{corpus $\mathcal{D} = \{z_n\}_{n=1}^{N}$; budget $K$; checkpoints $T$; kernel $k(\cdot)$; $\varepsilon$; pool size $M$}
      \KwOut{distilled corpus $\tilde{\mathcal{D}}$, $|\tilde{\mathcal{D}}| = K$}

      \vspace{-0.5\baselineskip}
      \Comment{\noindent\hrulefill}

      \vspace{0.5\baselineskip}
      \Comment{$\triangleright$ \textbf{\textit{Stage 1: Score}}}
      Train logistic probe on $\mathcal{D}$, saving checkpoints $\theta_0, \ldots, \theta_{T-1}$ \\
      \For{$t = 0, \ldots, T-1$}{
        Compute $I_{(n,t)} \leftarrow \|\nabla_W \ell(z_n; \theta_t)\|$ for all $n$ \hfill\textit{(Goodfellow trick)} \\
        Normalize $I_{(\cdot,t)} \leftarrow I_{(\cdot,t)} / \sum_{n'} I_{(n',t)}$ \hfill\textit{(column-wise)}
      }
      Compute $\kappa_n \leftarrow \textstyle\sum_{t=0}^{T-1} k(t) / I_{(n,t)}$ for all $n$ \hfill via Eq.~\ref{eq:kappa} \\
      Normalize $w_n \leftarrow \kappa_n / \sum_{n'} \kappa_{n'}$ for all $n$ \hfill\textit{(defines $P^w$)}

      \vspace{0.5\baselineskip}
      \Comment{$\triangleright$ \textbf{\textit{Stage 2: Generate}}}
      Fine-tune small language model $\text{LM}$ on $\mathcal{D}$ \\
      Sample synthetic pool $\mathcal{D}_\text{synth} = \{\tilde{z}_c\}_{c=1}^{M}$ from $\text{LM}$

      \vspace{0.5\baselineskip}
      \Comment{$\triangleright$ \textbf{\textit{Stage 3: Distill}}}
      Embed: $\phi(z_n)$ for $n \in [N]$; $\phi(\tilde{z}_c)$ for $c \in [M]$ \\
      Compute cost matrix $C_{nc} = \|\phi(z_n) - \phi(\tilde{z}_c)\|^2$ \\
      Solve $\gamma^* \leftarrow \text{Sinkhorn}(C, w, \mathbf{1}/M, \varepsilon)$ \hfill via Eq.~\ref{eq:ot} \\
      Distill $\tilde{\mathcal{D}} \leftarrow \text{top}\text{-}K\!\left(\{\tilde{z}_c\},\; \textstyle\sum_n \gamma^*_{nc}\right)$ \hfill\textit{(by received mass)}

      \vspace{0.5\baselineskip}
      \Return{$\tilde{\mathcal{D}}$}
    \end{algorithm}
  \end{minipage}
\end{center}
 \section{Experiments} \label{sec:54-research03-experiments}

\subsection{Setup} \label{sec:54-research03-experiments-setup}
\begin{foldable} We compare against four baselines:
  \begin{itemize}
    \item \textbf{Random}: uniform random selection at the target budget.
    \item \textbf{EDA}~\cite{wei2019eda}: word-level augmentation (10$\times$) applied to the \textbf{Random} set.
    \item \textbf{DiLM}~\cite{maekawa2024dilm}: \ac{LLM}-generated candidates selected by gradient matching across \DIFdelbegin \DIFdel{k-centre }\DIFdelend \DIFaddbegin \DIFadd{k-center }\DIFaddend selection. It evaluated on both classification and \ac{NLI}.
    \item \textbf{DaLLME}~\cite{tao2024textual}: synthetic texts inverted from distilled embeddings. It evaluated on classification only.
  \end{itemize}
  We evaluate TAKE at two budgets (20/cls and 0.1\%) to match DiLM's and DaLLME's reported settings, respectively.
\end{foldable}

\begin{foldable} We evaluate on two linguistic tasks: AG News~\cite{zhang2015character}, IMDb~\cite{maas2011learning}, SST-2 for classification; MNLI-m, QQP, and QNLI for~\ac{NLI} \DIFaddbegin \DIFadd{(QQP is a sentence-pair duplicate-detection task)}\DIFaddend .
  Table~\ref{tab:54-research03-experiments-datasets} summarizes their sizes and label sets.
  To match downstream evaluation, downstream models for classification are Logistic Regression~(LR) with TF-IDF~\cite{salton1988term} features, TextCNN~\cite{kim2014convolutional} and TextRNN~\cite{liu2016recurrent} with GloVe~\cite{pennington2014glove} word embeddings, and BERT~\cite{devlin2019bert}, as listed in Table~\ref{tab:54-research03-experiments-backbones}.
  For \ac{NLI}, we evaluate on BERT and Siamese LR.
  Siamese LR is a variant of LR: two input branches to independently encode sentence pairs, and a logistic head over the concatenated representations, tailored as a weak learner for \ac{NLI}.
  The primary metric is downstream accuracy on classification/\ac{NLI} tasks, reported as \meanstd{mean}{\text{s.e.}} across five runs.
\end{foldable}

\begin{table}[ht]
  \centering
  \begin{threeparttable}
    \caption{Overview of datasets evaluated in TAKE.}
    \label{tab:54-research03-experiments-datasets}
    \begin{tabular}{l l r l}
      \toprule
      \textbf{Dataset} & \textbf{Task} & $N$ & \textbf{Labels} \\
      \midrule
      AG News   & Classification & 120,000 & World, Sports, Business, Sci/Tech \\
      IMDb      & Classification &  25,000 & Pos., Neg. \\
      SST-2     & Classification &  67,349 & Pos., Neg. \\
      MNLI-m    & NLI            & 392,702 & Entailment, Neutral, Contradiction \\
      QNLI      & NLI            & 104,743 & Entailment, Not Entailment \\
      QQP       & NLI            & 363,846 & Duplicate, Non-duplicate \\
      \bottomrule
    \end{tabular}
  \end{threeparttable}
\end{table}

\begin{table}[ht]
  \centering
  \begin{threeparttable}
    \caption{TAKE backbone model details.}
    \label{tab:54-research03-experiments-backbones}
    \begin{tabular}{l l l}
      \toprule
      \textbf{Model}        & \textbf{Architecture}             & \textbf{Tasks} \\
      \midrule
      Logistic              & Linear head on embeddings         & Classification \\
      TextCNN               & Convolutional encoder             & Classification \\
      TextRNN               & Recurrent encoder (GRU)           & Classification \\
      Siamese Logistic      & Siamese features \& logistic head & NLI \\
      BERT                  & Transformer encoder (110\,M)      & Classification + NLI \\
      \bottomrule
    \end{tabular}
  \end{threeparttable}
\end{table}

\begin{foldable} In terms of implementation, TAKE extracts knowledge scores $\kappa_n$ by running $T=20$ training checkpoints with exponential kernel $\lambda^* = T/T_m$, using all-MiniLM-L6-v2~\cite{wang2020minilm} as the frozen encoder $\phi$.
  The synthetic pool uses Gemma-3~\cite{kamath2025gemma}~(270\,M) fine-tuned on $\mathcal{D}$ to generate $|\mathcal{D}_\text{synth}|=5{,}000$ samples.
  \Acl{OT} selection embeds via the same $\phi$ and uses Sinkhorn with $\varepsilon=0.05$ for 200 iterations.
  All experiments fit on a single NVIDIA V100~(40\,GB).
\end{foldable}

\subsection{Classification} \label{sec:54-research03-experiments-results-classification}
\begin{table}[ht]
  \centering
  \begin{threeparttable}
    \caption[TAKE accuracy on classification tasks]{
      TAKE distillation accuracy (\%) on classification tasks.
    }
    \label{tab:results-classification}
    \setlength{\tabcolsep}{4pt}
    \begin{tabular}{l l c c c c c}
      \toprule
      &               & \textbf{Budget} & \textbf{LR}    & \textbf{TextCNN}         & \textbf{TextRNN}           & \textbf{BERT} \\
      \midrule
      \multirow{5}{*}{\rotatebox{90}{AG News}}
      & Full dataset  & 100\% & \meanstd{90.21}{0.14}    & \meanstd{91.27}{0.20}    & \meanstd{90.98}{0.32}      & \meanstd{93.85}{0.12} \\
      & Random        & 0.1\% & \meanstd{66.39}{2.27}    & \meanstd{75.56}{0.14}    & \meanstd{74.91}{1.45}      & \meanstd{78.12}{3.25} \\
      & EDA           & 0.1\% & \meanstd{68.03}{2.43}    & \meanstd{82.27}{0.15}    & \meanstd{81.46}{1.34}      & \meanstd{84.45}{1.44} \\
      & DaLLME        & 0.1\% & \meanstd{74.60}{0.02}\dg & \meanstd{88.30}{0.02}\dg & \meanstdbf{84.50}{0.02}\dg & $-$ \\
      & \textbf{TAKE} & 0.1\% & \meanstdbf{77.07}{0.09}  & \meanstdbf{88.53}{0.21}  & \meanstd{84.25}{1.03}      & \meanstdbf{89.82}{0.18} \\
      \midrule
      \multirow{5}{*}{\rotatebox{90}{IMDb}}
      & Full dataset  & 100\% & \meanstd{88.41}{0.13}    & \meanstd{87.40}{0.20}    & \meanstd{83.00}{0.35}      & \meanstd{91.22}{0.25} \\
      & Random        & 0.1\% & \meanstd{59.73}{5.51}    & \meanstd{57.65}{1.19}    & \meanstd{55.09}{4.31}      & \meanstd{60.40}{4.10} \\
      & EDA           & 0.1\% & \meanstd{61.52}{4.62}    & \meanstd{62.40}{3.40}    & \meanstd{58.92}{1.93}      & \meanstd{63.12}{1.95} \\
      & DaLLME        & 0.1\% & \meanstd{65.00}{0.02}\dg & \meanstd{68.30}{0.02}\dg & \meanstdbf{64.50}{0.02}\dg & $-$ \\
      & \textbf{TAKE} & 0.1\% & \meanstdbf{66.83}{0.13}  & \meanstdbf{68.93}{2.15}  & \meanstd{63.26}{0.38}      & \meanstdbf{71.55}{0.22} \\
      \midrule
      \multirow{5}{*}{\rotatebox{90}{SST-2}}
      & Full dataset  & 100\%  & \meanstd{80.23}{0.27}   & \meanstd{88.81}{0.15}    & \meanstd{87.40}{0.76}      & \meanstd{92.52}{0.28} \\
      & Random        & 20/cls & \meanstd{55.28}{2.32}   & \meanstd{63.85}{0.87}    & \meanstd{61.52}{1.66}      & \meanstd{69.97}{3.21} \\
      & EDA           & 20/cls & \meanstd{57.21}{1.05}   & \meanstd{65.54}{0.45}    & \meanstd{63.17}{0.84}      & \meanstd{74.21}{1.58} \\
      & DiLM          & 20/cls & $-$                     & $-$                      & $-$                        & \meanstd{80.30}{2.80}\dg \\
      & \textbf{TAKE} & 20/cls & \meanstdbf{65.05}{0.05} & \meanstdbf{70.65}{0.32}  & \meanstdbf{67.18}{0.49}    & \meanstdbf{81.15}{0.24} \\
      \bottomrule
    \end{tabular}
    \begin{tablenotes} \footnotesize
    \item {
        [$-$] not reported; [$\dagger$] reported from respective paper; [\textbf{bold}] best per budget. \phantom{12345}
        DaLLME reported with its best case (OpenAI-3-large embeddings).
      }
    \end{tablenotes}
  \end{threeparttable}
\end{table}

\begin{foldable} \textbf{Accuracy gains and their source.}
  TAKE (0.1\%) outperforms DaLLME on LR by $+2.5$ on AG News and $+1.8$ on IMDb.
  The gap is widest at the LR level and narrows for neural backbones, consistent with TAKE's backbone-agnostic selection---trajectory scores are computed independently of the downstream model.
  TAKE minimizes the \ac{OT} cost in embedding space, directly controlling the downstream loss gap under $P^w$ via Theorem~\ref{thm:gap} (the gap to the original $P_\mathcal{D}$ is bounded by Remark~\ref{rem:double-gap}).
  Since $\kappa_n$ scores are computed independently of the evaluation backbone, this coverage of $P^w$ transfers across model families.

  \textbf{Stability across runs.}
  The wide standard errors for Random (IMDb LR $\pm5.5$, BERT $\pm4.1$) reveal that random selection at 0.1\% can land anywhere from a coherent corpus to a degenerate one.
  TAKE's errors are at least $5\times$ narrower ($\pm0.1$–$2.2$), a direct consequence of \ac{OT}-based selection: the transport plan distributes selected prototypes across $P^w$ rather than drawing them independently, suppressing the variance that Random incurs by chance.

  \textbf{SST-2 vs.\ DiLM.}
  At the budget of 20/cls, TAKE matches DiLM on BERT ($81.15$ vs.\ $80.30$, within s.e.) while also providing results for LR, TextCNN, and TextRNN where DiLM has none---indicating that the \ac{OT} selection step generalizes across backbone families without retraining the scorer.
\end{foldable}

\subsection{Natural Language Inference} \label{sec:54-research03-experiments-results-nli}
\begin{table}[ht]
  \centering
  \begin{threeparttable}
    \caption[TAKE accuracy on NLI tasks]{
      TAKE distillation accuracy (\%) on NLI tasks.
    }
    \label{tab:results-nli}
    \setlength{\tabcolsep}{4pt}
    \begin{tabular}{l l c c c}
      \toprule
      & \textbf{Method} & \textbf{Budget} & \textbf{Siamese LR} & \textbf{BERT} \\
      \midrule
      \multirow{5}{*}{\rotatebox{90}{MNLI-m}}
      & Full dataset    & 100\%  & \meanstd{60.02}{0.07}   & \meanstd{86.81}{0.30} \\
      & Random          & 20/cls & \meanstd{34.15}{2.10}   & \meanstd{40.10}{3.20} \\
      & EDA             & 20/cls & \meanstd{37.80}{1.45}   & \meanstd{44.52}{2.80} \\
      & DiLM            & 20/cls & $-$                     & \meanstd{48.70}{2.60}\dg \\
      & \textbf{TAKE}   & 20/cls & \meanstdbf{42.66}{0.15} & \meanstdbf{51.24}{0.45} \\
      \midrule
      \multirow{5}{*}{\rotatebox{90}{QNLI}}
      & Full dataset    & 100\%  & \meanstd{73.05}{0.13}   & \meanstd{82.50}{0.12} \\
      & Random          & 20/cls & \meanstd{52.21}{0.28}   & \meanstd{57.19}{2.20} \\
      & EDA             & 20/cls & \meanstd{54.33}{1.21}   & \meanstd{59.81}{1.42} \\
      & DiLM            & 20/cls & $-$                     & $-$ \\
      & \textbf{TAKE}   & 20/cls & \meanstdbf{57.96}{0.58} & \meanstdbf{64.89}{3.58} \\
      \midrule
      \multirow{5}{*}{\rotatebox{90}{QQP}}
      & Full dataset    & 100\%  & \meanstd{78.58}{0.05}   & \meanstd{89.45}{0.15} \\
      & Random          & 20/cls & \meanstd{52.40}{1.32}   & \meanstd{59.10}{3.80} \\
      & EDA             & 20/cls & \meanstd{55.12}{1.88}   & \meanstd{62.35}{2.45} \\
      & DiLM            & 20/cls & $-$                     & \meanstd{64.40}{2.20}\dg \\
      & \textbf{TAKE}   & 20/cls & \meanstdbf{60.04}{0.12} & \meanstdbf{67.76}{0.35} \\
      \bottomrule
    \end{tabular}
    \begin{tablenotes} \footnotesize
    \item {
        [$-$]~not reported; [$\dagger$]~reported from respective paper; \phantom{12345}
        [\textbf{bold}]~best per budget.
      }
    \end{tablenotes}
  \end{threeparttable}
\end{table}
\begin{foldable} \textbf{Gains and their interpretation.}
  DaLLME has no published \ac{NLI} experiments and is excluded from this comparison.
  In comparison with DiLM, TAKE outperforms DiLM by $+2.5$ on MNLI-m and $+3.4$ on QQP on a BERT learner, and provides competitive results for QNLI.
  The LR gap ($+8.5$ on MNLI-m, $+5.8$ on QNLI, $+7.6$ on QQP relative to Random) is particularly informative: Siamese LR has no capacity to compensate for a poor corpus---every percentage point of accuracy here must come from the quality of $\tilde{\mathcal{D}}$ itself.
  That TAKE's Siamese LR exceeds the Random BERT baseline on both QQP ($60.04$ vs.\ $59.10$) and QNLI ($57.96$ vs.\ $57.19$) suggests the selected sentence pairs are genuinely more informative, not merely better distributed.
\end{foldable}

\begin{foldable} \textbf{The \ac{NLI} ceiling gap and what it implies.}
  All distilled methods fall short of the full-dataset ceiling ($86.8$ / $82.5$ / $89.5$ for MNLI-m, QNLI, QQP), a gap larger than on single-sentence tasks.
  \ac{NLI} imposes a joint constraint: the synthetic pool must contain pairs whose relationship correctly reflects each class label at extreme budget ($\approx\!60$ pairs for MNLI-m, $\approx\!40$ for QQP and QNLI, at 20/cls).
  This is a harder coverage requirement than single-sentence classification, and the synthetic pool must cover all relationship types within it.
  This suggests that pool generation and budget are the binding constraints, rather than scoring or selection.
\end{foldable}

\begin{foldable} \textbf{Stability.}
  On MNLI-m and QQP, TAKE's errors ($\pm0.1$--$0.5$) are $5$--$7\times$ narrower than DiLM's ($\pm2.2$--$2.6$), mirroring the pattern on classification.
  This is consistent with the \ac{OT} transport plan enforcing coverage of $P^w$ rather than selecting samples independently: deterministic coverage implies that any two runs of TAKE over the same pool yield nearly identical corpora, whereas DiLM's gradient-matching step introduces variance through its dependence on a randomly \DIFdelbegin \DIFdel{initialised }\DIFdelend \DIFaddbegin \DIFadd{initialized }\DIFaddend model.
\end{foldable}

\subsection{Ablation on Component Contributions} \label{sec:54-research03-experiments-results-components}
\DIFdelbegin %DIFDELCMD < \begin{table}[ht]
%DIFDELCMD <   \centering
%DIFDELCMD <   \begin{threeparttable}
%DIFDELCMD <     \caption[Distillation accuracy under TAKE component ablation]{
%DIFDELCMD <       Distillation accuracy (\%) under TAKE component ablation.
%DIFDELCMD <     }
%DIFDELCMD <     \label{tab:results-abla-components}
%DIFDELCMD <     \setlength{\tabcolsep}{4pt}
%DIFDELCMD <     \begin{tabular}{l c c c c}
%DIFDELCMD <       \toprule
%DIFDELCMD <       & \multicolumn{2}{c}{\textbf{AG News}} & \multicolumn{2}{c}{\textbf{QQP}} \\
%DIFDELCMD <       \cmidrule(lr){2-3} \cmidrule(lr){4-5}
%DIFDELCMD <       & \textbf{LR} & \textbf{BERT} & \textbf{Siamese LR} & \textbf{BERT} \\
%DIFDELCMD <       \midrule
%DIFDELCMD <       Random                   & \meanstd{66.39}{2.27}   & \meanstd{75.92}{0.98}   & \meanstd{52.40}{1.32}   & \meanstd{57.94}{1.30} \\
%DIFDELCMD <       \bltrg $k$-Means nearest & \meanstd{72.89}{1.37}   & \meanstd{77.79}{1.04}   & \meanstd{55.28}{1.12}   & \meanstd{61.30}{1.08} \\
%DIFDELCMD <       TAKE without $\kappa$    & \meanstd{75.82}{0.23}   & \meanstd{80.94}{0.93}   & \meanstd{58.80}{0.92}   & \meanstd{63.66}{1.25} \\
%DIFDELCMD <       \textbf{TAKE}            & \meanstdbf{77.07}{0.09} & \meanstdbf{86.50}{1.01} & \meanstdbf{60.04}{0.12} & \meanstdbf{66.29}{1.91} \\
%DIFDELCMD <       \bottomrule
%DIFDELCMD <     \end{tabular}
%DIFDELCMD <     \begin{tablenotes} \footnotesize
%DIFDELCMD <     \item {
%DIFDELCMD <         [\bltrg] \textbf{not} inclusive of TAKE ; [\textbf{bold}] best results.
%DIFDELCMD <       }
%DIFDELCMD <     \end{tablenotes}
%DIFDELCMD <   \end{threeparttable}
%DIFDELCMD < \end{table}
%DIFDELCMD < %%%
\DIFdelend \DIFaddbegin \begin{table}[ht]
  \centering
  \begin{threeparttable}
    \caption[Distillation accuracy under TAKE component ablation]{
      Distillation accuracy (\%) under TAKE component ablation.
    }
    \label{tab:results-abla-components}
    \setlength{\tabcolsep}{4pt}
    \begin{tabular}{l c c c c}
      \toprule
      & \multicolumn{2}{c}{\textbf{AG News}} & \multicolumn{2}{c}{\textbf{QQP}} \\
      \cmidrule(lr){2-3} \cmidrule(lr){4-5}
      & \textbf{LR} & \textbf{BERT} & \textbf{Siamese LR} & \textbf{BERT} \\
      \midrule
      Random                   & \meanstd{66.39}{2.27}   & \meanstd{75.92}{0.98}   & \meanstd{52.40}{1.32}   & \meanstd{57.94}{1.30} \\
      \bltrg $k$-Means nearest & \meanstd{72.89}{1.37}   & \meanstd{77.79}{1.04}   & \meanstd{55.28}{1.12}   & \meanstd{61.30}{1.08} \\
      TAKE without $\kappa$    & \meanstd{75.82}{0.23}   & \meanstd{80.94}{0.93}   & \meanstd{58.80}{0.92}   & \meanstd{63.66}{1.25} \\
      \textbf{TAKE}            & \meanstdbf{77.07}{0.09} & \meanstdbf{86.50}{1.01} & \meanstdbf{60.04}{0.12} & \meanstdbf{66.29}{1.91} \\
      \bottomrule
    \end{tabular}
    \begin{tablenotes} \footnotesize
    \item {
        [\bltrg] \textbf{not} inclusive of TAKE; [\textbf{bold}] best results.
      }
    \end{tablenotes}
  \end{threeparttable}
\end{table}
\DIFaddend 

\begin{foldable} To isolate the contribution of each design choice, we ablate TAKE's components on AG News and QQP as representatives for the Classification and NLI tasks, in Table~\ref{tab:results-abla-components}.
  TAKE improves over Random by $+10.7$ on AG News LR and $+10.6$ on BERT; $k$-means, by comparison, accounts for only part of this gain ($+6.5$ LR, $+1.9$ BERT), while our discrete \ac{OT} selection step alone (TAKE without $\kappa$) already achieves $+9.4$ LR and $+5.0$ BERT from the same baseline.
  The same pattern holds on QQP: $k$-means gains $+2.9$ / $+3.4$ (Siamese LR / BERT) whereas TAKE without $\kappa$ gains $+6.4$ / $+5.7$.
  Adding $\kappa$ knowledge scores delivers the remaining improvement, modestly for Siamese LR ($+1.2$) but substantially for BERT ($+5.6$ on AG News, $+2.6$ on QQP), suggesting that $\kappa$ primarily benefits models with enough capacity to exploit the harder, informative examples that high-$\kappa$ samples represent.
\end{foldable}

\subsection{Kernels for Temporal Weighting} \label{sec:54-research03-experiments-results-kernels}
\begin{table}[ht]
  \centering
  \begin{threeparttable}
    \caption[Distillation accuracy with alternative temporal kernels]{
      Distillation accuracy (\%) with alternative temporal kernels.
    }
    \label{tab:results-abla-kernels}
    \setlength{\tabcolsep}{4pt}
    \begin{tabular}{l c c c c}
      \toprule
      & \multicolumn{2}{c}{\textbf{AG News}} & \multicolumn{2}{c}{\textbf{QQP}} \\
      \cmidrule(lr){2-3} \cmidrule(lr){4-5}
      \textbf{Kernel} & \textbf{LR} & \textbf{BERT} & \textbf{Siamese LR} & \textbf{BERT} \\
      \midrule
      TAKE without $\kappa$       & \meanstd{75.82}{0.23}   & \meanstd{80.94}{0.93}   & \meanstd{58.80}{0.92}   & \meanstd{63.66}{1.25} \\
      Last checkpoint             & \meanstd{75.24}{0.57}   & \meanstd{78.28}{0.66}   & \meanstd{58.32}{1.56}   & \meanstd{62.19}{0.57} \\
      Constant                    & \meanstd{76.13}{1.14}   & \meanstd{83.47}{0.71}   & \meanstd{59.44}{0.86}   & \meanstd{64.78}{1.51} \\
      Linear                      & \meanstd{76.20}{0.68}   & \meanstd{85.20}{1.00}   & \meanstd{59.61}{1.11}   & \meanstd{65.60}{1.35} \\
      Cosine                      & \meanstd{76.74}{0.53}   & \meanstd{85.46}{0.64}   & \meanstd{59.14}{0.85}   & \meanstd{66.16}{0.41} \\
      Exponential (\textbf{TAKE}) & \meanstdbf{77.07}{0.09} & \meanstdbf{86.50}{1.01} & \meanstdbf{60.04}{0.12} & \meanstdbf{66.29}{1.91} \\
      \bottomrule
    \end{tabular}
    \begin{tablenotes} \footnotesize
    \item {
        [\textbf{bold}] best results.
      }
    \end{tablenotes}
  \end{threeparttable}
\end{table}

\begin{foldable} Table~\ref{tab:results-abla-kernels} ablates the choice of temporal kernel defined in \S\ref{sec:53-research03-method-gradient-knowledge}.
  The last-checkpoint baseline is the critical exception: it falls \textit{below} the no-$\kappa$ baseline across all backbones, most severely on BERT ($-2.7$ on AG News, $-1.5$ on QQP).
  This is precisely the failure mode identified in Proposition~\ref{prop:traj}: when $\kappa$ is extracted at convergence only, the learner has fully memorized difficult samples, so the \textit{\acl{HSB}} is intensified rather than corrected, and the selected corpus is skewed \DIFdelbegin \DIFdel{towards }\DIFdelend \DIFaddbegin \DIFadd{toward }\DIFaddend harder, less representative examples.
  All other trajectory-based strategies lessen this effect and outperform the no-$\kappa$ baseline (OT only), with decreasing kernels (linear, cosine, exponential) doing so most strongly by concentrating mass on the early-to-mid training phase where the bias-correcting signal is strongest.
  The exponential kernel (TAKE default) achieves the best result in all cases, consistent with the theoretical preference for early-phase concentration of influence.
\end{foldable} \section{Limitations and Ethical Considerations} \label{sec:55-research03-limitations}

\begin{foldable} This method has three limitations.
  (\textbf{1})~\textbf{Gradient norm proxy.}
  $I_{(n,t)}$ uses plain gradient norms rather than Fisher-preconditioned scores; the Goodfellow trick~\cite{goodfellow2015efficient} keeps per-sample cost at $O(C + d_\phi)$ per checkpoint.
  K-FAC~\cite{martens2015optimizing} or EKFAC~\cite{george2018fast} preconditioning are extensions to consider when calibrated magnitude scores are needed.
  (\textbf{2})~\textbf{Encoder dependence.}
  The \ac{OT} cost matrix inherits the geometry of $\phi$; a mismatched encoder distorts transport costs.
  The $L$-Lipschitz assumption on $\ell$ (Theorem~\DIFdelbegin \DIFdel{1}\DIFdelend \DIFaddbegin \DIFadd{\ref{thm:gap}}\DIFaddend ) holds approximately when $\phi$ is task-aligned and embeddings are normalized, but remains an approximation for cross-entropy over discrete text.
  (\textbf{3})~\textbf{Synthetic pool coverage.}
  Pool diversity is bounded by the generator's domain coverage.
  Fine-tuning on $\mathcal{D}_\text{train}$ helps address this.
\end{foldable}

\begin{foldable} Two ethical considerations apply to TAKE.
  (\textbf{1})~\textbf{Bias propagation.}
  Distillation concentrates the statistical properties of $\mathcal{D}_\text{train}$, including demographic biases.
  Corpus auditing before distillation is recommended.
  (\textbf{2})~\textbf{Privacy risk.}
  Language models fine-tuned on sensitive data may memorize training examples, a risk not fully eliminated by synthetic data generation.
  Distance-to-closest-record (DCR) screening of $\mathcal{D}_\text{synth}$ can verify that synthetic samples are not verbatim copies; for regulated domains, differential privacy at the fine-tuning stage provides a formal guarantee against memorization.
\end{foldable}
 \section{Conclusion} \label{sec:56-research03-conclusion}

\begin{foldable} We propose TAKE, the first text \ac{DD} method to simultaneously satisfy all three desiderata: \textit{knowledge-based weighting} via trajectory-integrated influence scores that correct the \acl{HSB} and direct the budget toward informative samples; a \textit{dataset-level distillation objective} via discrete Optimal Transport that globally aligns the selected corpus with the source distribution; and \textit{human-readable output} via selection from \ac{LLM}-generated candidates, producing a distilled corpus that is inspectable, auditable, and transferable across architectures.
  These design choices address gaps that prior text \ac{DD} methods leave open: TAKE matches or exceeds baselines across six language benchmarks at extreme compression, with end-to-end theoretical formulation and guarantees.
  These results suggest that text dataset distillation can help address systemic cost pressures: reducing corpus size cuts cumulative training expenditure, lowers carbon footprint, and respects data-governance constraints without sacrificing transferability.
\end{foldable}
 \clearpage{}
\clearpage{}\acresetall

\chapter{Conclusion} \label{ch:60-conclusion}

\begin{foldable}
  This thesis addresses knowledge distillation along two complementary axes.
  The first is the \textit{model axis}: Chapters~\ref{ch:30-research01} and~\ref{ch:40-research02} compress a learned function under progressively tighter access restrictions---from few-shot black-box to fully data-free black-box---where the teacher withholds both its training data and its internal representations.
  The second is the \textit{data axis}: Chapter~\ref{ch:50-research03} compresses the training corpus itself, distilling a large text dataset into a compact human-readable surrogate under the additional constraint that tokens are discrete and non-differentiable.
  Together, these two axes trace a complete arc from model compression to data compression, both navigated under practical restrictions that standard \ac{KD} methods assume away.
\end{foldable}

\begin{foldable}
  Across both axes, \textit{distributional fidelity}---the degree to which the compact training signal faithfully represents the source distribution---is the binding constraint.
  Each contribution treats fidelity not as a property to hope for at sufficient scale, but as an objective to enforce by design: through adaptive confidence filtering, through structurally diverse synthesis priors, and through trajectory-aware influence weighting coupled with \DIFdelbegin \DIFdel{globally optimal transport}\DIFdelend \DIFaddbegin \DIFadd{a set-level optimal-transport objective that avoids redundant selection}\DIFaddend .
  Section~\ref{sec:61-conclusion-summary} \DIFdelbegin \DIFdel{synthesises }\DIFdelend \DIFaddbegin \DIFadd{synthesizes }\DIFaddend how each contribution \DIFdelbegin \DIFdel{realises }\DIFdelend \DIFaddbegin \DIFadd{realizes }\DIFaddend this principle; Section~\ref{sec:62-conclusion-future} maps what the \DIFdelbegin \DIFdel{programme }\DIFdelend \DIFaddbegin \DIFadd{program }\DIFaddend leaves open.
\end{foldable}

\section{Summary of Contributions} \label{sec:61-conclusion-summary}

\begin{foldable}
  The three contributions trace a progression of increasingly tight constraints, each answered by a more principled form of active fidelity enforcement.
  What unifies them is not a shared architecture but a shared design stance: when data is scarce, access is restricted, or tokens are discrete, fidelity to the source distribution must be built in, not assumed.
\end{foldable}

\begin{foldable}
  \textbf{DivBFKD} (Chapter~\ref{ch:30-research01}) addresses the few-shot black-box \ac{KD} problem, where the student has access to only a small \DIFdelbegin \DIFdel{unlabeled }\DIFdelend \DIFaddbegin \DIFadd{labeled }\DIFaddend image set and may query the teacher solely for \DIFdelbegin \DIFdel{top-1 predictions}\DIFdelend \DIFaddbegin \DIFadd{its predicted class probabilities}\DIFaddend .
  The key observation is that diversity of the training distribution---not merely its size---determines whether \DIFdelbegin \DIFdel{WGAN-synthesised }\DIFdelend \DIFaddbegin \DIFadd{WGAN-synthesized }\DIFaddend samples can provide useful gradient signal.
  DivBFKD enforces fidelity by introducing high-confidence images, those the teacher classifies with high certainty, as real-image exemplars in the WGAN discriminator loop on-the-fly.
  Adaptive per-class \DIFdelbegin \DIFdel{confidence thresholds prevent the teacher's class-wise accuracy bias from skewing which images are admitted, ensuring that the selected exemplars represent }\DIFdelend \DIFaddbegin \DIFadd{thresholds keep the easiest classes from dominating, so each class contributes reliably labeled exemplars from well inside }\DIFaddend the teacher's \DIFdelbegin \DIFdel{decision boundaries rather than its easiest cases}\DIFdelend \DIFaddbegin \DIFadd{class regions to anchor the generator}\DIFaddend .
  An ablation study \DIFdelbegin \DIFdel{confirms that removing adaptive thresholds---reverting to a global cutoff---degrades performance, validating that }\DIFdelend \DIFaddbegin \DIFadd{shows that the adaptive-threshold scheme improves student accuracy over a plain WGAN, and the teacher's }\DIFaddend class-wise \DIFdelbegin \DIFdel{calibration of the filtering criterion is load-bearing, not incidental}\DIFdelend \DIFaddbegin \DIFadd{confidence bias motivates per-class rather than global thresholds}\DIFaddend .
  DivBFKD achieves \DIFdelbegin \DIFdel{state-of-the-art performance among }\DIFdelend \DIFaddbegin \DIFadd{the highest accuracy among the compared }\DIFaddend few-shot black-box \ac{KD} methods \DIFdelbegin \DIFdel{across multiple }\DIFdelend \DIFaddbegin \DIFadd{(BBKD and FS-BBT) across the }\DIFaddend image classification datasets and architectures \DIFaddbegin \DIFadd{evaluated}\DIFaddend .
\end{foldable}

\begin{foldable}
  \textbf{DIPKD} (Chapter~\ref{ch:40-research02}) addresses the extreme case of access restriction: no real data at all, and a teacher that returns only top-1 labels.
  The fidelity problem becomes entirely generative---the student must be trained on a synthetic distribution that \DIFdelbegin \DIFdel{covers the }\DIFdelend \DIFaddbegin \DIFadd{aims to cover the inputs relevant to the }\DIFaddend teacher's \DIFdelbegin \DIFdel{decision surface }\DIFdelend \DIFaddbegin \DIFadd{decisions, }\DIFaddend without any real-image anchor.
  DIPKD addresses this through a three-phase pipeline where each phase serves a distinct purpose.
  The Synthesis phase constructs hierarchical multi-scale image priors from uniform noise with \DIFdelbegin \DIFdel{geometric transforms }\DIFdelend \DIFaddbegin \DIFadd{nonlinear filters }\DIFaddend and semantic cutmixing, \DIFaddbegin \DIFadd{whose masks are irregular regions rather than object segments, }\DIFaddend producing structural diversity before any teacher query.
  The Contrast phase passes these priors through a primer student backbone and applies \DIFdelbegin \DIFdel{a }\DIFdelend \DIFaddbegin \DIFadd{an instance-level }\DIFaddend contrastive objective that \DIFdelbegin \DIFdel{optimises }\DIFdelend \DIFaddbegin \DIFadd{optimizes }\DIFaddend the synthetic dataset directly, pushing \DIFdelbegin \DIFdel{candidates apart }\DIFdelend \DIFaddbegin \DIFadd{each synthetic image away from all others }\DIFaddend in embedding space to prevent mode collapse.
  The Distillation phase then \DIFdelbegin \DIFdel{recovers soft probabilistic labels }\DIFdelend \DIFaddbegin \DIFadd{takes soft outputs }\DIFaddend from the primer \DIFdelbegin \DIFdel{student---approximating }\DIFdelend \DIFaddbegin \DIFadd{student as surrogate soft supervision, in place of }\DIFaddend the dark knowledge the black-box teacher \DIFdelbegin \DIFdel{withholds---and }\DIFdelend \DIFaddbegin \DIFadd{withholds, and }\DIFaddend uses them alongside hard teacher labels to train the final student.
  DIPKD achieves state-of-the-art performance \DIFdelbegin \DIFdel{across twelve benchmarks , with the largest gains in fine-grained }\DIFdelend \DIFaddbegin \DIFadd{on 11 of the 12 benchmarks (all except Tiny-ImageNet), with substantial gains on several complex }\DIFaddend and domain-specific \DIFdelbegin \DIFdel{settings }\DIFdelend \DIFaddbegin \DIFadd{datasets }\DIFaddend where synthetic diversity is most constrained.
\end{foldable}

\begin{foldable}
  \textbf{TAKE} (Chapter~\ref{ch:50-research03}) shifts the problem domain from image model compression to text dataset distillation~(\ac{DD}), and with it the distributional fidelity problem shifts from generation to selection.
  The goal is to replace a large labeled text corpus with a compact surrogate that \DIFdelbegin \DIFdel{preserves }\DIFdelend \DIFaddbegin \DIFadd{approximates }\DIFaddend downstream task performance.
  TAKE is the first text \ac{DD} method to simultaneously satisfy three desiderata: knowledge-based weighting of samples by their task-aligned contribution, a dataset-level distillation objective that aligns the selected corpus with the full source distribution, and human-readable output that is auditable and transferable across architectures.
  Trajectory-aware influence scoring computes importance weights from reciprocal gradient norms accumulated across the training trajectory, correcting the hard-sample bias\DIFdelbegin \DIFdel{that convergence-point scoring introduces by assigning }\DIFdelend \DIFaddbegin \DIFadd{, whereby scoring at the convergence point tends to assign }\DIFaddend inflated importance to \DIFaddbegin \DIFadd{hard examples such as }\DIFaddend noisy or boundary \DIFdelbegin \DIFdel{examples}\DIFdelend \DIFaddbegin \DIFadd{ones}\DIFaddend .
  An \ac{LLM} generation pool produces diverse textual candidates, and discrete \ac{OT} via the Sinkhorn algorithm selects the subset that globally \DIFdelbegin \DIFdel{minimises the entropically regularised }\DIFdelend \DIFaddbegin \DIFadd{minimizes the entropically regularized }\DIFaddend transport cost between the candidate pool and the source distribution.
  The importance-weighted \ac{OT} objective combines task alignment and distributional coverage as complementary objectives\DIFdelbegin \DIFdel{: importance alone directs the budget but ignores global structure; coverage aloneensures spread but ignores contribution; together they close failure modes that neither addresses individually}\DIFdelend \DIFaddbegin \DIFadd{, and in our ablation, adding importance weights to the }\ac{OT} \DIFadd{objective improves over }\ac{OT} \DIFadd{alone}\DIFaddend .
  TAKE matches or exceeds prior text \ac{DD} baselines across six language benchmarks at extreme compression, establishing the first end-to-end text \ac{DD} method with theoretical guarantees.
\end{foldable}

\begin{foldable}
  \DIFaddbegin \DIFadd{Table~\ref{tab:contributions-comparison} compares the three contributions by assumptions, information, mechanism, key result, and limitation.
}\end{foldable}

\begin{table}[ht]
  \centering
  \begin{threeparttable}
    \caption{Comparison of the three thesis contributions.}
    \label{tab:contributions-comparison}
    \setlength{\tabcolsep}{8pt}
    \begin{tabular}{>{\raggedright\arraybackslash}p{2.1cm} >{\raggedright\arraybackslash}p{3.8cm} >{\raggedright\arraybackslash}p{3.8cm} >{\raggedright\arraybackslash}p{3.8cm}}
      \toprule
      \textbf{Aspect} & DivBFKD & DIPKD & TAKE \\
      \midrule
      Assumptions
        & A small labeled image set. A black-box teacher.
        & No real data. A black-box teacher.
        & A labeled text corpus. A fine-tunable \ac{LLM}. \\
      \addlinespace
      Information
        & Class probabilities for queried images.
        & Top-1 label for queried images.
        & Per-sample gradients over the training trajectory. An \ac{LLM}-generated candidate pool. \\
      \addlinespace
      Mechanism
        & WGAN with high-confidence images as extra real images. Adaptive per-class thresholds.
        & Image priors built from noise. Instance-level contrastive learning. A primer student for surrogate soft labels.
        & Trajectory-aware influence scores. Discrete \ac{OT} between the pool and the source distribution. \\
      \addlinespace
      Key result
        & Highest accuracy among the compared few-shot methods. A 17\% gain over standard \ac{KD} on CIFAR-10.
        & State of the art on 11 of 12 benchmarks. Tiny-ImageNet is the exception.
        & Matches or exceeds the compared baselines on six benchmarks at extreme compression. \\
      \addlinespace
      Limitation
        & Confidence is a proxy for coverage. Training queries the teacher more. GAN training can be unstable.
        & Priors are structurally diverse but not semantic. The pipeline needs extra queries and training.
        & The loss-gap bound holds only under assumptions. The encoder is fixed. The generator limits pool coverage. \\
      \bottomrule
    \end{tabular}
  \end{threeparttable}
\end{table}

\begin{foldable}
  \DIFaddend Across all three contributions, the fidelity-first design stance \DIFdelbegin \DIFdel{enables }\DIFdelend \DIFaddbegin \DIFadd{opens the potential for }\DIFaddend deployment in regimes that scale-dependent assumptions foreclose: edge and embedded devices for DivBFKD and DIPKD, and low-resource, privacy-constrained, and carbon-sensitive language model development for TAKE.
  Model compression reduces inference cost and latency; dataset distillation reduces training expenditure, data-governance risk, and cumulative carbon footprint across full development lifecycles.
  The practical impact scales with the tightness of the constraint: the more restricted the access regime, the larger the relative gain from enforcing fidelity actively rather than assuming it.
\end{foldable}
 \section{Limitations and Future Directions} \label{sec:62-conclusion-future}

\begin{foldable}
  Each contribution operates within boundaries that are not incidental but structural: they follow from the design choices that make each method tractable under its respective access constraints.
  Mapping these boundaries is therefore a prerequisite for extending the \DIFdelbegin \DIFdel{programme}\DIFdelend \DIFaddbegin \DIFadd{program}\DIFaddend , not a concession to its incompleteness.
\end{foldable}

\subsection{Technical Limitations} \label{subsec:62-conclusion-future-technical}

\begin{foldable}
  DivBFKD inherits the instability of adversarial training at scale: the WGAN loop is sensitive to generator-discriminator balance, and on-the-fly injection of high-confidence real images interacts with discriminator convergence in ways we have not fully \DIFdelbegin \DIFdel{characterised }\DIFdelend \DIFaddbegin \DIFadd{characterized }\DIFaddend at larger resolutions.
  The adaptive thresholds are calibrated against the teacher's class-wise accuracy on the available few-shot set and therefore correct for distributional bias within those classes; they are blind to domain shift between the few-shot set and the teacher's original training distribution.
  More fundamentally, high-confidence filtering \DIFdelbegin \DIFdel{guarantees teacher agreement but not }\DIFdelend \DIFaddbegin \DIFadd{favors images on which the teacher is confident but does not ensure }\DIFaddend task-relevant diversity: a set of confident images may still be clustered in feature space if the teacher is confidently right for consistent visual reasons.
\end{foldable}

\begin{foldable}
  DIPKD's Synthesis phase generates priors from hierarchical multi-scale uniform noise with \DIFdelbegin \DIFdel{geometric transforms }\DIFdelend \DIFaddbegin \DIFadd{nonlinear filters }\DIFaddend and semantic cutmixing.
  This procedure yields structural diversity but is semantically agnostic: the priors may not cover the natural image manifold in \DIFdelbegin \DIFdel{fine-grained }\DIFdelend \DIFaddbegin \DIFadd{complex }\DIFaddend or domain-specific tasks where intra-class variation is subtle.
  \DIFaddbegin \DIFadd{In particular, semantic cutmixing uses irregular masks rather than object segments, so a teacher label on a synthetic image does not make it a meaningful instance of that class.
  }\DIFaddend The primer student is pre-selected by architecture without a criterion matching its geometric capacity to the synthetic distribution it will process, and the soft labels it produces are an approximation of the teacher's dark knowledge, with no principled bound on the gap between primer and teacher distributions.
\end{foldable}

\begin{foldable}
  TAKE's trajectory-aware influence scores are computed from reciprocal gradient norms; trajectory integration substantially mitigates hard-sample bias, but the scores lack Fisher preconditioning and rankings among samples with similar gradient magnitudes may be imperfectly calibrated.
  The discrete \ac{OT} cost matrix is defined in the encoder's embedding space, so encoder misalignment propagates directly into the transport plan.
  The \ac{LLM} generation pool fine-tuned on the source corpus risks memorizing rare examples, raising privacy concerns under membership inference, and the Sinkhorn solver scales quadratically in pool size, constraining candidate diversity at extreme compression.
  Across all three contributions, the shared cross-cutting limitation is the assumption of teacher correctness: errors, biases, and long-tail failures in the teacher propagate through each framework without detection or correction.
\end{foldable}

\DIFaddbegin \begin{foldable}
  \DIFadd{Each method pays for fidelity with extra computation or teacher access, and we did not fully measure these costs.
  DivBFKD queries the teacher on generated images throughout adversarial training to compute confidence scores, and again to label the synthetic distillation set.
  DIPKD needs teacher labels for its synthetic images and trains a primer student and runs a contrastive phase before the final student, which is more training than standard distillation.
  TAKE adds per-sample gradient tracking over the training trajectory, }\ac{LLM} \DIFadd{generation, and an optimal-transport solve whose per-iteration cost grows with the product of the source and pool sizes.
  Whether the accuracy gains justify these costs depends on how expensive teacher queries and real data are in a given deployment.
}\end{foldable}

\begin{foldable}
  \DIFadd{The evidence for generalization is bounded by what we tested.
  Claims of generality across architectures, datasets, or domains should be read as limited to the settings evaluated in each chapter, and the theoretical guarantees hold only under their stated assumptions.
}\end{foldable}

\DIFaddend \subsection{Theoretical Open Problems} \label{subsec:62-conclusion-future-theoretical}

\begin{foldable}
  The fidelity-efficiency Pareto frontier is \DIFdelbegin \DIFdel{uncharacterised }\DIFdelend \DIFaddbegin \DIFadd{uncharacterized }\DIFaddend in all three regimes.
  \DIFdelbegin \DIFdel{No }\DIFdelend \DIFaddbegin \DIFadd{Beyond the assumption-dependent loss-gap bound for TAKE (Theorem~\ref{thm:gap}), no general }\DIFaddend formal bound relates the distributional coverage of a synthetic or selected set to the generalization gap of a student trained on it, because the relevant bound must account jointly for the synthesis or selection procedure, the student architecture, and the task loss.
  \DIFaddbegin \DIFadd{Relatedly, distributional fidelity is a design principle that the three methods operationalize differently, not a general theoretical framework, and no common quantitative measure yet links it to downstream performance across the three settings.
  }\DIFaddend Privacy-fidelity coupling is similarly unexplored: methods that enforce fidelity by anchoring synthesis or selection to real data---as both DivBFKD and TAKE do---may inadvertently increase leakage under membership inference, yet the privacy-fidelity trade-off has no formal treatment in the current framework.
\end{foldable}

\begin{foldable}
  The hard-sample bias that TAKE's trajectory integration corrects is empirically documented but causally unresolved: it is not known whether the bias is a selection effect, in which high-gradient samples are disproportionately noisy, or a representation effect, in which convergence-point gradients systematically underweight samples that shaped early feature learning.
  Resolving this distinction has direct implications for when and how trajectory integration should be designed.
  Finally, evaluation protocols for black-box \ac{KD} and text \ac{DD} are \DIFdelbegin \DIFdel{unstandardised }\DIFdelend \DIFaddbegin \DIFadd{unstandardized }\DIFaddend across the field---different works vary teacher architectures, student capacity ratios, compression factors, and test procedures---making reliable cross-method comparison and \DIFdelbegin \DIFdel{characterisation }\DIFdelend \DIFaddbegin \DIFadd{characterization }\DIFaddend of fidelity gains difficult.
\end{foldable}

\subsection{Future Directions} \label{subsec:62-conclusion-future-directions}

\begin{foldable}
  The most direct extension of DivBFKD is to replace the WGAN generator with a diffusion model conditioned on teacher predictions, which would simultaneously address training instability and semantic agnosticism of the synthesis prior.
  Per-class confidence gates could be learned via bilevel \DIFdelbegin \DIFdel{optimisation}\DIFdelend \DIFaddbegin \DIFadd{optimization}\DIFaddend , with the outer loop \DIFdelbegin \DIFdel{minimising }\DIFdelend \DIFaddbegin \DIFadd{minimizing }\DIFaddend student-teacher disagreement on a held-out set, replacing the current heuristic threshold with a criterion directly tied to the distillation objective.
  \DIFdelbegin \DIFdel{Analysing }\DIFdelend \DIFaddbegin \DIFadd{Analyzing }\DIFaddend the robustness of DivBFKD-distilled models to covariate shift and adversarial perturbation---as flagged in Chapter~\ref{ch:30-research01}---is a prerequisite for deployment assurance, since capability transfer and safety transfer are not guaranteed to coincide.
\end{foldable}

\begin{foldable}
  DIPKD's synthesis prior could be grounded in teacher-gradient saliency maps, even under black-box access, using gradient-free attribution methods to condition noise \DIFdelbegin \DIFdel{initialisation }\DIFdelend \DIFaddbegin \DIFadd{initialization }\DIFaddend on the teacher's discriminative regions.
  Neural architecture search applied to primer student selection, guided by a diversity criterion on the synthetic distribution, would replace unprincipled pre-selection.
  Replacing logit matching in the Distillation phase with metric learning on primer embeddings would reduce sensitivity to the primer-teacher distribution gap by targeting structural similarity rather than output matching.
\end{foldable}

\begin{foldable}
  For TAKE, Fisher-preconditioning the influence scores is the natural next step: it would account for loss-landscape curvature and provide better-calibrated rankings, at the cost of Hessian estimation.
  Differential privacy for the \ac{LLM} generation pool under explicit $\varepsilon$ budgets would formalize the privacy guarantee; greedy lexicographic selection or neural hashing for candidate embeddings would reduce the quadratic scaling of discrete \ac{OT} and enable larger, more diverse pools.
  The broader agenda is integrating distributional fidelity with robustness, fairness, and calibration as co-equal design constraints from the outset---not as post-hoc evaluation criteria applied after the fidelity objective has already been \DIFdelbegin \DIFdel{optimised}\DIFdelend \DIFaddbegin \DIFadd{optimized}\DIFaddend .
\end{foldable}
 \clearpage{}

\backmatter
\bibliography{references}

\end{document}
